% ============================================================================
% letsgolegacy-framework.tex -- dokumentacja techniczna frameworku letsgolegacy.
%
% CZYM JEST TEN DOKUMENT. Technicznym opisem całego frameworku letsgolegacy --
% wszystkich publicznych repozytoriów letsgolegacy.* -- i procedurą pracy z nim,
% po polsku. To NIE jest research paper w rozumieniu architecture-standards
% (docs/research/00-RESEARCH-DOCUMENTATION.md): nie stawia pytań badawczych i nie
% wprowadza własnych liczb ani faktów. Jest prezentacją dokumentów, które są
% źródłem prawdy: plików Markdown, schematów, skryptów, workflowów i kodu w
% repozytoriach letsgolegacy.* na commitach wymienionych w sekcji "Zakres, wersja
% i źródła". Gdy ten dokument i źródło się różnią, rację ma źródło, a ten dokument
% trzeba poprawić. Każda podsekcja kończy się listą źródeł (\zrodla).
%
% \paperstatus JEST PRZEDEFINIOWANY. W szablonie oznacza DRAFT albo VERIFIED
% badania; tutaj niesie fazę frameworku, którą dokument opisuje, i wersję samego
% dokumentu: commit, z którego zbudowano PDF. Workflow zapisuje ją w version.tex
% przed kompilacją; lokalnie tego pliku nie ma i nagłówek mówi "build lokalny".
%
% PRZYJĘTY DRYF. Dokument jest ręcznie redagowaną prezentacją źródeł; nic
% mechanicznie nie pilnuje, że zmiana w źródle dotrze tutaj.
%
% BUDOWANIE. .github/workflows/build-docs-pdf.yml. Lokalnie:
%   npm ci && node scripts/render-diagrams.mjs
%   cd docs/papers && latexmk -pdf letsgolegacy-framework.tex
% ============================================================================
\documentclass[11pt,a4paper]{article}

\def\houselang{polish}
\IfFileExists{version.tex}{\input{version.tex}}{\def\docversion{build lokalny}}
\newcommand{\paperstatus}{Faza 01 --- \docversion}
\newcommand{\headerleft}{letsgolegacy --- dokumentacja techniczna frameworku}
\newcommand{\pdftitleline}{letsgolegacy: dokumentacja techniczna frameworku}
\newcommand{\pdfauthorline}{letsgolegacy}
% ============================================================================
% HOUSE-PREAMBLE.tex -- the estate's house LaTeX style, as a file you copy in
% and \input rather than a block you copy out.
%
% See 00-RESEARCH-DOCUMENTATION.md, "The house preamble is a file, not a
% block". Copy this file to <repo>/docs/papers/house-preamble.tex (or
% docs/research/papers/ for a study paper) and \input it from every document
% in the repository that wants the house look. One copy per repository; every
% document in that repository \inputs the same one.
%
% WHY A FILE. The estate's first attempt at this shipped a template whose
% preamble was meant to be copied into each new document. Copying is what
% drifts: `ab-ove`'s two overview editions carried none of the seven house
% markers -- not the colors, the section formatting, the header, the status
% marker, the path macro, the geometry or the link style -- and nothing said
% so, because a convention about copying correctly has nothing to fail. The
% Beamer theme, which was always distributed as a file, propagated
% byte-identically over the same period. Distribute the file.
%
% ----------------------------------------------------------------------------
% THE CONTRACT. The including document defines these BEFORE \input, and
% carries its own \documentclass and \begin{document} -- this file has
% neither, so that one preamble can serve documents with different classes.
%
%   \paperstatus    Right-hand header, and the title block's status line.
%                   A study paper: DRAFT or VERIFIED, mirroring the study's
%                   own header. Any other document: its own status, named in
%                   that document's header comment as repurposed.
%   \headerleft     Left-hand header. Usually "Project --- Document title".
%   \pdftitleline   PDF metadata title.
%   \pdfauthorline  PDF metadata author.
%
% Optional, both defaulted below if the document does not define them:
%
%   \houselang      babel language for this edition, e.g. \def\houselang{polish}.
%                   Defaults to english. Set it in the edition file, not here:
%                   it is the one thing that is genuinely per-edition.
%   \editionsuffix  This edition's diagram-file suffix, e.g.
%                   \def\editionsuffix{.pl}. Defaults to empty. \dgm reads it,
%                   so a figure is called by slug and resolves per edition --
%                   see "Figures" below.
%
% A minimal English document:
%
%   \documentclass[11pt,a4paper]{article}
%   \newcommand{\paperstatus}{DRAFT}
%   \newcommand{\headerleft}{Project --- Overview}
%   \newcommand{\pdftitleline}{Project: an overview}
%   \newcommand{\pdfauthorline}{Author Name}
%   % ============================================================================
% HOUSE-PREAMBLE.tex -- the estate's house LaTeX style, as a file you copy in
% and \input rather than a block you copy out.
%
% See 00-RESEARCH-DOCUMENTATION.md, "The house preamble is a file, not a
% block". Copy this file to <repo>/docs/papers/house-preamble.tex (or
% docs/research/papers/ for a study paper) and \input it from every document
% in the repository that wants the house look. One copy per repository; every
% document in that repository \inputs the same one.
%
% WHY A FILE. The estate's first attempt at this shipped a template whose
% preamble was meant to be copied into each new document. Copying is what
% drifts: `ab-ove`'s two overview editions carried none of the seven house
% markers -- not the colors, the section formatting, the header, the status
% marker, the path macro, the geometry or the link style -- and nothing said
% so, because a convention about copying correctly has nothing to fail. The
% Beamer theme, which was always distributed as a file, propagated
% byte-identically over the same period. Distribute the file.
%
% ----------------------------------------------------------------------------
% THE CONTRACT. The including document defines these BEFORE \input, and
% carries its own \documentclass and \begin{document} -- this file has
% neither, so that one preamble can serve documents with different classes.
%
%   \paperstatus    Right-hand header, and the title block's status line.
%                   A study paper: DRAFT or VERIFIED, mirroring the study's
%                   own header. Any other document: its own status, named in
%                   that document's header comment as repurposed.
%   \headerleft     Left-hand header. Usually "Project --- Document title".
%   \pdftitleline   PDF metadata title.
%   \pdfauthorline  PDF metadata author.
%
% Optional, both defaulted below if the document does not define them:
%
%   \houselang      babel language for this edition, e.g. \def\houselang{polish}.
%                   Defaults to english. Set it in the edition file, not here:
%                   it is the one thing that is genuinely per-edition.
%   \editionsuffix  This edition's diagram-file suffix, e.g.
%                   \def\editionsuffix{.pl}. Defaults to empty. \dgm reads it,
%                   so a figure is called by slug and resolves per edition --
%                   see "Figures" below.
%
% A minimal English document:
%
%   \documentclass[11pt,a4paper]{article}
%   \newcommand{\paperstatus}{DRAFT}
%   \newcommand{\headerleft}{Project --- Overview}
%   \newcommand{\pdftitleline}{Project: an overview}
%   \newcommand{\pdfauthorline}{Author Name}
%   % ============================================================================
% HOUSE-PREAMBLE.tex -- the estate's house LaTeX style, as a file you copy in
% and \input rather than a block you copy out.
%
% See 00-RESEARCH-DOCUMENTATION.md, "The house preamble is a file, not a
% block". Copy this file to <repo>/docs/papers/house-preamble.tex (or
% docs/research/papers/ for a study paper) and \input it from every document
% in the repository that wants the house look. One copy per repository; every
% document in that repository \inputs the same one.
%
% WHY A FILE. The estate's first attempt at this shipped a template whose
% preamble was meant to be copied into each new document. Copying is what
% drifts: `ab-ove`'s two overview editions carried none of the seven house
% markers -- not the colors, the section formatting, the header, the status
% marker, the path macro, the geometry or the link style -- and nothing said
% so, because a convention about copying correctly has nothing to fail. The
% Beamer theme, which was always distributed as a file, propagated
% byte-identically over the same period. Distribute the file.
%
% ----------------------------------------------------------------------------
% THE CONTRACT. The including document defines these BEFORE \input, and
% carries its own \documentclass and \begin{document} -- this file has
% neither, so that one preamble can serve documents with different classes.
%
%   \paperstatus    Right-hand header, and the title block's status line.
%                   A study paper: DRAFT or VERIFIED, mirroring the study's
%                   own header. Any other document: its own status, named in
%                   that document's header comment as repurposed.
%   \headerleft     Left-hand header. Usually "Project --- Document title".
%   \pdftitleline   PDF metadata title.
%   \pdfauthorline  PDF metadata author.
%
% Optional, both defaulted below if the document does not define them:
%
%   \houselang      babel language for this edition, e.g. \def\houselang{polish}.
%                   Defaults to english. Set it in the edition file, not here:
%                   it is the one thing that is genuinely per-edition.
%   \editionsuffix  This edition's diagram-file suffix, e.g.
%                   \def\editionsuffix{.pl}. Defaults to empty. \dgm reads it,
%                   so a figure is called by slug and resolves per edition --
%                   see "Figures" below.
%
% A minimal English document:
%
%   \documentclass[11pt,a4paper]{article}
%   \newcommand{\paperstatus}{DRAFT}
%   \newcommand{\headerleft}{Project --- Overview}
%   \newcommand{\pdftitleline}{Project: an overview}
%   \newcommand{\pdfauthorline}{Author Name}
%   \input{house-preamble}
%   \begin{document}
%   ...
%
% Its Polish sibling differs by two lines, both before the \input:
%
%   \def\houselang{polish}
%   \def\editionsuffix{.pl}
%
% ----------------------------------------------------------------------------
% PACKAGES THIS FILE DOES NOT LOAD. The house set below is what every document
% needs. A document that needs more -- tikz, pifont, longtable, fancyvrb,
% amssymb -- loads it in its own preamble AFTER the \input, and says in a
% comment that it is an addition rather than part of the house set, so the
% next reader can tell the two apart at a glance.
% ============================================================================

% ----------------------------------------------------------------------------
% Encoding, then language. fontenc before inputenc before babel: babel reads
% the active font encoding when it sets up its shorthands, so loading it first
% gives an edition that is translated but typeset wrong.
% ----------------------------------------------------------------------------
\usepackage[T1]{fontenc}
\usepackage[utf8]{inputenc}

% The edition's language, defaulted so a monolingual document defines nothing.
% \PassOptionsToPackage rather than \usepackage[\houselang]{babel}: the option
% list here is a macro, and PassOptionsToPackage is the mechanism built for
% passing one.
\providecommand{\houselang}{english}
\PassOptionsToPackage{\houselang}{babel}
\usepackage{babel}

% ----------------------------------------------------------------------------
% The house package set.
% ----------------------------------------------------------------------------
\usepackage{geometry}
\usepackage{xcolor}
\usepackage{hyperref}
\usepackage{graphicx}
\usepackage{enumitem}
\usepackage{fancyhdr}
\usepackage{titlesec}
\usepackage{parskip}
\usepackage{booktabs}
\usepackage{array}
\usepackage{tabularx}

% ----------------------------------------------------------------------------
% Page geometry.
% ----------------------------------------------------------------------------
\geometry{
    a4paper,
    left=2.2cm,
    right=2.2cm,
    top=2.5cm,
    bottom=2.5cm
}

% ----------------------------------------------------------------------------
% House colors. Extracted from the estate's first formal documents
% (business_analysis.tex, pitch-deck-demium.tex) and unchanged since: a
% document from any repository should look like it came from the same shop.
% ----------------------------------------------------------------------------
\definecolor{primaryblue}{RGB}{41,128,185}
\definecolor{darkblue}{RGB}{31,78,121}
\definecolor{lightgray}{RGB}{236,240,241}
\definecolor{darkgray}{RGB}{52,73,94}
\definecolor{accentgreen}{RGB}{39,174,96}
\definecolor{accentorange}{RGB}{230,126,34}
\definecolor{accentred}{RGB}{231,76,60}

% ----------------------------------------------------------------------------
% Section formatting.
% ----------------------------------------------------------------------------
\titleformat{\section}
    {\color{primaryblue}\Large\bfseries\sffamily}
    {\thesection.}{0.5em}{}[\titlerule]

\titleformat{\subsection}
    {\color{darkblue}\large\bfseries\sffamily}
    {\thesubsection}{0.5em}{}

\titleformat{\subsubsection}
    {\color{darkgray}\normalsize\bfseries\sffamily}
    {\thesubsubsection}{0.5em}{}

% ----------------------------------------------------------------------------
% Header and footer. \paperstatus rides in the right-hand header of every
% page, so a draft cannot be mistaken for a final document at any page a
% reader happens to open it at -- which is the whole reason the marker is in
% the header rather than only on the title page.
% ----------------------------------------------------------------------------
\pagestyle{fancy}
\fancyhf{}
\renewcommand{\headrulewidth}{0.4pt}
\fancyhead[L]{\small\sffamily\color{darkgray} \headerleft}
\fancyhead[R]{\small\sffamily\color{accentred} \paperstatus\ --- \today}
\fancyfoot[C]{\thepage}

% ----------------------------------------------------------------------------
% Links and PDF metadata.
% ----------------------------------------------------------------------------
\hypersetup{
    colorlinks=true,
    linkcolor=primaryblue,
    urlcolor=primaryblue,
    pdfauthor={\pdfauthorline},
    pdftitle={\pdftitleline}
}

% ----------------------------------------------------------------------------
% Figures. A diagram's source is one .mmd file rendered to vector PDF by the
% repository's render step; the document includes the rendered file.
%
% \dgm resolves the path from a slug and the edition, so the body text names a
% diagram the way the diagrams directory does and never writes a path:
%
%   \includegraphics[width=0.78\linewidth]{\dgm{a1-system-context}}
%
% resolves to ../diagrams/rendered/a1-system-context.pdf in the default
% edition and ../diagrams/rendered/a1-system-context.pl.pdf in an edition that
% set \def\editionsuffix{.pl}. A repository whose editions share one set of
% diagrams simply never sets \editionsuffix, and every edition resolves to the
% same file.
%
% Hardcoding the path instead is what the estate did first, and it costs one
% edit per figure per edition: `ab-ove` maintained `b1-reader-loop.pdf` in the
% English file against `b1-reader-loop.pl.pdf` in the Polish one, by hand, in
% both. At three figures that is survivable; `agent-eval-bench` has 23.
%
% \housediagramdir is overridable for a document that does not sit one level
% below the diagrams directory -- define it before the \input.
% ----------------------------------------------------------------------------
\providecommand{\editionsuffix}{}
\providecommand{\housediagramdir}{../diagrams/rendered}
\newcommand{\dgm}[1]{\housediagramdir/#1\editionsuffix.pdf}

% ----------------------------------------------------------------------------
% House convenience commands.
%
% \repopath is for a path, a filename or an identifier copied from the code;
% it is deliberately the only one of these, so that a document does not grow a
% private synonym for it. \finding marks the sentence a section exists to
% deliver.
% ----------------------------------------------------------------------------
\newcommand{\repopath}[1]{\colorbox{lightgray}{\texttt{#1}}}
\newcommand{\finding}[1]{{\color{primaryblue}\textbf{#1}}}

% ----------------------------------------------------------------------------
% The house title block, as a command so that every document's first page
% looks the same and none of them reimplements it.
%
%   \housetitle{Title}{Subtitle or empty}{Author}{Affiliation or repo URL}{Provenance line}
%
% The provenance line is the one this estate keeps paying for the absence of:
% a document that presents another document names it here, so a reader who
% finds the two disagreeing knows which one is the source. For a study paper
% it names the companion study and the commit; for a project overview, the
% markdown it presents.
% ----------------------------------------------------------------------------
\newcommand{\housetitle}[5]{%
    \thispagestyle{plain}%
    \begin{center}
        {\LARGE\bfseries\sffamily\color{darkblue} #1}

        \ifx\relax#2\relax\else
            \vspace{0.3cm}

            {\Large\sffamily\color{darkgray} #2}
        \fi

        \vspace{0.6cm}

        {\large #3}

        \vspace{0.2cm}

        {\normalsize\color{darkgray} #4}

        \vspace{0.2cm}

        {\normalsize\color{accentred}\sffamily \paperstatus\ --- \today\ --- #5}
    \end{center}

    \vspace{0.4cm}
}

%   \begin{document}
%   ...
%
% Its Polish sibling differs by two lines, both before the \input:
%
%   \def\houselang{polish}
%   \def\editionsuffix{.pl}
%
% ----------------------------------------------------------------------------
% PACKAGES THIS FILE DOES NOT LOAD. The house set below is what every document
% needs. A document that needs more -- tikz, pifont, longtable, fancyvrb,
% amssymb -- loads it in its own preamble AFTER the \input, and says in a
% comment that it is an addition rather than part of the house set, so the
% next reader can tell the two apart at a glance.
% ============================================================================

% ----------------------------------------------------------------------------
% Encoding, then language. fontenc before inputenc before babel: babel reads
% the active font encoding when it sets up its shorthands, so loading it first
% gives an edition that is translated but typeset wrong.
% ----------------------------------------------------------------------------
\usepackage[T1]{fontenc}
\usepackage[utf8]{inputenc}

% The edition's language, defaulted so a monolingual document defines nothing.
% \PassOptionsToPackage rather than \usepackage[\houselang]{babel}: the option
% list here is a macro, and PassOptionsToPackage is the mechanism built for
% passing one.
\providecommand{\houselang}{english}
\PassOptionsToPackage{\houselang}{babel}
\usepackage{babel}

% ----------------------------------------------------------------------------
% The house package set.
% ----------------------------------------------------------------------------
\usepackage{geometry}
\usepackage{xcolor}
\usepackage{hyperref}
\usepackage{graphicx}
\usepackage{enumitem}
\usepackage{fancyhdr}
\usepackage{titlesec}
\usepackage{parskip}
\usepackage{booktabs}
\usepackage{array}
\usepackage{tabularx}

% ----------------------------------------------------------------------------
% Page geometry.
% ----------------------------------------------------------------------------
\geometry{
    a4paper,
    left=2.2cm,
    right=2.2cm,
    top=2.5cm,
    bottom=2.5cm
}

% ----------------------------------------------------------------------------
% House colors. Extracted from the estate's first formal documents
% (business_analysis.tex, pitch-deck-demium.tex) and unchanged since: a
% document from any repository should look like it came from the same shop.
% ----------------------------------------------------------------------------
\definecolor{primaryblue}{RGB}{41,128,185}
\definecolor{darkblue}{RGB}{31,78,121}
\definecolor{lightgray}{RGB}{236,240,241}
\definecolor{darkgray}{RGB}{52,73,94}
\definecolor{accentgreen}{RGB}{39,174,96}
\definecolor{accentorange}{RGB}{230,126,34}
\definecolor{accentred}{RGB}{231,76,60}

% ----------------------------------------------------------------------------
% Section formatting.
% ----------------------------------------------------------------------------
\titleformat{\section}
    {\color{primaryblue}\Large\bfseries\sffamily}
    {\thesection.}{0.5em}{}[\titlerule]

\titleformat{\subsection}
    {\color{darkblue}\large\bfseries\sffamily}
    {\thesubsection}{0.5em}{}

\titleformat{\subsubsection}
    {\color{darkgray}\normalsize\bfseries\sffamily}
    {\thesubsubsection}{0.5em}{}

% ----------------------------------------------------------------------------
% Header and footer. \paperstatus rides in the right-hand header of every
% page, so a draft cannot be mistaken for a final document at any page a
% reader happens to open it at -- which is the whole reason the marker is in
% the header rather than only on the title page.
% ----------------------------------------------------------------------------
\pagestyle{fancy}
\fancyhf{}
\renewcommand{\headrulewidth}{0.4pt}
\fancyhead[L]{\small\sffamily\color{darkgray} \headerleft}
\fancyhead[R]{\small\sffamily\color{accentred} \paperstatus\ --- \today}
\fancyfoot[C]{\thepage}

% ----------------------------------------------------------------------------
% Links and PDF metadata.
% ----------------------------------------------------------------------------
\hypersetup{
    colorlinks=true,
    linkcolor=primaryblue,
    urlcolor=primaryblue,
    pdfauthor={\pdfauthorline},
    pdftitle={\pdftitleline}
}

% ----------------------------------------------------------------------------
% Figures. A diagram's source is one .mmd file rendered to vector PDF by the
% repository's render step; the document includes the rendered file.
%
% \dgm resolves the path from a slug and the edition, so the body text names a
% diagram the way the diagrams directory does and never writes a path:
%
%   \includegraphics[width=0.78\linewidth]{\dgm{a1-system-context}}
%
% resolves to ../diagrams/rendered/a1-system-context.pdf in the default
% edition and ../diagrams/rendered/a1-system-context.pl.pdf in an edition that
% set \def\editionsuffix{.pl}. A repository whose editions share one set of
% diagrams simply never sets \editionsuffix, and every edition resolves to the
% same file.
%
% Hardcoding the path instead is what the estate did first, and it costs one
% edit per figure per edition: `ab-ove` maintained `b1-reader-loop.pdf` in the
% English file against `b1-reader-loop.pl.pdf` in the Polish one, by hand, in
% both. At three figures that is survivable; `agent-eval-bench` has 23.
%
% \housediagramdir is overridable for a document that does not sit one level
% below the diagrams directory -- define it before the \input.
% ----------------------------------------------------------------------------
\providecommand{\editionsuffix}{}
\providecommand{\housediagramdir}{../diagrams/rendered}
\newcommand{\dgm}[1]{\housediagramdir/#1\editionsuffix.pdf}

% ----------------------------------------------------------------------------
% House convenience commands.
%
% \repopath is for a path, a filename or an identifier copied from the code;
% it is deliberately the only one of these, so that a document does not grow a
% private synonym for it. \finding marks the sentence a section exists to
% deliver.
% ----------------------------------------------------------------------------
\newcommand{\repopath}[1]{\colorbox{lightgray}{\texttt{#1}}}
\newcommand{\finding}[1]{{\color{primaryblue}\textbf{#1}}}

% ----------------------------------------------------------------------------
% The house title block, as a command so that every document's first page
% looks the same and none of them reimplements it.
%
%   \housetitle{Title}{Subtitle or empty}{Author}{Affiliation or repo URL}{Provenance line}
%
% The provenance line is the one this estate keeps paying for the absence of:
% a document that presents another document names it here, so a reader who
% finds the two disagreeing knows which one is the source. For a study paper
% it names the companion study and the commit; for a project overview, the
% markdown it presents.
% ----------------------------------------------------------------------------
\newcommand{\housetitle}[5]{%
    \thispagestyle{plain}%
    \begin{center}
        {\LARGE\bfseries\sffamily\color{darkblue} #1}

        \ifx\relax#2\relax\else
            \vspace{0.3cm}

            {\Large\sffamily\color{darkgray} #2}
        \fi

        \vspace{0.6cm}

        {\large #3}

        \vspace{0.2cm}

        {\normalsize\color{darkgray} #4}

        \vspace{0.2cm}

        {\normalsize\color{accentred}\sffamily \paperstatus\ --- \today\ --- #5}
    \end{center}

    \vspace{0.4cm}
}

%   \begin{document}
%   ...
%
% Its Polish sibling differs by two lines, both before the \input:
%
%   \def\houselang{polish}
%   \def\editionsuffix{.pl}
%
% ----------------------------------------------------------------------------
% PACKAGES THIS FILE DOES NOT LOAD. The house set below is what every document
% needs. A document that needs more -- tikz, pifont, longtable, fancyvrb,
% amssymb -- loads it in its own preamble AFTER the \input, and says in a
% comment that it is an addition rather than part of the house set, so the
% next reader can tell the two apart at a glance.
% ============================================================================

% ----------------------------------------------------------------------------
% Encoding, then language. fontenc before inputenc before babel: babel reads
% the active font encoding when it sets up its shorthands, so loading it first
% gives an edition that is translated but typeset wrong.
% ----------------------------------------------------------------------------
\usepackage[T1]{fontenc}
\usepackage[utf8]{inputenc}

% The edition's language, defaulted so a monolingual document defines nothing.
% \PassOptionsToPackage rather than \usepackage[\houselang]{babel}: the option
% list here is a macro, and PassOptionsToPackage is the mechanism built for
% passing one.
\providecommand{\houselang}{english}
\PassOptionsToPackage{\houselang}{babel}
\usepackage{babel}

% ----------------------------------------------------------------------------
% The house package set.
% ----------------------------------------------------------------------------
\usepackage{geometry}
\usepackage{xcolor}
\usepackage{hyperref}
\usepackage{graphicx}
\usepackage{enumitem}
\usepackage{fancyhdr}
\usepackage{titlesec}
\usepackage{parskip}
\usepackage{booktabs}
\usepackage{array}
\usepackage{tabularx}

% ----------------------------------------------------------------------------
% Page geometry.
% ----------------------------------------------------------------------------
\geometry{
    a4paper,
    left=2.2cm,
    right=2.2cm,
    top=2.5cm,
    bottom=2.5cm
}

% ----------------------------------------------------------------------------
% House colors. Extracted from the estate's first formal documents
% (business_analysis.tex, pitch-deck-demium.tex) and unchanged since: a
% document from any repository should look like it came from the same shop.
% ----------------------------------------------------------------------------
\definecolor{primaryblue}{RGB}{41,128,185}
\definecolor{darkblue}{RGB}{31,78,121}
\definecolor{lightgray}{RGB}{236,240,241}
\definecolor{darkgray}{RGB}{52,73,94}
\definecolor{accentgreen}{RGB}{39,174,96}
\definecolor{accentorange}{RGB}{230,126,34}
\definecolor{accentred}{RGB}{231,76,60}

% ----------------------------------------------------------------------------
% Section formatting.
% ----------------------------------------------------------------------------
\titleformat{\section}
    {\color{primaryblue}\Large\bfseries\sffamily}
    {\thesection.}{0.5em}{}[\titlerule]

\titleformat{\subsection}
    {\color{darkblue}\large\bfseries\sffamily}
    {\thesubsection}{0.5em}{}

\titleformat{\subsubsection}
    {\color{darkgray}\normalsize\bfseries\sffamily}
    {\thesubsubsection}{0.5em}{}

% ----------------------------------------------------------------------------
% Header and footer. \paperstatus rides in the right-hand header of every
% page, so a draft cannot be mistaken for a final document at any page a
% reader happens to open it at -- which is the whole reason the marker is in
% the header rather than only on the title page.
% ----------------------------------------------------------------------------
\pagestyle{fancy}
\fancyhf{}
\renewcommand{\headrulewidth}{0.4pt}
\fancyhead[L]{\small\sffamily\color{darkgray} \headerleft}
\fancyhead[R]{\small\sffamily\color{accentred} \paperstatus\ --- \today}
\fancyfoot[C]{\thepage}

% ----------------------------------------------------------------------------
% Links and PDF metadata.
% ----------------------------------------------------------------------------
\hypersetup{
    colorlinks=true,
    linkcolor=primaryblue,
    urlcolor=primaryblue,
    pdfauthor={\pdfauthorline},
    pdftitle={\pdftitleline}
}

% ----------------------------------------------------------------------------
% Figures. A diagram's source is one .mmd file rendered to vector PDF by the
% repository's render step; the document includes the rendered file.
%
% \dgm resolves the path from a slug and the edition, so the body text names a
% diagram the way the diagrams directory does and never writes a path:
%
%   \includegraphics[width=0.78\linewidth]{\dgm{a1-system-context}}
%
% resolves to ../diagrams/rendered/a1-system-context.pdf in the default
% edition and ../diagrams/rendered/a1-system-context.pl.pdf in an edition that
% set \def\editionsuffix{.pl}. A repository whose editions share one set of
% diagrams simply never sets \editionsuffix, and every edition resolves to the
% same file.
%
% Hardcoding the path instead is what the estate did first, and it costs one
% edit per figure per edition: `ab-ove` maintained `b1-reader-loop.pdf` in the
% English file against `b1-reader-loop.pl.pdf` in the Polish one, by hand, in
% both. At three figures that is survivable; `agent-eval-bench` has 23.
%
% \housediagramdir is overridable for a document that does not sit one level
% below the diagrams directory -- define it before the \input.
% ----------------------------------------------------------------------------
\providecommand{\editionsuffix}{}
\providecommand{\housediagramdir}{../diagrams/rendered}
\newcommand{\dgm}[1]{\housediagramdir/#1\editionsuffix.pdf}

% ----------------------------------------------------------------------------
% House convenience commands.
%
% \repopath is for a path, a filename or an identifier copied from the code;
% it is deliberately the only one of these, so that a document does not grow a
% private synonym for it. \finding marks the sentence a section exists to
% deliver.
% ----------------------------------------------------------------------------
\newcommand{\repopath}[1]{\colorbox{lightgray}{\texttt{#1}}}
\newcommand{\finding}[1]{{\color{primaryblue}\textbf{#1}}}

% ----------------------------------------------------------------------------
% The house title block, as a command so that every document's first page
% looks the same and none of them reimplements it.
%
%   \housetitle{Title}{Subtitle or empty}{Author}{Affiliation or repo URL}{Provenance line}
%
% The provenance line is the one this estate keeps paying for the absence of:
% a document that presents another document names it here, so a reader who
% finds the two disagreeing knows which one is the source. For a study paper
% it names the companion study and the commit; for a project overview, the
% markdown it presents.
% ----------------------------------------------------------------------------
\newcommand{\housetitle}[5]{%
    \thispagestyle{plain}%
    \begin{center}
        {\LARGE\bfseries\sffamily\color{darkblue} #1}

        \ifx\relax#2\relax\else
            \vspace{0.3cm}

            {\Large\sffamily\color{darkgray} #2}
        \fi

        \vspace{0.6cm}

        {\large #3}

        \vspace{0.2cm}

        {\normalsize\color{darkgray} #4}

        \vspace{0.2cm}

        {\normalsize\color{accentred}\sffamily \paperstatus\ --- \today\ --- #5}
    \end{center}

    \vspace{0.4cm}
}


% Dodatki spoza zestawu house (pakiety, styl bloków kodu, \zrodla, checklisty):
% ============================================================================
% Dodatki do zestawu house (additions, NOT part of the house set).
%
% Ładowane przez letsgolegacy-framework.tex PO % ============================================================================
% HOUSE-PREAMBLE.tex -- the estate's house LaTeX style, as a file you copy in
% and \input rather than a block you copy out.
%
% See 00-RESEARCH-DOCUMENTATION.md, "The house preamble is a file, not a
% block". Copy this file to <repo>/docs/papers/house-preamble.tex (or
% docs/research/papers/ for a study paper) and \input it from every document
% in the repository that wants the house look. One copy per repository; every
% document in that repository \inputs the same one.
%
% WHY A FILE. The estate's first attempt at this shipped a template whose
% preamble was meant to be copied into each new document. Copying is what
% drifts: `ab-ove`'s two overview editions carried none of the seven house
% markers -- not the colors, the section formatting, the header, the status
% marker, the path macro, the geometry or the link style -- and nothing said
% so, because a convention about copying correctly has nothing to fail. The
% Beamer theme, which was always distributed as a file, propagated
% byte-identically over the same period. Distribute the file.
%
% ----------------------------------------------------------------------------
% THE CONTRACT. The including document defines these BEFORE \input, and
% carries its own \documentclass and \begin{document} -- this file has
% neither, so that one preamble can serve documents with different classes.
%
%   \paperstatus    Right-hand header, and the title block's status line.
%                   A study paper: DRAFT or VERIFIED, mirroring the study's
%                   own header. Any other document: its own status, named in
%                   that document's header comment as repurposed.
%   \headerleft     Left-hand header. Usually "Project --- Document title".
%   \pdftitleline   PDF metadata title.
%   \pdfauthorline  PDF metadata author.
%
% Optional, both defaulted below if the document does not define them:
%
%   \houselang      babel language for this edition, e.g. \def\houselang{polish}.
%                   Defaults to english. Set it in the edition file, not here:
%                   it is the one thing that is genuinely per-edition.
%   \editionsuffix  This edition's diagram-file suffix, e.g.
%                   \def\editionsuffix{.pl}. Defaults to empty. \dgm reads it,
%                   so a figure is called by slug and resolves per edition --
%                   see "Figures" below.
%
% A minimal English document:
%
%   \documentclass[11pt,a4paper]{article}
%   \newcommand{\paperstatus}{DRAFT}
%   \newcommand{\headerleft}{Project --- Overview}
%   \newcommand{\pdftitleline}{Project: an overview}
%   \newcommand{\pdfauthorline}{Author Name}
%   % ============================================================================
% HOUSE-PREAMBLE.tex -- the estate's house LaTeX style, as a file you copy in
% and \input rather than a block you copy out.
%
% See 00-RESEARCH-DOCUMENTATION.md, "The house preamble is a file, not a
% block". Copy this file to <repo>/docs/papers/house-preamble.tex (or
% docs/research/papers/ for a study paper) and \input it from every document
% in the repository that wants the house look. One copy per repository; every
% document in that repository \inputs the same one.
%
% WHY A FILE. The estate's first attempt at this shipped a template whose
% preamble was meant to be copied into each new document. Copying is what
% drifts: `ab-ove`'s two overview editions carried none of the seven house
% markers -- not the colors, the section formatting, the header, the status
% marker, the path macro, the geometry or the link style -- and nothing said
% so, because a convention about copying correctly has nothing to fail. The
% Beamer theme, which was always distributed as a file, propagated
% byte-identically over the same period. Distribute the file.
%
% ----------------------------------------------------------------------------
% THE CONTRACT. The including document defines these BEFORE \input, and
% carries its own \documentclass and \begin{document} -- this file has
% neither, so that one preamble can serve documents with different classes.
%
%   \paperstatus    Right-hand header, and the title block's status line.
%                   A study paper: DRAFT or VERIFIED, mirroring the study's
%                   own header. Any other document: its own status, named in
%                   that document's header comment as repurposed.
%   \headerleft     Left-hand header. Usually "Project --- Document title".
%   \pdftitleline   PDF metadata title.
%   \pdfauthorline  PDF metadata author.
%
% Optional, both defaulted below if the document does not define them:
%
%   \houselang      babel language for this edition, e.g. \def\houselang{polish}.
%                   Defaults to english. Set it in the edition file, not here:
%                   it is the one thing that is genuinely per-edition.
%   \editionsuffix  This edition's diagram-file suffix, e.g.
%                   \def\editionsuffix{.pl}. Defaults to empty. \dgm reads it,
%                   so a figure is called by slug and resolves per edition --
%                   see "Figures" below.
%
% A minimal English document:
%
%   \documentclass[11pt,a4paper]{article}
%   \newcommand{\paperstatus}{DRAFT}
%   \newcommand{\headerleft}{Project --- Overview}
%   \newcommand{\pdftitleline}{Project: an overview}
%   \newcommand{\pdfauthorline}{Author Name}
%   \input{house-preamble}
%   \begin{document}
%   ...
%
% Its Polish sibling differs by two lines, both before the \input:
%
%   \def\houselang{polish}
%   \def\editionsuffix{.pl}
%
% ----------------------------------------------------------------------------
% PACKAGES THIS FILE DOES NOT LOAD. The house set below is what every document
% needs. A document that needs more -- tikz, pifont, longtable, fancyvrb,
% amssymb -- loads it in its own preamble AFTER the \input, and says in a
% comment that it is an addition rather than part of the house set, so the
% next reader can tell the two apart at a glance.
% ============================================================================

% ----------------------------------------------------------------------------
% Encoding, then language. fontenc before inputenc before babel: babel reads
% the active font encoding when it sets up its shorthands, so loading it first
% gives an edition that is translated but typeset wrong.
% ----------------------------------------------------------------------------
\usepackage[T1]{fontenc}
\usepackage[utf8]{inputenc}

% The edition's language, defaulted so a monolingual document defines nothing.
% \PassOptionsToPackage rather than \usepackage[\houselang]{babel}: the option
% list here is a macro, and PassOptionsToPackage is the mechanism built for
% passing one.
\providecommand{\houselang}{english}
\PassOptionsToPackage{\houselang}{babel}
\usepackage{babel}

% ----------------------------------------------------------------------------
% The house package set.
% ----------------------------------------------------------------------------
\usepackage{geometry}
\usepackage{xcolor}
\usepackage{hyperref}
\usepackage{graphicx}
\usepackage{enumitem}
\usepackage{fancyhdr}
\usepackage{titlesec}
\usepackage{parskip}
\usepackage{booktabs}
\usepackage{array}
\usepackage{tabularx}

% ----------------------------------------------------------------------------
% Page geometry.
% ----------------------------------------------------------------------------
\geometry{
    a4paper,
    left=2.2cm,
    right=2.2cm,
    top=2.5cm,
    bottom=2.5cm
}

% ----------------------------------------------------------------------------
% House colors. Extracted from the estate's first formal documents
% (business_analysis.tex, pitch-deck-demium.tex) and unchanged since: a
% document from any repository should look like it came from the same shop.
% ----------------------------------------------------------------------------
\definecolor{primaryblue}{RGB}{41,128,185}
\definecolor{darkblue}{RGB}{31,78,121}
\definecolor{lightgray}{RGB}{236,240,241}
\definecolor{darkgray}{RGB}{52,73,94}
\definecolor{accentgreen}{RGB}{39,174,96}
\definecolor{accentorange}{RGB}{230,126,34}
\definecolor{accentred}{RGB}{231,76,60}

% ----------------------------------------------------------------------------
% Section formatting.
% ----------------------------------------------------------------------------
\titleformat{\section}
    {\color{primaryblue}\Large\bfseries\sffamily}
    {\thesection.}{0.5em}{}[\titlerule]

\titleformat{\subsection}
    {\color{darkblue}\large\bfseries\sffamily}
    {\thesubsection}{0.5em}{}

\titleformat{\subsubsection}
    {\color{darkgray}\normalsize\bfseries\sffamily}
    {\thesubsubsection}{0.5em}{}

% ----------------------------------------------------------------------------
% Header and footer. \paperstatus rides in the right-hand header of every
% page, so a draft cannot be mistaken for a final document at any page a
% reader happens to open it at -- which is the whole reason the marker is in
% the header rather than only on the title page.
% ----------------------------------------------------------------------------
\pagestyle{fancy}
\fancyhf{}
\renewcommand{\headrulewidth}{0.4pt}
\fancyhead[L]{\small\sffamily\color{darkgray} \headerleft}
\fancyhead[R]{\small\sffamily\color{accentred} \paperstatus\ --- \today}
\fancyfoot[C]{\thepage}

% ----------------------------------------------------------------------------
% Links and PDF metadata.
% ----------------------------------------------------------------------------
\hypersetup{
    colorlinks=true,
    linkcolor=primaryblue,
    urlcolor=primaryblue,
    pdfauthor={\pdfauthorline},
    pdftitle={\pdftitleline}
}

% ----------------------------------------------------------------------------
% Figures. A diagram's source is one .mmd file rendered to vector PDF by the
% repository's render step; the document includes the rendered file.
%
% \dgm resolves the path from a slug and the edition, so the body text names a
% diagram the way the diagrams directory does and never writes a path:
%
%   \includegraphics[width=0.78\linewidth]{\dgm{a1-system-context}}
%
% resolves to ../diagrams/rendered/a1-system-context.pdf in the default
% edition and ../diagrams/rendered/a1-system-context.pl.pdf in an edition that
% set \def\editionsuffix{.pl}. A repository whose editions share one set of
% diagrams simply never sets \editionsuffix, and every edition resolves to the
% same file.
%
% Hardcoding the path instead is what the estate did first, and it costs one
% edit per figure per edition: `ab-ove` maintained `b1-reader-loop.pdf` in the
% English file against `b1-reader-loop.pl.pdf` in the Polish one, by hand, in
% both. At three figures that is survivable; `agent-eval-bench` has 23.
%
% \housediagramdir is overridable for a document that does not sit one level
% below the diagrams directory -- define it before the \input.
% ----------------------------------------------------------------------------
\providecommand{\editionsuffix}{}
\providecommand{\housediagramdir}{../diagrams/rendered}
\newcommand{\dgm}[1]{\housediagramdir/#1\editionsuffix.pdf}

% ----------------------------------------------------------------------------
% House convenience commands.
%
% \repopath is for a path, a filename or an identifier copied from the code;
% it is deliberately the only one of these, so that a document does not grow a
% private synonym for it. \finding marks the sentence a section exists to
% deliver.
% ----------------------------------------------------------------------------
\newcommand{\repopath}[1]{\colorbox{lightgray}{\texttt{#1}}}
\newcommand{\finding}[1]{{\color{primaryblue}\textbf{#1}}}

% ----------------------------------------------------------------------------
% The house title block, as a command so that every document's first page
% looks the same and none of them reimplements it.
%
%   \housetitle{Title}{Subtitle or empty}{Author}{Affiliation or repo URL}{Provenance line}
%
% The provenance line is the one this estate keeps paying for the absence of:
% a document that presents another document names it here, so a reader who
% finds the two disagreeing knows which one is the source. For a study paper
% it names the companion study and the commit; for a project overview, the
% markdown it presents.
% ----------------------------------------------------------------------------
\newcommand{\housetitle}[5]{%
    \thispagestyle{plain}%
    \begin{center}
        {\LARGE\bfseries\sffamily\color{darkblue} #1}

        \ifx\relax#2\relax\else
            \vspace{0.3cm}

            {\Large\sffamily\color{darkgray} #2}
        \fi

        \vspace{0.6cm}

        {\large #3}

        \vspace{0.2cm}

        {\normalsize\color{darkgray} #4}

        \vspace{0.2cm}

        {\normalsize\color{accentred}\sffamily \paperstatus\ --- \today\ --- #5}
    \end{center}

    \vspace{0.4cm}
}

%   \begin{document}
%   ...
%
% Its Polish sibling differs by two lines, both before the \input:
%
%   \def\houselang{polish}
%   \def\editionsuffix{.pl}
%
% ----------------------------------------------------------------------------
% PACKAGES THIS FILE DOES NOT LOAD. The house set below is what every document
% needs. A document that needs more -- tikz, pifont, longtable, fancyvrb,
% amssymb -- loads it in its own preamble AFTER the \input, and says in a
% comment that it is an addition rather than part of the house set, so the
% next reader can tell the two apart at a glance.
% ============================================================================

% ----------------------------------------------------------------------------
% Encoding, then language. fontenc before inputenc before babel: babel reads
% the active font encoding when it sets up its shorthands, so loading it first
% gives an edition that is translated but typeset wrong.
% ----------------------------------------------------------------------------
\usepackage[T1]{fontenc}
\usepackage[utf8]{inputenc}

% The edition's language, defaulted so a monolingual document defines nothing.
% \PassOptionsToPackage rather than \usepackage[\houselang]{babel}: the option
% list here is a macro, and PassOptionsToPackage is the mechanism built for
% passing one.
\providecommand{\houselang}{english}
\PassOptionsToPackage{\houselang}{babel}
\usepackage{babel}

% ----------------------------------------------------------------------------
% The house package set.
% ----------------------------------------------------------------------------
\usepackage{geometry}
\usepackage{xcolor}
\usepackage{hyperref}
\usepackage{graphicx}
\usepackage{enumitem}
\usepackage{fancyhdr}
\usepackage{titlesec}
\usepackage{parskip}
\usepackage{booktabs}
\usepackage{array}
\usepackage{tabularx}

% ----------------------------------------------------------------------------
% Page geometry.
% ----------------------------------------------------------------------------
\geometry{
    a4paper,
    left=2.2cm,
    right=2.2cm,
    top=2.5cm,
    bottom=2.5cm
}

% ----------------------------------------------------------------------------
% House colors. Extracted from the estate's first formal documents
% (business_analysis.tex, pitch-deck-demium.tex) and unchanged since: a
% document from any repository should look like it came from the same shop.
% ----------------------------------------------------------------------------
\definecolor{primaryblue}{RGB}{41,128,185}
\definecolor{darkblue}{RGB}{31,78,121}
\definecolor{lightgray}{RGB}{236,240,241}
\definecolor{darkgray}{RGB}{52,73,94}
\definecolor{accentgreen}{RGB}{39,174,96}
\definecolor{accentorange}{RGB}{230,126,34}
\definecolor{accentred}{RGB}{231,76,60}

% ----------------------------------------------------------------------------
% Section formatting.
% ----------------------------------------------------------------------------
\titleformat{\section}
    {\color{primaryblue}\Large\bfseries\sffamily}
    {\thesection.}{0.5em}{}[\titlerule]

\titleformat{\subsection}
    {\color{darkblue}\large\bfseries\sffamily}
    {\thesubsection}{0.5em}{}

\titleformat{\subsubsection}
    {\color{darkgray}\normalsize\bfseries\sffamily}
    {\thesubsubsection}{0.5em}{}

% ----------------------------------------------------------------------------
% Header and footer. \paperstatus rides in the right-hand header of every
% page, so a draft cannot be mistaken for a final document at any page a
% reader happens to open it at -- which is the whole reason the marker is in
% the header rather than only on the title page.
% ----------------------------------------------------------------------------
\pagestyle{fancy}
\fancyhf{}
\renewcommand{\headrulewidth}{0.4pt}
\fancyhead[L]{\small\sffamily\color{darkgray} \headerleft}
\fancyhead[R]{\small\sffamily\color{accentred} \paperstatus\ --- \today}
\fancyfoot[C]{\thepage}

% ----------------------------------------------------------------------------
% Links and PDF metadata.
% ----------------------------------------------------------------------------
\hypersetup{
    colorlinks=true,
    linkcolor=primaryblue,
    urlcolor=primaryblue,
    pdfauthor={\pdfauthorline},
    pdftitle={\pdftitleline}
}

% ----------------------------------------------------------------------------
% Figures. A diagram's source is one .mmd file rendered to vector PDF by the
% repository's render step; the document includes the rendered file.
%
% \dgm resolves the path from a slug and the edition, so the body text names a
% diagram the way the diagrams directory does and never writes a path:
%
%   \includegraphics[width=0.78\linewidth]{\dgm{a1-system-context}}
%
% resolves to ../diagrams/rendered/a1-system-context.pdf in the default
% edition and ../diagrams/rendered/a1-system-context.pl.pdf in an edition that
% set \def\editionsuffix{.pl}. A repository whose editions share one set of
% diagrams simply never sets \editionsuffix, and every edition resolves to the
% same file.
%
% Hardcoding the path instead is what the estate did first, and it costs one
% edit per figure per edition: `ab-ove` maintained `b1-reader-loop.pdf` in the
% English file against `b1-reader-loop.pl.pdf` in the Polish one, by hand, in
% both. At three figures that is survivable; `agent-eval-bench` has 23.
%
% \housediagramdir is overridable for a document that does not sit one level
% below the diagrams directory -- define it before the \input.
% ----------------------------------------------------------------------------
\providecommand{\editionsuffix}{}
\providecommand{\housediagramdir}{../diagrams/rendered}
\newcommand{\dgm}[1]{\housediagramdir/#1\editionsuffix.pdf}

% ----------------------------------------------------------------------------
% House convenience commands.
%
% \repopath is for a path, a filename or an identifier copied from the code;
% it is deliberately the only one of these, so that a document does not grow a
% private synonym for it. \finding marks the sentence a section exists to
% deliver.
% ----------------------------------------------------------------------------
\newcommand{\repopath}[1]{\colorbox{lightgray}{\texttt{#1}}}
\newcommand{\finding}[1]{{\color{primaryblue}\textbf{#1}}}

% ----------------------------------------------------------------------------
% The house title block, as a command so that every document's first page
% looks the same and none of them reimplements it.
%
%   \housetitle{Title}{Subtitle or empty}{Author}{Affiliation or repo URL}{Provenance line}
%
% The provenance line is the one this estate keeps paying for the absence of:
% a document that presents another document names it here, so a reader who
% finds the two disagreeing knows which one is the source. For a study paper
% it names the companion study and the commit; for a project overview, the
% markdown it presents.
% ----------------------------------------------------------------------------
\newcommand{\housetitle}[5]{%
    \thispagestyle{plain}%
    \begin{center}
        {\LARGE\bfseries\sffamily\color{darkblue} #1}

        \ifx\relax#2\relax\else
            \vspace{0.3cm}

            {\Large\sffamily\color{darkgray} #2}
        \fi

        \vspace{0.6cm}

        {\large #3}

        \vspace{0.2cm}

        {\normalsize\color{darkgray} #4}

        \vspace{0.2cm}

        {\normalsize\color{accentred}\sffamily \paperstatus\ --- \today\ --- #5}
    \end{center}

    \vspace{0.4cm}
}
, zgodnie
% z kontraktem preambuły: pakiety, których house set nie zawiera, dokument
% ładuje sam i oznacza je jako dodatek.
% ============================================================================

% Bloki kodu, poleceń i konfiguracji. listings, nie minted: minted wymaga
% shell-escape i Pygments w obrazie CI, a dokument nie potrzebuje kolorowania.
\usepackage{listings}

% Długie tabele z kolumnami X (longtable + tabularx w jednym środowisku).
\usepackage{xltabular}

% Znaki Unicode, których inputenc nie zna bez deklaracji.
\DeclareUnicodeCharacter{2192}{\ensuremath{\rightarrow}}
\DeclareUnicodeCharacter{2190}{\ensuremath{\leftarrow}}
\DeclareUnicodeCharacter{2194}{\ensuremath{\leftrightarrow}}
\DeclareUnicodeCharacter{2264}{\ensuremath{\leq}}
\DeclareUnicodeCharacter{2265}{\ensuremath{\geq}}
\DeclareUnicodeCharacter{2260}{\ensuremath{\neq}}
\DeclareUnicodeCharacter{2248}{\ensuremath{\approx}}
\DeclareUnicodeCharacter{00D7}{\ensuremath{\times}}
\DeclareUnicodeCharacter{2212}{\ensuremath{-}}

% Styl bloków kodu: tło i ramka w kolorach house, łamanie długich linii,
% polskie znaki w komentarzach (literate), cudzysłowy proste (upquote).
\lstdefinestyle{kod}{
  basicstyle=\ttfamily\footnotesize,
  columns=fullflexible,
  keepspaces=true,
  breaklines=true,
  frame=single,
  rulecolor=\color{lightgray},
  backgroundcolor=\color{lightgray!40},
  xleftmargin=0.4em,
  xrightmargin=0.4em,
  aboveskip=0.6em,
  belowskip=0.6em,
  showstringspaces=false,
  upquote=true,
  literate=
    {ą}{{\k{a}}}1 {ć}{{\'c}}1 {ę}{{\k{e}}}1 {ł}{{\l{}}}1 {ń}{{\'n}}1
    {ó}{{\'o}}1 {ś}{{\'s}}1 {ź}{{\'z}}1 {ż}{{\.z}}1
    {Ą}{{\k{A}}}1 {Ć}{{\'C}}1 {Ę}{{\k{E}}}1 {Ł}{{\L{}}}1 {Ń}{{\'N}}1
    {Ó}{{\'O}}1 {Ś}{{\'S}}1 {Ź}{{\'Z}}1 {Ż}{{\.Z}}1
    {—}{{---}}1 {–}{{--}}1 {„}{{,,}}1 {”}{{''}}1 {→}{{$\rightarrow$}}1
}
\lstset{style=kod}

% Ścieżka źródłowa na końcu podsekcji: skąd pochodzi to, co napisano wyżej.
% Dokument niczego nie dodaje do źródeł, więc każda podsekcja je nazywa.
\newcommand{\zrodla}[1]{\par\smallskip{\small\raggedright\color{darkgray}\textit{Źródła:} #1\par}\medskip}

% \path (pakiet url, ładowany przez hyperref) służy do ścieżek zbyt długich na
% niełamliwe pole \repopath -- w tabelach i w \zrodla. To nie jest prywatny
% synonim \repopath: \repopath zostaje dla krótkich identyfikatorów w tekście.
\urlstyle{tt}

% Tabele: odrobinę więcej powietrza między wierszami.
\renewcommand{\arraystretch}{1.15}

% Długie nazwy repozytoriów i ścieżki: \path może łamać się także na myślnikach
% (letsgolegacy.secondkey-nopcommerce-bench), a akapit z niełamliwym \repopath
% dostaje trochę luzu, zamiast wychodzić na margines.
\expandafter\def\expandafter\UrlBreaks\expandafter{\UrlBreaks\do\-}
\setlength{\emergencystretch}{3em}

% Checklisty (kształt przewodników estate: sekcja kończy się listą „- [ ]”):
% pusta ramka zamiast punktora, bez dodatkowych pakietów.
\newcommand{\polewyboru}{\raisebox{0.05ex}{\fbox{\rule{0pt}{0.42em}\rule{0.42em}{0pt}}}}
\newenvironment{checklista}{\begin{itemize}[label=\polewyboru, leftmargin=1.7em]}{\end{itemize}}


\begin{document}

\housetitle{letsgolegacy}
  {Dokumentacja techniczna frameworku}
  {Architektura, komponenty, formaty i procedura pracy krok po kroku}
  {\url{https://github.com/konradcinkusz/letsgolegacy.secondkey}}
  {prezentuje repozytoria letsgolegacy.* na commitach z sekcji~\ref{sec:wprowadzenie-wersja}}

\tableofcontents
\clearpage

\section{Wprowadzenie}\label{sec:wprowadzenie}

Ten dokument opisuje framework letsgolegacy w całości: problem, który rozwiązuje,
architekturę, każdy komponent i każde repozytorium, formaty plików, polecenia oraz
procedurę pracy krok po kroku. Jest przeznaczony dla inżyniera, który ma z frameworkiem
pracować, rozwijać go albo ocenić jego wynik. Opisuje stan repozytoriów na konkretnych
commitach (sekcja~\ref{sec:wprowadzenie-wersja}) i niczego do nich nie dodaje.

\subsection{Problem}\label{sec:wprowadzenie-problem}

Systemy biznesowe na .NET Framework są przenoszone na współczesny .NET coraz częściej
przez agentów migracyjnych, takich jak modernize-dotnet (GitHub Copilot app
modernization for .NET). Kryterium odbioru, którym posługuje się samo narzędzie, to:
kandydat się buduje, a istniejące testy przechodzą. To kryterium nie mówi, czy system po
migracji zachowuje się tak samo wobec użytkownika. Migracja z .NET Framework na .NET
zmienia na przykład bibliotekę globalizacji z NLS na ICU, więc porównywanie i sortowanie
napisów zależne od kultury może dać inny wynik, choć build i testy są zielone. Tę
pułapkę bench nopCommerce traktuje jako główny kandydat na moment, wokół którego zbudowana
jest demonstracja: agent migruje, build i testy przechodzą, a weryfikacja i tak znajduje
realną różnicę zachowania.

Rdzeń opisuje swoje zadanie jednym zdaniem: niezależny, deterministyczny dowód, że system
.NET zmigrowany przez agenta AI nadal dotrzymuje kontraktu biznesowego. Nazwa Second Key
pochodzi z obrazu dwóch kluczy. Agent migracyjny przekręca pierwszy klucz: migruje kod.
Second Key przekręca drugi: nagrywa, jak obecny system naprawdę się zachowuje, odtwarza to
zachowanie na wersji zmigrowanej, sprawdza każdą różnicę względem spisanego kontraktu
tego, co musi się dziać i czego nie wolno, i daje osobie zatwierdzającej evidence pack
(pakiet dowodowy) dla każdej zmiany. \finding{Second Key weryfikuje; niczego nie
przepisuje.}

\zrodla{\path{letsgolegacy.secondkey/README.md},
\path{letsgolegacy.secondkey-nopcommerce-bench/docs/WORKPLAN.md},
\path{letsgolegacy.secondkey-nopcommerce-bench/docs/P6-RUNBOOK.md}}

\subsection{Czym jest framework letsgolegacy}\label{sec:wprowadzenie-framework}

Framework letsgolegacy to metoda weryfikowalnej modernizacji systemów legacy .NET razem z
narzędziami, które ją realizują. Narzędzia mieszkają w repozytoriach o prefiksie
\repopath{letsgolegacy.}:

\begin{itemize}
  \item \textbf{Second Key} --- łańcuch weryfikacji: nagranie ruchu (capture), kontrakt,
        odtworzenie (replay) na obu wersjach, porównanie (compare), werdykt i pakiet
        dowodowy. Rdzeń w \path{letsgolegacy.secondkey}, rozszerzenia fazy 02 i 03 w
        \path{secondkey-capture}, \path{secondkey-sandbox} i \path{secondkey-portal}.
  \item \textbf{Portcullis} --- bramka na pull request: deterministyczne analyzery Roslyn
        uruchamiane na zmienionych liniach, z wynikiem w SARIF i plikiem baseline
        (\path{letsgolegacy.portcullis}).
  \item \textbf{secondkey-standards} --- standardy architektury i migracji spisane raz w
        Markdown i dostarczane z jednego źródła agentowi migracyjnemu (jako skill) oraz
        bramce (jako konfiguracja analyzerów).
  \item \textbf{legacy-risk-scan} --- strona, która z pliku projektu (\repopath{.csproj},
        \repopath{packages.config}, \repopath{composer.json}) liczy w przeglądarce raport
        końca wsparcia (EOL) i znanych podatności.
  \item \textbf{bench nopCommerce} --- publiczny okaz: nopCommerce 3.90 na .NET Framework,
        na którym cały łańcuch jest uruchamiany w CI i na którym będzie oceniona migracja
        wykonana przez agenta.
\end{itemize}

Komponenty wymieniają się wyłącznie plikami: nagraniem \repopath{*.skcap}, kontraktem
\repopath{contract.yaml}, przebiegiem \repopath{*.skrun}, werdyktem
\repopath{verdict.json}, logami bramki \repopath{*.sarif} i pakietem dowodowym. Żaden
komponent nie sięga do wnętrza innego, więc każdy da się wymienić, jeśli czyta i pisze te
same artefakty (rozdział~\ref{sec:architektura}).

\zrodla{\path{letsgolegacy.secondkey/README.md},
\path{letsgolegacy.secondkey/docs/WORKPLAN.md},
\path{letsgolegacy.portcullis/docs/WORKPLAN.md},
\path{letsgolegacy.secondkey-standards/README.md},
\path{letsgolegacy.legacy-risk-scan/README.md},
\path{letsgolegacy.secondkey-nopcommerce-bench/README.md}}

\subsection{Nazewnictwo}\label{sec:wprowadzenie-nazwy}

Źródła nazywają cały łańcuch „Second Key” i numerują jego komponenty od C0 do C12 ---
także te, które mieszkają poza rdzeniem: standardy to komponent C4, a Portcullis to
komponent C5, „bramka Second Key”. Ten dokument przyjmuje nazwę nadrzędną
\textbf{letsgolegacy} dla całego frameworku: metody i wszystkich repozytoriów
\repopath{letsgolegacy.*}, łącznie z legacy-risk-scan i benchem, które nie są
komponentami łańcucha. Nazwa \textbf{Second Key} oznacza tu łańcuch weryfikacji i jego
rdzeń (CLI \repopath{sk}, formaty artefaktów), a \textbf{Portcullis} --- bramkę.

Nazwy w kodzie i w formatach się nie zmieniają: przestrzenie nazw \repopath{SecondKey.*},
polecenie \repopath{sk}, plik konfiguracji \repopath{secondkey.yaml}, wersja formatu
kontraktu \path{apiVersion: secondkey/v1}, nazwy repozytoriów i pakietów. Gdy
dokument cytuje źródło, które mówi „Second Key” w szerszym sensie, zaznacza to w
kontekście.

\zrodla{\path{letsgolegacy.secondkey/README.md},
\path{letsgolegacy.portcullis/docs/WORKPLAN.md},
\path{letsgolegacy.secondkey-standards/README.md}}

\subsection{Zakres, wersja i źródła}\label{sec:wprowadzenie-wersja}

Dokument prezentuje publiczne repozytoria frameworku w stanie z tabeli poniżej (gałąź
\repopath{main}). Każda podsekcja kończy się listą plików, z których pochodzi jej treść.
Źródłem prawdy są te pliki: dokumentacja w Markdown, schematy, skrypty, workflowy i kod.
\finding{Gdy ten dokument i źródło się różnią, rację ma źródło}, a dokument należy
poprawić. Nic mechanicznie nie pilnuje, że zmiana w źródle trafi do tego dokumentu --- to
przyjęty koszt ręcznie redagowanej prezentacji; listy źródeł pozwalają sprawdzić każde
twierdzenie osobno.

\begin{tabularx}{\linewidth}{@{}>{\raggedright\arraybackslash}p{6.2cm} l l X@{}}
\toprule
Repozytorium & Commit & Data & Rola \\
\midrule
\path{letsgolegacy.secondkey} & \texttt{fe374c7} & 2026-09-26 & rdzeń Second Key \\
\path{letsgolegacy.portcullis} & \texttt{6a018a2} & 2026-09-26 & bramka C5 \\
\path{letsgolegacy.secondkey-standards} & \texttt{07520bd} & 2026-09-26 & standardy C4 \\
\path{letsgolegacy.secondkey-nopcommerce-bench} & \texttt{e6d4fe9} & 2026-09-26 & publiczny okaz \\
\path{letsgolegacy.legacy-risk-scan} & \texttt{49f8a4c} & 2026-09-26 & raport EOL i CVE \\
\path{letsgolegacy.secondkey-capture} & \texttt{e56068e} & 2026-09-25 & C1, faza 02 \\
\path{letsgolegacy.secondkey-sandbox} & \texttt{75042c6} & 2026-09-25 & C12, faza 02 \\
\path{letsgolegacy.secondkey-portal} & \texttt{e73546f} & 2026-09-25 & C11, faza 03 \\
\bottomrule
\end{tabularx}

Poza zakresem są sprawy nietechniczne i wszystko, czego nie ma w publicznych
repozytoriach. Dokument nie opisuje też pracy w toku na gałęziach, które nie trafiły do
\repopath{main} na wymienionych commitach.

Wersję samego dokumentu --- commit, z którego zbudowano PDF --- podaje nagłówek każdej
strony. PDF buduje workflow \path{.github/workflows/build-docs-pdf.yml} repozytorium
rdzenia przy każdym pushu do \repopath{main}; źródła dokumentu są w
\path{docs/papers/} i \path{docs/diagrams/} (dodatek~\ref{app:ci}).

\zrodla{\path{letsgolegacy.*} na podanych commitach;
\path{architecture-standards/docs/research/00-RESEARCH-DOCUMENTATION.md}}

\subsection{Jak czytać ten dokument}\label{sec:wprowadzenie-czytanie}

\begin{tabularx}{\linewidth}{@{}p{1.9cm} X@{}}
\toprule
Rozdział & Zawartość \\
\midrule
\ref{sec:architektura} & Architektura: łańcuch, artefakty, komponenty C0--C12, zasady projektowe, fazy \\
\ref{sec:procedura} & \textbf{Procedura pracy krok po kroku} \\
\ref{sec:rdzen} & Rdzeń Second Key: projekty, decyzje, testy, CI \\
\ref{sec:formaty} & Formaty artefaktów, w tym pełny język kontraktu i klasy werdyktu \\
\ref{sec:cli} & CLI \texttt{sk}: polecenia, opcje, \texttt{secondkey.yaml}, kody wyjścia \\
\ref{sec:portcullis} & Portcullis: reguły, diff gate, baseline, SARIF, dystrybucja \\
\ref{sec:standardy} & Standardy migracji: 20 reguł, generator, skill, drift check \\
\ref{sec:risk-scan} & legacy-risk-scan \\
\ref{sec:faza02} & Repozytoria fazy 02 i 03: capture, sandbox, portal \\
\ref{sec:bench} & Bench nopCommerce: topologia, skrypty, nagranie, kontrakt, wyniki \\
\ref{sec:bezpieczenstwo} & Bezpieczeństwo i higiena repozytoriów \\
\ref{sec:status} & Status i plan \\
\ref{sec:slownik} & Słownik pojęć \\
\ref{app:klauzule}, \ref{app:ci} & Dodatki: klauzule kontraktu benchu; workflowy CI \\
\bottomrule
\end{tabularx}

Kto chce zacząć pracę, czyta rozdziały~\ref{sec:architektura} i~\ref{sec:procedura};
rozdziały~\ref{sec:rdzen}--\ref{sec:bench} są referencją, do której procedura odsyła.

Konwencje: nazwy techniczne zostają po angielsku, gdy są nazwami własnymi,
identyfikatorami albo ustalonymi terminami (framework, pipeline, pull request, replay,
snapshot); identyfikatory, polecenia i pliki są zapisane krojem kodu, np.
\repopath{sk compare}; wartości z kodu i logów są przepisane bez zmian.

\zrodla{ten dokument}

\clearpage
\section{Architektura frameworku}\label{sec:architektura}

Ten rozdział pokazuje całość: jakie kroki składają się na weryfikację migracji, jakie
pliki przechodzą między krokami, które komponenty je produkują, w których repozytoriach
mieszkają i jakie zasady projektowe obowiązują wszystkie z nich. Szczegóły każdego
komponentu są w rozdziałach~\ref{sec:rdzen}--\ref{sec:bench}.

\subsection{Łańcuch weryfikacji}\label{sec:architektura-lancuch}

Łańcuch ma siedem kroków. Każdy kończy się plikiem, który czyta następny krok.

\begin{tabularx}{\linewidth}{@{}p{2.3cm} >{\raggedright\arraybackslash}X >{\raggedright\arraybackslash}p{3.3cm}@{}}
\toprule
Krok & Komponent & Wynik \\
\midrule
1 Capture & C1 --- proxy nagrywające przed systemem legacy, bez zmiany kodu & \repopath{*.skcap} \\
2 Contract & C2 (miner, faza 02) i człowiek; C3 --- silnik kontraktu & \repopath{contract.yaml} \\
3 Standards & C4 --- standardy dostarczane agentowi migracyjnemu jako skills & pull request agenta \\
4 Gate & C5 --- analyzery Roslyn na zmienionych liniach (Portcullis) & \repopath{*.sarif} \\
5 Replay & C6 --- obie wersje w identycznych warunkach; C7 porównuje pole po polu względem kontraktu & \repopath{*.skrun}, \repopath{verdict.json} \\
6 Mutations & C9 --- wstrzyknięte defekty, które kontrakt musi wykryć & kill rate per klauzula \\
7 Evidence & C10 --- jeden pakiet na pull request & \repopath{evidence/} \\
\bottomrule
\end{tabularx}

\begin{figure}[htbp]
  \centering
  \includegraphics[width=\linewidth,height=0.62\textheight,keepaspectratio]{\dgm{a2-lancuch-weryfikacji}}
  \caption{Łańcuch weryfikacji i artefakty między krokami.}
  \label{fig:lancuch}
\end{figure}

Strona legacy i kandydat są uruchamiane, ale nie są modyfikowane przez łańcuch: capture
stoi przed systemem legacy jako proxy, a replay wysyła nagrane żądania do obu wersji.
Wszystko, co następuje po replay, działa na plikach.

\zrodla{\path{letsgolegacy.secondkey/README.md},
\path{letsgolegacy.secondkey/docs/WORKPLAN.md}}

\subsection{Artefakty są interfejsem}\label{sec:architektura-artefakty}

\finding{Każdy komponent czyta i pisze pliki, nigdy wnętrza innego komponentu.} Dowolny
komponent można wymienić, dopóki czyta i pisze te same artefakty. Dlatego formaty są
wersjonowane, opisane schematami JSON i pilnowane testami na próbkach
(rozdział~\ref{sec:formaty}).

\begin{xltabular}{\linewidth}{@{}>{\raggedright\arraybackslash}p{2.9cm} >{\raggedright\arraybackslash}X >{\raggedright\arraybackslash}X >{\raggedright\arraybackslash}X@{}}
\toprule
Artefakt & Format & Produkuje & Czyta \\
\midrule
\endfirsthead
\toprule
Artefakt & Format & Produkuje & Czyta \\
\midrule
\endhead
\repopath{*.skcap} & JSONL, jedno zdarzenie na linię, korelowane identyfikatorem wymiany (\emph{exchange id}) & C1 capture & C2 miner, C6 replay \\
\repopath{contract.yaml} & YAML walidowany JSON Schema & C2 miner i człowiek & C3 silnik, C7 komparator, C9 mutacje \\
\repopath{standards/} & Markdown → skills i NuGet & C4 (\path{letsgolegacy.secondkey-standards}) & agent migracyjny, C5 bramka \\
\repopath{*.sarif} & SARIF 2.1 & C5 bramka (\path{letsgolegacy.portcullis}) & C10 evidence \\
\repopath{*.skrun} & JSONL, wyniki obu stron dla każdego scenariusza & C6 replay & C7 komparator \\
\repopath{verdict.json} & JSON & C7 komparator, C9 mutacje & C10 evidence, C11 portal \\
\repopath{evidence.dsse} & in-toto / DSSE & C10 evidence (podpis od fazy 03) & C11 portal, audytor \\
\bottomrule
\end{xltabular}

Konsekwencja widoczna w benchu: granica między maszynami może przebiegać wszędzie tam,
gdzie da się skopiować plik. Tylko capture i replay muszą mieć dostęp HTTP do działającego
systemu; kontrakt, porównanie i pakiet dowodowy to „plik na wejściu, plik na wyjściu”
(sekcja~\ref{sec:architektura-topologia}).

\zrodla{\path{letsgolegacy.secondkey/docs/WORKPLAN.md},
\path{letsgolegacy.secondkey/docs/adr/0001-artifacts-are-the-interface.md},
\path{letsgolegacy.secondkey/schemas/skcap.schema.json},
\path{letsgolegacy.secondkey-nopcommerce-bench/docs/TOPOLOGY.md}}

\subsection{Komponenty C0--C12 i repozytoria}\label{sec:architektura-komponenty}

\begin{xltabular}{\linewidth}{@{}l >{\raggedright\arraybackslash}X >{\raggedright\arraybackslash}p{5.2cm} l@{}}
\toprule
Id & Komponent & Repozytorium & Faza \\
\midrule
\endfirsthead
\toprule
Id & Komponent & Repozytorium & Faza \\
\midrule
\endhead
C0 & CLI \repopath{sk} & \path{letsgolegacy.secondkey} & 01 \\
C1 & capture HTTP (proxy YARP, tylko ruch przychodzący) & \path{letsgolegacy.secondkey} & 01 \\
C1 & capture SQL, kolejek, tokenizacja PII & \path{letsgolegacy.secondkey-capture} & 02 \\
C2 & miner kontraktu (\repopath{sk mine}) & \path{letsgolegacy.secondkey} & 02 \\
C3 & silnik kontraktu & \path{letsgolegacy.secondkey} & 01 \\
C4 & pakiet standardów & \path{letsgolegacy.secondkey-standards} & 01 \\
C5 & bramka na PR (Portcullis) & \path{letsgolegacy.portcullis} & 01 \\
C6 & replay i reset stanu & \path{letsgolegacy.secondkey} & 01 \\
C7 & komparator i werdykt & \path{letsgolegacy.secondkey} & 01 \\
C8 & doradcza triage różnic (nigdy nie zmienia werdyktu) & \path{letsgolegacy.secondkey} & 02--03 \\
C9 & silnik mutacji (\repopath{sk mutate}) & \path{letsgolegacy.secondkey} & 02 \\
C10 & pakiet dowodowy (bez podpisu w fazie 01, podpis w fazie 03) & \path{letsgolegacy.secondkey} & 01, 03 \\
C11 & portal dowodów & \path{letsgolegacy.secondkey-portal} & 03 \\
C12 & pakowanie w granicy klienta (sandbox) & \path{letsgolegacy.secondkey-sandbox} & 02 \\
\bottomrule
\end{xltabular}

\begin{figure}[htbp]
  \centering
  \includegraphics[width=\linewidth,height=0.7\textheight,keepaspectratio]{\dgm{a1-mapa-repozytoriow}}
  \caption{Repozytoria frameworku letsgolegacy i to, co przepływa między nimi. Nazwy bez wspólnego prefiksu \texttt{letsgolegacy.}}
  \label{fig:mapa}
\end{figure}

W fazie 01 \repopath{sk mine} i \repopath{sk mutate} istnieją jako stuby, które kończą
się kodem 70 (\repopath{NotImplemented}); kontrakt pisze człowiek.

\zrodla{\path{letsgolegacy.secondkey/README.md},
\path{letsgolegacy.secondkey/docs/WORKPLAN.md},
\path{letsgolegacy.secondkey/docs/cli.md},
\path{letsgolegacy.secondkey-capture/docs/WORKPLAN.md},
\path{letsgolegacy.secondkey-sandbox/docs/WORKPLAN.md},
\path{letsgolegacy.secondkey-portal/docs/WORKPLAN.md}}

\subsection{Zasady projektowe}\label{sec:architektura-zasady}

Rdzeń egzekwuje siedem zasad projektowych. Pozostałe repozytoria się na nie powołują.

\begin{enumerate}
  \item \textbf{Werdykt jest deterministyczny.} O pass/fail decyduje wyłącznie kod bez
        modelu językowego. Model może szkicować klauzule kontraktu albo porządkować
        różnice do przeglądu, ale nigdy nie decyduje. Nic w \repopath{SecondKey.Contract}
        ani \repopath{SecondKey.Compare} nie woła modelu językowego; ten sam kontrakt i
        ten sam przebieg dają zawsze ten sam werdykt (poza polem \repopath{createdAt}).
  \item \textbf{Wszystko działa w granicy klienta; domyślnie zero ruchu wychodzącego.}
        Repozytorium \path{letsgolegacy.secondkey-sandbox} pakuje łańcuch dla różnych
        rodzajów takiej granicy.
  \item \textbf{Tylko technologie, które odbiorca już ma}: .NET, SQL Server, Entra ID,
        GitHub albo Azure DevOps, Kubernetes.
  \item \textbf{System legacy jest nietykalny}: capture nie wymaga zmiany kodu ani
        rekompilacji. Czas i losowość w kodzie legacy (\repopath{DateTime.Now},
        \repopath{Guid.NewGuid()}) obsługuje normalizacja w komparatorze, a nie zmiana
        systemu legacy.
  \item \textbf{Dowód jest artefaktem kryptograficznym, nie dashboardem}: pakiet zawiera
        skróty SHA-256 wszystkich plików i niepodpisany statement in-toto; podpis dochodzi
        w fazie 03 bez zmiany formatu.
  \item \textbf{Brak własnej usługi tożsamości}; używany jest IdP klienta.
  \item \textbf{Każdy krok działa z CLI w pipeline}; portal jest oknem na dowody, a nie
        interfejsem.
\end{enumerate}

Obok siedmiu zasad cały łańcuch wiążą trzy decyzje:

\begin{itemize}
  \item \textbf{Różnica bez wyjaśnienia zawodzi (fail closed).} Różnica, której nie tłumaczy
        żadna klauzula ani reguła normalizacji, jest regresją \repopath{uncovered-difference}
        i blokuje, dopóki przeglądnięta zmiana kontraktu nie doda reguły normalizacji
        (różnica nie ma znaczenia) albo klauzuli (ma znaczenie). Strona bez odpowiedzi to
        regresja \repopath{missing-result} (ADR 0003, decyzja D-2 analizy fazy 01).
  \item \textbf{Kontrakt nie czyta prozy i nie przechodzi pusto.} Wartości ze stron HTML
        wyciągają nazwane ekstraktory z jawnym typem i kulturą; klauzula, do której nie
        trafiło żadne żądanie, jest \repopath{unexercised}, nigdy spełniona (ADR 0002).
  \item \textbf{Agent i bramka dostają te same standardy.} Standardy są spisane raz
        i kompilowane do skilla agenta oraz do konfiguracji analyzerów bramki; rozbieżność
        między tym, co usłyszał agent, a tym, co egzekwuje bramka, jest błędem CI, a nie
        znaleziskiem (rozdział~\ref{sec:standardy}).
\end{itemize}

Decyzje ADR rdzenia opisuje podrozdział~\ref{sec:rdzen-adr}.

\zrodla{\path{letsgolegacy.secondkey/docs/WORKPLAN.md},
\path{letsgolegacy.secondkey/CLAUDE.md},
\path{letsgolegacy.secondkey/docs/formats/verdict.md},
\path{letsgolegacy.secondkey/docs/analysis/phase-01.md},
\path{letsgolegacy.secondkey/docs/adr/0002-contract-language.md},
\path{letsgolegacy.secondkey/docs/adr/0003-verdict-classes.md},
\path{letsgolegacy.secondkey-standards/docs/WORKPLAN.md},
\path{letsgolegacy.secondkey-portal/docs/WORKPLAN.md}}

\subsection{Fazy i bramka fazy 01}\label{sec:architektura-fazy}

Praca jest podzielona na fazy. Faza 01 to demonstracja; jej bramką jest publiczny pakiet
dowodowy na nopCommerce 3.x, czyli ticket P9 benchu. Ścieżka krytyczna fazy 01:
S0 → S1 → S3--S6 → S7--S14 → S15 → (bench P3) → (bench P4--P8) → (bench P9). Faza 02
zaczyna się dopiero po bramce fazy 01.

\begin{xltabular}{\linewidth}{@{}l >{\raggedright\arraybackslash}X@{}}
\toprule
Faza & Zakres według planów pracy repozytoriów \\
\midrule
\endfirsthead
\toprule
Faza & Zakres według planów pracy repozytoriów \\
\midrule
\endhead
01 & Rdzeń: CLI, capture HTTP, silnik kontraktu, replay z resetem snapshotem SQL Server,
komparator, niepodpisany pakiet dowodowy, job end-to-end (S0--S15; S16 przeniesiony do
fazy 02). Portcullis: import, zmiana nazwy, reguły migracyjne, SARIF z filtrem zmienionych
linii i baseline (R0--R3); wsparcie P8 (reguły Portcullis na pull requeście kandydata
benchu) zaplanowane. Standardy: format, generator, drift check, skill dla modernize-dotnet
(C4-a--c, R4). Bench: P1--P10. legacy-risk-scan: L18. Sandbox: decyzja o topologii
demonstracyjnej D7 w stanie „proposed”, zrobiona, gdy ADR zostanie przyjęty. \\
02 & \repopath{sk mine} (miner kontraktu), \repopath{sk mutate} (silnik mutacji), capture
SQL, kolejek i PII w \path{secondkey-capture}, triage doradcza C8, pakowanie sandbox (Helm,
Hyper-V, Bicep, bundle air-gapped), reguły warstw ArchUnitNET w Portcullis, standardy
klienta jako skills, kanały wydań standardów przez NuGet i tagi git z wersją przypiętą
w każdym repozytorium (02--03), outbound stubs (S16). \\
03 & Podpis pakietu (in-toto w kopertach DSSE, cosign z kluczem klienta), biblioteka
fragmentów kontraktów, portal dowodów C11, paczki reguł per estate, pakiet Azure Managed
Application. \\
04 & Java: capture usług Java, analyzery Java w Portcullis (modernize-java). \\
\bottomrule
\end{xltabular}

\zrodla{\path{letsgolegacy.secondkey/docs/WORKPLAN.md},
\path{letsgolegacy.portcullis/docs/WORKPLAN.md},
\path{letsgolegacy.secondkey-standards/docs/WORKPLAN.md},
\path{letsgolegacy.secondkey-nopcommerce-bench/docs/WORKPLAN.md},
\path{letsgolegacy.secondkey-capture/docs/WORKPLAN.md},
\path{letsgolegacy.secondkey-sandbox/docs/WORKPLAN.md},
\path{letsgolegacy.secondkey-portal/docs/WORKPLAN.md}}

\subsection{Topologia: Windows dla legacy, Linux dla reszty}\label{sec:architektura-topologia}

System legacy na .NET Framework potrzebuje Windows: nopCommerce 3.90 to ASP.NET MVC 5
serwowany przez IIS z bazą SQL Server. Łańcuch Second Key jest napisany na .NET 10 i
działa wszędzie. Demonstracyjna topologia (decyzja D7, w stanie „proposed”) umieszcza więc
stronę legacy na jednym hoście Windows, a resztę łańcucha na Linuksie. Tylko dwa kroki
muszą sięgać do działającego sklepu przez HTTP: capture i replay. Działają obok IIS na
hoście Windows, a ich wyniki (\repopath{*.skcap}, \repopath{*.skrun}) wychodzą jako
artefakty CI do jobu linuksowego, który liczy werdykt i pakiet. Hostowane runnery GitHub
nie współdzielą sieci, więc job linuksowy i tak nie mógłby sięgnąć do IIS na runnerze
Windows.

\begin{figure}[htbp]
  \centering
  \includegraphics[width=\linewidth,height=0.6\textheight,keepaspectratio]{\dgm{a7-topologia-bench}}
  \caption{Topologia benchu: host Windows ze stroną legacy, capture i replay; Linux
           dla werdyktu i pakietu.}
  \label{fig:topologia}
\end{figure}

\zrodla{\path{letsgolegacy.secondkey-nopcommerce-bench/docs/TOPOLOGY.md},
\path{letsgolegacy.secondkey-sandbox/docs/WORKPLAN.md}}

\subsection{Znane trudne problemy}\label{sec:architektura-problemy}

Plan pracy rdzenia wymienia problemy, które są śledzone, a nie ukrywane:

\begin{itemize}
  \item \textbf{C1:} przypisanie instrukcji SQL do żądania HTTP bez zmian w kodzie to
        heurystyka z jawnie raportowanym marginesem niepewności.
  \item \textbf{C2:} fałszywe niezmienniki z małych prób; miner raportuje wsparcie
        statystyczne dla każdej klauzuli i milczy przy zbyt małej liczbie danych.
  \item \textbf{C6:} czas i losowość w kodzie legacy bez wstrzykiwania zależności są
        obsługiwane normalizacją w C7, a nie ingerencją w system legacy.
  \item \textbf{C9/C5:} czas działania na dużych solution; mutacje i analiza tylko kodu
        zmienionego w PR.
  \item \textbf{C10:} postać raportu, którą zaakceptuje audytor, jest ustalana z
        pierwszymi partnerami, a nie zgadywana.
\end{itemize}

\zrodla{\path{letsgolegacy.secondkey/docs/WORKPLAN.md}}

\clearpage
\section{Procedura pracy z frameworkiem krok po kroku}\label{sec:procedura}

Ten rozdział prowadzi przez pracę z frameworkiem od oceny systemu legacy do
opublikowanego pakietu dowodowego. Kroki odpowiadają łańcuchowi z
rozdziału~\ref{sec:architektura}. Tam, gdzie framework ma implementację referencyjną ---
skrypty benchu nopCommerce, runbook P6, tutorial Portcullis, przewodnik standardów ---
procedura przytacza ją dosłownie i odsyła do rozdziału z pełną referencją. Każdy krok
kończy się checklistą, a tam, gdzie źródła nazywają typowe awarie, także tabelą trybów
awarii (objaw → przyczyna); krok 8 ma zamiast niej tabelę decyzji.

\begin{figure}[htbp]
  \centering
  \includegraphics[height=0.85\textheight]{\dgm{a6-procedura}}
  \caption{Kroki procedury i pętle powrotne.}
  \label{fig:procedura}
\end{figure}

\begin{xltabular}{\linewidth}{@{}p{1.2cm} >{\raggedright\arraybackslash}X >{\raggedright\arraybackslash}p{3.6cm} >{\raggedright\arraybackslash}p{3.3cm}@{}}
\toprule
Krok & Cel & Wynik & Referencja \\
\midrule
\endfirsthead
\toprule
Krok & Cel & Wynik & Referencja \\
\midrule
\endhead
P & Próba łańcucha na próbkach o znanym wyniku: sklep przykładowy rdzenia, rozgrzewka benchu & pakiet dowodowy próbki & rozdz.~\ref{sec:cli}, \ref{sec:bench} \\
0 & Inwentaryzacja ryzyka technologicznego (EOL, CVE) & raport w przeglądarce & rozdz.~\ref{sec:risk-scan} \\
1 & System legacy w znanym, powtarzalnym stanie & działający system, snapshot bazy & rozdz.~\ref{sec:bench} \\
2 & Nagranie zachowania systemu legacy & \repopath{*.skcap} & rozdz.~\ref{sec:cli}, \ref{sec:formaty} \\
3 & Spisanie kontraktu: co musi się dziać, czego nie wolno & \repopath{contract.yaml} & rozdz.~\ref{sec:formaty} \\
4 & Dowód A/A: kontrakt trzyma na legacy, a łańcuch nie wnosi własnych różnic & werdykt \repopath{pass} & rozdz.~\ref{sec:cli} \\
4a & Czułość kontraktu: zabite mutanty & liczba zabitych mutantów przy M00 \repopath{pass} & rozdz.~\ref{sec:bench} \\
5 & Migracja przez agenta ze standardami & pull request kandydata & rozdz.~\ref{sec:standardy} \\
6 & Bramka na zmienionych liniach & \repopath{*.sarif} & rozdz.~\ref{sec:portcullis} \\
7 & Replay legacy kontra kandydat & \repopath{*.skrun}, \repopath{verdict.json} & rozdz.~\ref{sec:cli} \\
8 & Analiza różnic i decyzje & poprawiony kandydat albo przeglądnięta zmiana kontraktu & rozdz.~\ref{sec:formaty} \\
9 & Pakiet dowodowy & \repopath{evidence/} & rozdz.~\ref{sec:formaty} \\
10 & Publikacja wyniku i liczb & opublikowany pakiet & rozdz.~\ref{sec:bench} \\
\bottomrule
\end{xltabular}

\zrodla{\path{letsgolegacy.secondkey/README.md},
\path{letsgolegacy.secondkey-nopcommerce-bench/docs/WORKPLAN.md}}

% ---------------------------------------------------------------------------
\subsection{Role i wymagania wstępne}\label{sec:procedura-wymagania}

Źródła opisują następujących uczestników procedury:

\begin{itemize}
  \item \textbf{operator łańcucha} --- uruchamia skrypty strony legacy i polecenia
        \repopath{sk}; w benchu robi to CI, a na stacji roboczej człowiek, tymi samymi
        skryptami;
  \item \textbf{autor kontraktu} --- człowiek; w fazie 02 klauzule może proponować miner
        C2, ale klauzula decyduje dopiero po akceptacji;
  \item \textbf{operator agenta migracyjnego} --- osoba z subskrypcją GitHub Copilot,
        która uruchamia modernize-dotnet (P6 w benchu to praca człowieka);
  \item \textbf{osoba zatwierdzająca} --- ktoś, kto na podstawie pakietu dowodowego
        decyduje o zmigrowanym systemie bez uruchamiania Second Key (na przykład komitet
        zmian, druga linia obrony, audytor).
\end{itemize}

Narzędzia potrzebne w poszczególnych krokach:

\begin{xltabular}{\linewidth}{@{}>{\raggedright\arraybackslash}p{4.4cm} >{\raggedright\arraybackslash}X >{\raggedright\arraybackslash}p{2.6cm}@{}}
\toprule
Narzędzie & Do czego & Kroki \\
\midrule
\endfirsthead
\toprule
Narzędzie & Do czego & Kroki \\
\midrule
\endhead
przeglądarka; lokalnie Node 22 & strona legacy-risk-scan & 0 \\
Windows z IIS i SQL Server; Visual Studio 2022 lub Build Tools 2022 z workloadem web; git & build i hosting systemu legacy na .NET Framework (w benchu SQL Server 2022 Express) & P, 1, 2, 4, 4a, 7 \\
Windows PowerShell 5.1 lub PowerShell 7 & skrypty benchu & P, 1--9 \\
.NET 10 SDK & \repopath{sk}, Portcullis, generator standardów, build kandydata & P, 2--9 \\
GitHub Copilot z trybem agenta; Visual Studio, VS Code albo Copilot CLI z modernize-dotnet; GitHub CLI & migracja & 5 \\
przeglądarka oparta na Chromium & \repopath{report.pdf} w pakiecie & P, 9 \\
curl, python3 & \path{scripts/e2e.sh} rdzenia (oba) i sprawdzenie skrótów pakietu (python3); przy rozwoju rdzenia także \path{scripts/mutation.sh} (python3); \path{scripts/setup.sh} rdzenia ich nie sprawdza & P, 9 \\
gitleaks; w standardach i legacy-risk-scan zamiast niego działający Docker & skan sekretów przed commitem w każdym repozytorium (sekcja~\ref{sec:bezpieczenstwo-skan}) & wszystkie \\
Docker (tylko przy rozwoju rdzenia) & testy resetu SQL Server & --- \\
PSScriptAnalyzer (tylko przy rozwoju benchu) & lint skryptów benchu & --- \\
Chromium dla Playwright (tylko przy rozwoju legacy-risk-scan) & \repopath{npm run e2e} & --- \\
\bottomrule
\end{xltabular}

Skąd wziąć narzędzia --- wskazówki z plików onboardingowych repozytoriów:

\begin{xltabular}{\linewidth}{@{}>{\raggedright\arraybackslash}p{4.4cm} >{\raggedright\arraybackslash}X@{}}
\toprule
Narzędzie & Instalacja i sprawdzenie \\
\midrule
\endfirsthead
\toprule
Narzędzie & Instalacja i sprawdzenie \\
\midrule
\endhead
.NET 10 SDK & \url{https://dotnet.microsoft.com/download/dotnet/10.0}; rdzeń przypina \texttt{10.0.100} z \texttt{rollForward: latestFeature} (\path{global.json}); \texttt{./scripts/setup.sh --check} w rdzeniu tylko raportuje braki i niczego nie zmienia \\
gitleaks & \url{https://github.com/gitleaks/gitleaks/releases}; macOS: \texttt{brew install gitleaks}; Go: \texttt{go install github.com/zricethezav/gitleaks/v8@latest}; w standardach i legacy-risk-scan zastępuje go działający Docker \\
Docker & \url{https://docs.docker.com/get-docker/} \\
Node.js 22 lub nowszy & \url{https://nodejs.org/en/download}; \path{scripts/setup.sh} legacy-risk-scan odmawia starszej wersji; renderer diagramów tego dokumentu instaluje \repopath{npm ci} (sekcja~\ref{sec:rdzen-ci-pdf}) \\
Chromium dla Playwright & \texttt{npx playwright install chromium}, raz, dla \repopath{npm run e2e} w legacy-risk-scan \\
Chrome, Edge albo Chromium & PDF pakietu dowodowego; gdy przeglądarki nie ma na PATH, jej ścieżka w \repopath{SK\_CHROME\_PATH} (kolejność wyszukiwania: sekcja~\ref{sec:cli-evidence}) \\
IIS z ASP.NET 4.x, SQL Server 2022 Express & instaluje je skrypt 3 benchu; nośnik SQL Server jest przypięty w \path{pins.json}, a każde pobranie jest sprawdzane wartościami z tego pliku \\
agent modernize-dotnet & instalacja i sprawdzenie dla Visual Studio, VS Code, Copilot CLI, aplikacji GitHub Copilot i agenta w chmurze: podrozdział~\ref{sec:standardy-modernize} (krok 4) \\
\bottomrule
\end{xltabular}

\repopath{sk} i Portcullis nie mają jeszcze wydanych pakietów. Bench buduje oba z
przypiętych commitów (\path{pins.json}, sekcja \repopath{chain}) i uruchamia je jako
\repopath{dotnet <assembly>}. Ręcznie robi się to tak:

\begin{lstlisting}
git clone https://github.com/konradcinkusz/letsgolegacy.secondkey.git secondkey
git -C secondkey checkout 967c49c0005c4412180914f43437d75b43f148dc   # chain.secondKey.commit
dotnet build secondkey/src/SecondKey.Cli/SecondKey.Cli.csproj -c Release -o <work>/tools/chain/secondkey
git clone https://github.com/konradcinkusz/letsgolegacy.portcullis.git portcullis
git -C portcullis checkout 84e1925fe0d85cf23415f0d33c98590f86128722  # chain.portcullis.commit
dotnet build portcullis/src/Portcullis.Cli/Portcullis.Cli.csproj -c Release -o <work>/tools/chain/portcullis
\end{lstlisting}

Skrypty benchu znajdują narzędzia w \path{<work>\tools\chain\} albo pod ścieżkami ze
zmiennych \repopath{NOPBENCH\_SK} i \repopath{NOPBENCH\_PORTCULLIS}. W rdzeniu CLI
uruchamia się też bez instalacji:
\repopath{dotnet run --project src/SecondKey.Cli -- --help}.

\textbf{Przypięte narzędzia benchu są starsze niż commity opisane w tym dokumencie.}
\texttt{967c49c} poprzedza commit rdzenia z sekcji~\ref{sec:wprowadzenie-wersja}
(\texttt{7939e80}). Później doszły między innymi: odrzucanie przez \repopath{sk validate}
operatora jednej wartości na całej liście (\texttt{05ba3b3}), dokończenie wymian w locie
przy zatrzymaniu \repopath{sk capture} (\texttt{51e1754}), końce linii
\texttt{\textbackslash n} w każdym dokumencie JSON na każdej platformie (\texttt{b2e1bb7})
i sprawdzanie wartości \repopath{secondkey.yaml} przy wczytaniu pliku (\texttt{7939e80}).
Portcullis \texttt{84e1925} poprzedza \texttt{52d9f1e}, w którym parser skanera przeszedł
z Roslyn 4.11 na 5.9: projekt .NET 10 domyślnie kompiluje C\# 14, a starszy Roslyn czyta go
jako błędy składni (podrozdział~\ref{sec:portcullis-roslyn}). Zachowanie opisane
w rozdziałach~\ref{sec:cli} i \ref{sec:portcullis} dotyczy commitów z tabeli; bramka na
kandydacie .NET 10 (krok~6) potrzebuje Portcullis co najmniej z commita \texttt{52d9f1e}.

\zrodla{\path{letsgolegacy.secondkey-nopcommerce-bench/scripts/README.md},
\path{letsgolegacy.secondkey-nopcommerce-bench/pins.json},
\path{letsgolegacy.secondkey-nopcommerce-bench/docs/P6-RUNBOOK.md},
\path{letsgolegacy.secondkey/docs/formats/evidence.md},
\path{letsgolegacy.secondkey/docs/cli.md},
\path{letsgolegacy.secondkey/docs/analysis/phase-01.md},
\path{letsgolegacy.secondkey/global.json},
\path{letsgolegacy.secondkey/scripts/setup.sh},
\path{letsgolegacy.secondkey/scripts/e2e.sh},
\path{letsgolegacy.portcullis/src/Portcullis.Engine/Portcullis.Engine.csproj},
\path{letsgolegacy.secondkey-standards/scripts/setup.sh},
\path{letsgolegacy.secondkey-standards/scripts/hooks/pre-commit},
\path{letsgolegacy.secondkey-standards/docs/USING-WITH-MODERNIZE-DOTNET.md},
\path{letsgolegacy.legacy-risk-scan/README.md},
\path{letsgolegacy.legacy-risk-scan/scripts/setup.sh}}

% ---------------------------------------------------------------------------
\subsection{Krok przygotowawczy: łańcuch na próbkach}\label{sec:procedura-proba}

\textbf{Cel.} Zanim łańcuch spotka prawdziwy system legacy, uruchomić go raz od początku do
końca na próbce, której wynik jest znany. Źródła dają dwie takie próby.

\textbf{Sklep przykładowy rdzenia} (.NET 10 SDK, \texttt{curl}, \texttt{python3}; w klonie
\path{letsgolegacy.secondkey}):
\begin{lstlisting}
./scripts/setup.sh     # once: checks the tools, installs the secret-scanning commit hook
./scripts/e2e.sh       # the whole chain on the sample shop
\end{lstlisting}
Skrypt uruchamia sklep przykładowy dwa razy (strona legacy i kandydat), nagrywa ruch przez
proxy capture, odtwarza go na obu stronach, porównuje z \path{samples/contract.yaml}
i zapisuje pakiet dowodowy w \path{.secondkey/e2e/evidence/}. \repopath{sk compare} ma się
tu skończyć kodem 1: próbka jest zbudowana tak, żeby werdykt zawierał każdą klasę.
\texttt{SK\_PDF=required ./scripts/e2e.sh} działa jak CI: brak PDF jest błędem.
Szczegóły: podrozdział~\ref{sec:cli-e2e}.

\textbf{Rozgrzewka benchu na eShopLegacyMVC} (P3): skrypty 1--6 z \path{scripts/eshop} na
maszynie Windows z IIS, Visual Studio 2022 z workloadem web i środowiskiem .NET 10,
w PowerShell z uprawnieniami administratora, z narzędziami łańcucha zbudowanymi jak
w sekcji~\ref{sec:procedura-wymagania}. Krok 6 rozgrzewki kończy się błędem, jeśli werdykt
A/A nie jest \repopath{pass} z każdą zaakceptowaną klauzulą utrzymaną i żadną
niesprawdzoną --- dopiero po zapisaniu pakietu, więc pakiet wyjaśnia porażkę. Polecenia
i skrypty: podrozdziały~\ref{sec:bench-rozgrzewka} i \ref{sec:bench-skrypty-eshop}.

\textbf{Checklista}
\begin{checklista}
  \item \repopath{./scripts/e2e.sh} kończy się linią \texttt{e2e: ok}; pakiet jest
        w \path{.secondkey/e2e/evidence/}.
  \item Rozgrzewka: werdykt A/A \repopath{pass}, pakiet zapisany.
\end{checklista}

\zrodla{\path{letsgolegacy.secondkey/README.md},
\path{letsgolegacy.secondkey/scripts/e2e.sh},
\path{letsgolegacy.secondkey-nopcommerce-bench/README.md},
\path{letsgolegacy.secondkey-nopcommerce-bench/warmup/eshop/README.md}}

% ---------------------------------------------------------------------------
\subsection{Krok 0: ocena ryzyka technologicznego}\label{sec:procedura-krok0}

\textbf{Cel.} Zanim zapadnie decyzja o migracji, spisać, które składniki systemu są po
końcu wsparcia (EOL) albo mają znane podatności. Krok jest niezależny od łańcucha: nie
produkuje artefaktu, który czytają dalsze kroki.

\textbf{Wejście.} Plik projektu albo zależności: \path{.csproj} lub \path{.vbproj}
(SDK-style i klasyczny), \path{Directory.Packages.props}, \path{packages.config},
\path{composer.json}, \path{composer.lock}, \path{package.json}.

\textbf{Postępowanie.}
\begin{enumerate}
  \item Otwórz stronę \url{https://konradcinkusz.github.io/letsgolegacy.legacy-risk-scan/}.
        Publikuje ją \path{deploy.yml} na GitHub Pages z gałęzi \repopath{main}; według
        planu pracy ticket L18 jest scalony, a za ukończony uzna go dopiero przejście jobu
        \repopath{verify} po pierwszym udanym wdrożeniu. Lokalnie stronę uruchamia \repopath{./scripts/setup.sh}, a potem
        \repopath{npm start}; serwer udostępnia ją pod
        \url{http://127.0.0.1:4173/letsgolegacy.legacy-risk-scan/}.
  \item Wklej plik albo wybierz go z dysku; „Wypróbuj na przykładzie” pokazuje raport dla
        pliku przykładowego. Parsowanie odbywa się w przeglądarce; na zewnątrz trafiają
        tylko współrzędne pakietów (nazwa, ekosystem, wersja) --- zapytania
        \repopath{querybatch} do OSV.dev, najwyżej 1000 pakietów w każdym, potem po jednym
        \repopath{GET} na każde zwrócone advisory. Wersje Microsoft SQL Server i Windows
        Server, których plik projektu nie zawiera, wybiera się w formularzu; zmiana po
        raporcie przebudowuje go bez sieci.
  \item Przeczytaj raport: status każdej pozycji („Koniec wsparcia”, „Koniec wsparcia
        w ciągu roku”, „Znane podatności”, „Nieustalone”, „Bez zastrzeżeń”) liczony z dat
        w dniu raportu, advisories z OSV oraz datę, z której pochodzą dane EOL
        (\repopath{generatedAt}). Gdy OSV było nieosiągalne, raport ma ostrzeżenie,
        a „Sprawdź podatności ponownie” powtarza skan. Raport zapisuje przycisk „Drukuj lub
        zapisz jako PDF”.
  \item Sprawdź ograniczenia, zanim uznasz raport za pełny (tabela niżej). Ten krok przycisk
        po przycisku, z uwagami do pliku, rozwija podrozdział~\ref{sec:risk-scan-procedura}.
\end{enumerate}

\begin{xltabular}{\linewidth}{@{}>{\raggedright\arraybackslash}p{5.4cm} >{\raggedright\arraybackslash}X@{}}
\toprule
Objaw & Przyczyna \\
\midrule
\endfirsthead
\toprule
Objaw & Przyczyna \\
\midrule
\endhead
Pakiet NuGet bez advisories, choć powinien je mieć & OSV porównuje identyfikatory NuGet z uwzględnieniem wielkości liter; ręcznie edytowany \repopath{PackageReference} może mieć inną wielkość liter niż kanoniczna \\
Brak Bootstrapa lub wtyczek jQuery w raporcie z klasycznego \repopath{.csproj} & klasyczny projekt wymienia tylko pakiety z DLL; pakiety samych treści widać tylko w \repopath{packages.config} \\
Wersja w raporcie starsza niż zainstalowana & ograniczenia (\repopath{composer.json}, \repopath{package.json}, zakresy NuGet) są sprawdzane w najniższej dozwolonej wersji; dokładny wynik dają pliki lock \\
Pozycje „Nieustalone” & DLL skopiowane ręcznie, komponenty z GAC, pakiety bez wersji albo z wersją wyznaczoną zakresem nie dają się ocenić automatycznie; wymagają ręcznej weryfikacji \\
\bottomrule
\end{xltabular}

\textbf{Checklista}
\begin{checklista}
  \item Przeskanowany plik z dokładnymi wersjami (lock albo \repopath{packages.config}), jeśli istnieje.
  \item Zapisana data danych EOL z raportu.
  \item Pozycje „Nieustalone” zweryfikowane ręcznie.
\end{checklista}

\zrodla{\path{letsgolegacy.legacy-risk-scan/README.md},
\path{letsgolegacy.legacy-risk-scan/docs/WORKPLAN.md},
\path{letsgolegacy.legacy-risk-scan/src/index.html},
\path{letsgolegacy.legacy-risk-scan/src/lib/messages.js}}

% ---------------------------------------------------------------------------
\subsection{Krok 1: system legacy w znanym stanie}\label{sec:procedura-krok1}

\textbf{Cel.} Uruchomić system legacy tak, żeby każdy scenariusz --- przy nagrywaniu i
przy każdym replay, na każdej stronie --- zaczynał się od dokładnie tego samego stanu
danych i działał na dokładnie tej samej platformie. \finding{System legacy nie jest przy
tym zmieniany: żadnej zmiany kodu, rekompilacji ani modułu IIS} (zasada projektowa 4).
Implementacją referencyjną tego kroku są skrypty 1--6 benchu i początek skryptu 7.

\textbf{Zasady, które krok realizuje, i ich powody.}
\begin{enumerate}
  \item \textbf{Przypięta wersja źródeł.} Bench przypina nopCommerce
        \repopath{release-3.90} do commita \repopath{12d1f01} w \path{pins.json}; skrypt 1
        weryfikuje checkout i odmawia, gdy tag upstream się przesunął.
  \item \textbf{Powtarzalny build z manifestem.} Skrypt 2 zapisuje SHA-256 każdego
        opublikowanego pliku (\path{state\legacy-site.sha256}, zgodny z
        \repopath{sha256sum -c}) i \path{state\build.json}; nagranie ruchu podaje potem
        skrót manifestu, więc wiadomo, jaki build obserwowało.
  \item \textbf{Ustalona platforma.} To, co host dokłada do zachowania, jest przypięte
        (tabela poniżej).
  \item \textbf{Znany stan danych.} Świeża baza z danymi przykładowymi instalatora, potem
        konfiguracja potrzebna scenariuszom --- \textbf{przez UI samego systemu, nigdy
        przez SQL}, z odczytem każdej zmiany. Tak każdy wiersz, który system zapisuje, jest
        wierszem, który zapisałby sam, razem z cache i zdarzeniami.
  \item \textbf{Timery w aplikacji wyłączone przed resetem.} nopCommerce 3.90 uruchamia
        zadania harmonogramu na timerach w aplikacji webowej; restore bazy pod działającym
        zadaniem wywraca proces roboczy i sklep odpowiada 500 do czasu restartu. Bench
        wyłącza cztery zadania włączone w instalacji przykładowej i restartuje aplikację.
  \item \textbf{Snapshot bazy jako stan startowy.} Po konfiguracji powstaje snapshot SQL
        Server (\repopath{nopcommerce\_legacy\_snapshot}), przywracany przed każdym
        scenariuszem --- przy nagrywaniu i w replay po obu stronach.
\end{enumerate}

\begin{xltabular}{\linewidth}{@{}>{\raggedright\arraybackslash}p{3.2cm} >{\raggedright\arraybackslash}p{4.6cm} >{\raggedright\arraybackslash}X@{}}
\toprule
Ustawienie hosta & Wartość w benchu & Dlaczego ma znaczenie \\
\midrule
\endfirsthead
\toprule
Ustawienie hosta & Wartość w benchu & Dlaczego ma znaczenie \\
\midrule
\endhead
Collation SQL & \path{SQL_Latin1_General_CP1_CI_AS} dla serwera i bazy & porządek i równość w zapytaniach (sortowanie katalogu, wyszukiwanie) \\
Runtime .NET Framework & 4.8 (hosta), aplikacja zbudowana dla 4.5.1 & runtime in-place, tak jak na każdym produkcyjnym hoście tego sklepu \\
Kultura & \repopath{en-US}, ustawiana per żądanie z języka sklepu & formatowanie cen i dat; porównanie NLS → ICU w P7 dotyczy runtime, więc obie strony muszą widzieć tę samą kulturę \\
Strefa czasowa & UTC na hostowanych runnerach & ważność rabatów, daty zamówień, „nowe” produkty \\
Kompresja & niezainstalowana & nagrane ciała odpowiedzi zostają czytelne \\
Zegar & czas rzeczywisty & nie da się go kontrolować bez zmiany kodu legacy; komparator normalizuje czasy \\
\bottomrule
\end{xltabular}

\textbf{Polecenia (bench, host Windows).} Klon benchu:
\begin{lstlisting}
git clone https://github.com/konradcinkusz/letsgolegacy.secondkey-nopcommerce-bench.git
\end{lstlisting}
W katalogu głównym klonu, w PowerShell 5.1 albo 7:

\begin{lstlisting}
.\scripts\1-fetch-nopcommerce.ps1
.\scripts\2-build-legacy.ps1
# od tego miejsca w PowerShell z uprawnieniami administratora:
.\scripts\3-install-iis-sql.ps1
.\scripts\4-deploy-and-install.ps1
.\scripts\5-smoke.ps1
.\scripts\6-configure-store.ps1
\end{lstlisting}

Gdy polityka blokuje uruchamianie skryptów:
\begin{lstlisting}
powershell -ExecutionPolicy Bypass -File .\scripts\<script>.ps1
\end{lstlisting}
Każdy krok sprawdza swoje warunki wstępne i mówi, czego brakuje; nic, co zapisuje, nie
trafia do repozytorium. Każdy krok zapisuje wynik w \path{<work>\state\<krok>.json}, a
następny krok go czyta, gdy nie podano parametru: najpierw jawny parametr, potem plik
przekazania, potem wartość domyślna. Katalog roboczy (\repopath{-WorkRoot} albo
\repopath{NOPBENCH\_WORK}, domyślnie \path{C:\nopbench}) powinien mieć krótką ścieżkę,
bo pipeline publikacji zagnieżdża ścieżki głęboko, a starsze zadania MSBuild wciąż trafiają
w limit 260 znaków.

Pozostałe przepisy runbooka skryptów --- ponowny build bez pobierania, build od zera,
kontrola witryny z manifestem po stronie linuksowej, ponowna instalacja sklepu i wdrożenie
buildu zrobionego gdzie indziej --- podaje sekcja~\ref{sec:bench-przepisy}; parametry,
zmienne i pliki przekazania --- podrozdziały~\ref{sec:bench-stan} i \ref{sec:bench-skrypty}.

\textbf{Jedyny sekret.} Skrypt 4 tworzy administratora sklepu z hasłem generowanym na
przebieg (albo wziętym z \repopath{NOPBENCH\_ADMIN\_PASSWORD}). Hasło jest maskowane w
logach CI, nigdy nie jest drukowane i trafia tylko do
\path{<work>\secrets\legacy-admin.json}. Baza nie ma hasła: witryna łączy się tożsamością
puli aplikacji (uwierzytelnianie Windows).

\textbf{Stacja robocza zamiast CI.} Te same skrypty działają na maszynie wirtualnej
Windows Server 2022 pod Hyper-V (wersja ewaluacyjna wystarcza; 4 vCPU i 8~GB są wygodne)
albo na Windows 10/11 Pro lub Enterprise z IIS. Żeby strona linuksowa (WSL2 albo maszyna
wirtualna) dotarła do sklepu, trzeba otworzyć port na hoście Windows; potem smoke test
działa także z Linuksa:

\begin{lstlisting}
New-NetFirewallRule -DisplayName 'nopbench legacy 8080' -Direction Inbound -Protocol TCP -LocalPort 8080 -Action Allow
pwsh scripts/5-smoke.ps1 -BaseUrl http://<windows-host>:8080/
\end{lstlisting}

\begin{xltabular}{\linewidth}{@{}>{\raggedright\arraybackslash}p{5.6cm} >{\raggedright\arraybackslash}X@{}}
\toprule
Objaw & Przyczyna i działanie \\
\midrule
\endfirsthead
\toprule
Objaw & Przyczyna i działanie \\
\midrule
\endhead
Każda kwota podatku i wysyłki w nagraniu wynosi \$0.00 & pominięta konfiguracja sklepu: dane przykładowe nie mają stawek \\
Pierwsze żądania po restore dostają 500, potem sklep wraca & restore pod działającym zadaniem harmonogramu; wyłącz zadania i zrestartuj aplikację (skrypt 6) \\
Skrypt 6 kończy się błędem na sklepie już skonfigurowanym & oczekuje sklepu w stanie po instalacji; najpierw ponowna instalacja (skrypt 4) \\
Build się nie udaje & otwórz \path{<work>\logs\build.binlog} (albo \path{publish.binlog}) w MSBuild Structured Log Viewer; CI publikuje te pliki jako artefakt \repopath{legacy-build-logs} \\
Dowolna awaria hosta & \path{.\scripts\collect-diagnostics.ps1} zbiera logi IIS, dziennika zdarzeń, SQL Server i nopCommerce do \path{<work>\diagnostics} (w CI: artefakt \repopath{legacy-iis-diagnostics}) \\
\bottomrule
\end{xltabular}

\textbf{Checklista}
\begin{checklista}
  \item Źródła legacy przypięte do commita i zweryfikowane.
  \item Build z manifestem SHA-256; skrót manifestu zapisany.
  \item Collation, runtime, kultura, strefa czasowa i kompresja ustalone i zapisane.
  \item Konfiguracja wykonana przez UI systemu, każda zmiana odczytana.
  \item Timery działające na bazie wyłączone, aplikacja zrestartowana.
  \item Smoke test przechodzi.
\end{checklista}

\zrodla{\path{letsgolegacy.secondkey-nopcommerce-bench/scripts/README.md},
\path{letsgolegacy.secondkey-nopcommerce-bench/docs/TOPOLOGY.md},
\path{letsgolegacy.secondkey-nopcommerce-bench/docs/P4-TRAFFIC-PLAN.md},
\path{letsgolegacy.secondkey-nopcommerce-bench/docs/adr/0005-nopcommerce-traffic-recording.md},
\path{letsgolegacy.secondkey-nopcommerce-bench/results/p4/README.md},
\path{letsgolegacy.secondkey/README.md},
\path{letsgolegacy.secondkey/docs/WORKPLAN.md}}

% ---------------------------------------------------------------------------
\subsection{Krok 2: nagranie ruchu}\label{sec:procedura-krok2}

\textbf{Cel.} Zapisać, jak system legacy naprawdę się zachowuje, jako plik
\repopath{*.skcap}: zestaw scenariuszy, z których każdy jest jedną sesją, nagrany przez
proxy \repopath{sk capture} stojące przed systemem.

\textbf{Planowanie scenariuszy.} Bench spisuje zestaw w planie ruchu (P4-TRAFFIC-PLAN),
zanim cokolwiek nagra: obszary (katalog, koszyk, rabaty, podatki, wysyłka, checkout,
zamówienia), persony, kroki i regułę biznesową, którą każdy scenariusz ćwiczy --- z
odwołaniem do kodu systemu legacy, który tę regułę realizuje. Reguły opisują, jak
zachowuje się \emph{ten} system, a nie jak powinien zachowywać się sklep w ogóle. Zasady
nagrywania z planu:

\begin{enumerate}
  \item \textbf{Każdy scenariusz startuje ze snapshotu} --- scenariusze są niezależne od
        siebie i od kolejności, więc każdy da się nagrać i odtworzyć osobno.
  \item \textbf{Tylko przez proxy.} Żądanie, które omija proxy, nie jest dowodem.
  \item \textbf{Jeden scenariusz --- jedna sesja}: własny cookie jar i nagłówek
        \repopath{x-bench-scenario: <id>} w każdym żądaniu. Ten nagłówek jest session
        key capture, więc scenariusz jest jedną sesją w \repopath{*.skcap} i jednym
        scenariuszem w replay. Historia z dwiema osobami to dwa scenariusze.
  \item \textbf{Jak przeglądarka}: linki i formularze samego systemu, każdy formularz z
        wszystkimi polami, token anti-forgery czytany ze strony, która go wydała.
  \item \textbf{Sondy (probes).} Język kontraktu zawęża klauzulę tylko metodą, ścieżką i
        query, więc żądanie, na którego odpowiedzi kontrakt stwierdza fakt właściwy
        jednemu scenariuszowi, dostaje w query \repopath{probe=<nazwa>}
        (\repopath{/cart?probe=T11}); system ignoruje nieznany parametr, a parametr tylko
        nazywa żądanie.
  \item \textbf{Deterministyczne dane wejściowe}: stałe produkty, ilości, adresy, kody
        kuponów, klienci; rejestracja zawsze tym samym adresem (snapshot sprawia, że jest
        wolny).
  \item \textbf{Brak wywołań wychodzących}: tylko płatności i wysyłka offline, więc
        outbound stubs (S16) nie są potrzebne.
\end{enumerate}

\textbf{Polecenie.} Ogólna postać (szczegóły opcji: rozdział~\ref{sec:cli}):

\begin{lstlisting}
sk capture --target http://127.0.0.1:5080 --listen http://127.0.0.1:8080 --session-key ASP.NET_SessionId
\end{lstlisting}

Plik nagrania (\repopath{--out}, domyślnie \path{.secondkey/traffic.skcap}) jest tworzony od
nowa przy każdym starcie: \repopath{sk capture} bez pytania nadpisuje nagranie o tej samej
ścieżce (sekcja~\ref{sec:cli-wywolanie}).

W benchu robi to skrypt 7, a proxy słucha na \repopath{http://127.0.0.1:8000/} i przekazuje
do \repopath{http://localhost:8080/}:

\begin{lstlisting}
.\scripts\7-record-traffic.ps1              # cały zestaw
.\scripts\7-record-traffic.ps1 -Only T14    # jeden scenariusz; nagranie oznaczone jako częściowe
\end{lstlisting}

Skrypt uruchamia \repopath{sk capture} z opcjami \repopath{--listen}, \repopath{--target},
\repopath{--out} i \repopath{--session-key}, czeka na linię
\repopath{sk capture: recording}, prowadzi scenariusze, a na końcu czeka, aż plik będzie
zawierał tyle wymian, ile żądań wysłano, i dopiero wtedy zatrzymuje proces. Jeśli liczba
się nie zgadza, nagranie jest odrzucane: żądanie bez odpowiedzi legacy nie zostało
zapisane, a żądanie od kogoś innego „nie jest nasze”. Scenariusz, który się nie powiódł,
nie zatrzymuje pozostałych; skrypt drukuje log nopCommerce dla nieudanego scenariusza
(następny restore go skasuje) i kończy się błędem.

\textbf{Zatrzymanie capture.} Ctrl+C albo SIGTERM zatrzymuje proxy czysto: nie przyjmuje
nowych żądań, pozwala dokończyć żądania w locie (do 30 sekund), zamyka plik, drukuje
podsumowanie i kończy się kodem 0:

\begin{lstlisting}
sk capture: recorded 290 exchange(s), skipped 0 the legacy system did not answer.
\end{lstlisting}

Żądanie, na które system legacy nie odpowiedział, dostaje od proxy 502 albo 504 i
\textbf{nie jest nagrywane}, więc w nagraniu brakuje go w scenariuszu --- skryptowane
nagranie musi mieć \repopath{skipped 0}. Każda wymiana jest zapisywana i spłukiwana na
dysk od razu, więc zabity proces zachowuje to, co nagrał, ale zabicie w trakcie zapisu
zostawia uciętą ostatnią linię, którą \repopath{sk validate} odrzuca. Tam, gdzie pipeline
może proces tylko zabić, podaje się \repopath{--exit-after <n>} z liczbą wymian, które
skrypt wyśle.

\textbf{Po nagraniu.}
\begin{lstlisting}
sk validate traffic.skcap
\end{lstlisting}
Zapisz liczbę scenariuszy i skrót manifestu buildu legacy razem z nagraniem (bench:
\path{state\traffic.json} --- ścieżka i skrót nagrania, lista scenariuszy z
identyfikatorami sesji i liczbą żądań, snapshot, odpowiedzi na pierwsze żądanie po każdym
restore, skrót manifestu witryny).

\textbf{Prywatność nagrania.} Capture redaguje nagłówki przed zapisem: zawsze
\repopath{authorization}, \repopath{proxy-authorization}, \repopath{cookie},
\repopath{set-cookie} i \repopath{x-api-key}, oraz każdy dodany opcją
\repopath{--redact-header}. Ciała ponad \repopath{--max-body-bytes} (domyślnie 1~MiB) są
ucinane i oznaczane jako obcięte. Redakcja dotyczy tylko nagłówków: ciała żądań
i odpowiedzi, także pola formularzy, trafiają do nagrania
(sekcja~\ref{sec:bezpieczenstwo-dane}). Dlatego bench konfiguruje sklep z pominięciem proxy
nagrywającego --- hasło administratora nigdy nie może trafić do nagrania --- a scenariusze
logują się wyłącznie kontami sklepu przykładowego. Nagrania systemów rzeczywistych nigdy
nie trafiają do git (\repopath{.gitignore} wyklucza \repopath{*.skcap} i \repopath{*.skrun}
poza syntetycznymi próbkami).

\begin{xltabular}{\linewidth}{@{}>{\raggedright\arraybackslash}p{5.6cm} >{\raggedright\arraybackslash}X@{}}
\toprule
Objaw & Przyczyna \\
\midrule
\endfirsthead
\toprule
Objaw & Przyczyna \\
\midrule
\endhead
Replay rozjeżdża się od pierwszego scenariusza & replay nie startował ze snapshotu albo nagranie ominęło proxy \\
Liczba wymian w pliku mniejsza niż liczba wysłanych żądań & system legacy nie odpowiedział na część żądań (proxy zwróciło 502/504 i pominęło je) \\
\repopath{sk validate} odrzuca ostatnią linię & proces zabity w trakcie zapisu; użyj \repopath{--exit-after} albo zatrzymaj sygnałem \\
Reguła \repopath{never} przechodzi, niczego nie sprawdzając & scenariusz nie dotarł do tego, czego reguła pilnuje; każda reguła \repopath{never} potrzebuje pozytywnego bliźniaka w tym samym scenariuszu \\
Identyfikatory zamówień różnią się między stronami & scenariusze startowały z różnych snapshotów albo dodatkowe żądanie utworzyło zamówienie \\
Absolutne URL-e w odpowiedziach wskazują port systemu, nie proxy & capture wysyła do systemu nagłówek host celu (domyślne zachowanie YARP); skrypty scenariuszy używają wyłącznie ścieżek względnych, więc nic nie omija nagrania \\
\bottomrule
\end{xltabular}

\textbf{Checklista}
\begin{checklista}
  \item Konfiguracja wykonana, potem snapshot; każdy scenariusz startuje ze snapshotu.
  \item Każde żądanie każdego scenariusza przeszło przez proxy.
  \item Każda persona ma własny cookie jar; tokeny anti-forgery czytane ze stron.
  \item Każda reguła planowana do kontraktu jest ćwiczona przez nagrany scenariusz, a
        każda reguła \repopath{never} ma pozytywnego bliźniaka.
  \item \repopath{skipped 0}; liczba wymian równa liczbie wysłanych żądań;
        \repopath{sk validate} przechodzi.
  \item Liczba scenariuszy i skrót manifestu legacy zapisane razem z nagraniem.
  \item Sondy NLS → ICU z planu ruchu są w nagraniu.
\end{checklista}

\zrodla{\path{letsgolegacy.secondkey-nopcommerce-bench/docs/P4-TRAFFIC-PLAN.md},
\path{letsgolegacy.secondkey-nopcommerce-bench/docs/TOPOLOGY.md},
\path{letsgolegacy.secondkey-nopcommerce-bench/scripts/README.md},
\path{letsgolegacy.secondkey-nopcommerce-bench/scripts/lib/NopBench.Chain.ps1},
\path{letsgolegacy.secondkey/docs/cli.md},
\path{letsgolegacy.secondkey/docs/formats/README.md},
\path{letsgolegacy.secondkey/docs/formats/contract.md},
\path{letsgolegacy.secondkey/src/SecondKey.Artifacts/Capture/CaptureWriter.cs},
\path{letsgolegacy.secondkey-nopcommerce-bench/scripts/6-configure-store.ps1},
\path{letsgolegacy.secondkey-nopcommerce-bench/scripts/lib/NopBench.Shop.ps1},
\path{letsgolegacy.secondkey/.gitignore}}

% ---------------------------------------------------------------------------
\subsection{Krok 3: kontrakt}\label{sec:procedura-krok3}

\textbf{Cel.} Spisać w \repopath{contract.yaml}, co w nagranym zachowaniu jest regułą
biznesową: co musi się dziać (klauzule \repopath{must}) i czego nigdy nie wolno
(klauzule \repopath{never}), oraz które różnice nie mają znaczenia (normalizacja). Pełna
referencja języka: rozdział~\ref{sec:formaty}.

\textbf{Postępowanie.}
\begin{enumerate}
  \item \textbf{Weź reguły z planu ruchu.} Każda reguła ma identyfikator, rodzaj
        (\repopath{must}/\repopath{never}), scenariusze, które ją ćwiczą, i miejsce w kodzie
        legacy, z którego pochodzi (w benchu R01--R46).
  \item \textbf{Wartości ze stron HTML wyciągaj ekstraktorami.} Klauzule nigdy nie czytają
        prozy: wartość, która żyje na stronie HTML, najpierw wyciąga ekstraktor (selektor
        CSS, JSON, nagłówek albo wyrażenie regularne z jedną grupą). Przy domyślnym
        \repopath{compare.html: extracts} strona HTML jest porównywana między wersjami
        wyłącznie przez wyciągnięte wartości --- markup nie jest zachowaniem. Gdy
        ekstraktor ma \repopath{all: true}, lista jest całością: liczy się ją
        (\repopath{countAtLeast} i pokrewne), przeszukuje (\repopath{contains}) albo testuje
        wartość po wartości przez \repopath{[*]}; operator jednej wartości
        (\repopath{gt}, \repopath{between}, \repopath{matches}\ldots{}) na całej liście
        \repopath{sk validate} odrzuca. Ekstraktor deklaruje typ (\repopath{as}:
        \repopath{string}, \repopath{decimal}, \repopath{integer}, \repopath{boolean})
        i kulturę, w której strona zapisuje liczby (\repopath{culture}, domyślnie
        niezmienna). Wartość, której ekstraktor nie odczyta, liczy się przeciw stronie, na
        której wystąpiła.
  \item \textbf{Zawężaj klauzule} przez \repopath{when} (metoda, ścieżka, query) --- dla
        faktów jednego scenariusza przez sondę \repopath{probe=<nazwa>}.
  \item \textbf{Pisz klauzule nieobecności.} Kontrakt, który mówi tylko, co musi się
        dziać, powiedział połowę. Bench celuje w około 20\,\% klauzul \repopath{never}
        (wynik: 15 z 76, 19,7\,\%). \finding{Każda klauzula \repopath{never} potrzebuje
        pozytywnego bliźniaka w tym samym scenariuszu}, inaczej może przejść bez
        sprawdzenia czegokolwiek.
  \item \textbf{Normalizuj tylko to, co różni się między dwoma przebiegami tego samego
        systemu}: ciasteczka i GUID-y klientów, tokeny anti-forgery, ciasteczka sesji i
        uwierzytelniania, GUID-y zamówień i znaczniki czasu, nagłówki transportowe
        wynikające z ciał. Nie normalizuj tego, co musi się zgadzać --- w benchu numery
        zamówień, które z tego samego snapshotu i tych samych żądań muszą być równe po obu
        stronach.
  \item \textbf{Nadaj pochodzenie} (\repopath{provenance}): klauzula \repopath{proposed}
        jest liczona i pokazywana, ale o niczym nie decyduje; decyduje dopiero
        \repopath{accepted}.
  \item \textbf{Podbij \repopath{metadata.revision}} przy każdej zmianie klauzuli.
  \item \textbf{Zwaliduj}:
\begin{lstlisting}
sk validate contract.yaml secondkey.yaml
\end{lstlisting}
        Błąd walidacji wskazuje miejsce (ścieżkę w dokumencie) i powód. Walidacja
        \repopath{secondkey.yaml} sprawdza wyrażenia korelacji (muszą się kompilować i mieć grupę),
        ale nie zmienne z \repopath{connectionStringEnv}; ich brak zgłasza dopiero
        \repopath{sk replay} (sekcja~\ref{sec:cli-konfiguracja}).
\end{enumerate}

Plan ruchu benchu zawiera też sondy NLS → ICU dla P7: miejsca w nopCommerce, w których
dane kulturowe docierają do klienta (formatowanie ujemnych kwot walutowych, zamiana
wielkości liter i przycinanie kodów kuponów, porównanie kodów pocztowych, sortowanie
nazw, daty i godziny). Źródło traktuje je jako hipotezy do potwierdzenia albo obalenia w
P7, a nie jako znaleziska.

\begin{xltabular}{\linewidth}{@{}>{\raggedright\arraybackslash}p{5.6cm} >{\raggedright\arraybackslash}X@{}}
\toprule
Objaw & Przyczyna i działanie \\
\midrule
\endfirsthead
\toprule
Objaw & Przyczyna i działanie \\
\midrule
\endhead
\repopath{sk validate}: \texttt{the contract has no accepted 'never' clause} & kontrakt musi zawierać co najmniej jedną zaakceptowaną klauzulę \repopath{never} \\
\repopath{sk validate} odrzuca \repopath{gt}, \repopath{between}, \repopath{matches}\ldots{} na ekstraktorze z \repopath{all: true} & operator jednej wartości nigdy nie decyduje o wartościach listy; użyj ścieżki z \repopath{[*]}, którą podaje komunikat \\
Błąd YAML przy \repopath{select} ze ścieżką w nawiasach kwadratowych & w mapowaniu w nawiasach klamrowych \texttt{[} i \texttt{]} są składnią; ścieżkę zapisuje się w cudzysłowie: \texttt{\{ select: "response.headers.location[0]", \ldots{} \}} \\
Klauzula na ekstraktorze regresuje jako błąd, choć wartość jest na stronie (np. \texttt{(\$16.00)} jako \repopath{decimal} w \repopath{en-US}) & ekstraktor nie odczytał wartości, a wartość nieczytelna liczy się przeciw stronie, na której wystąpiła; sprawdź \repopath{as} i \repopath{culture} \\
Różnica w dacie mimo \repopath{masks.timestamps} & maska obejmuje tylko daty ISO-8601; inne zapisy (\texttt{/Date(\ldots)/}, \texttt{26.09.2026}) zostają celowo --- dodaj maskę \repopath{patterns}, jeśli zapis legacy ma być ukryty po obu stronach \\
\bottomrule
\end{xltabular}

Pełna lista komunikatów walidacji kontraktu: podrozdział~\ref{sec:formaty-bledy}.

\textbf{Checklista}
\begin{checklista}
  \item Każda klauzula ma \repopath{why} i pochodzenie; reguła w planie wskazuje kod legacy.
  \item Udział klauzul \repopath{never} około 20\,\%; każda ma pozytywnego bliźniaka.
  \item Normalizacja obejmuje tylko wartości zmienne między przebiegami tego samego systemu.
  \item \repopath{sk validate} przechodzi dla kontraktu i \repopath{secondkey.yaml}.
\end{checklista}

\zrodla{\path{letsgolegacy.secondkey/docs/formats/contract.md},
\path{letsgolegacy.secondkey/docs/adr/0002-contract-language.md},
\path{letsgolegacy.secondkey/docs/formats/verdict.md},
\path{letsgolegacy.secondkey-nopcommerce-bench/docs/P4-TRAFFIC-PLAN.md},
\path{letsgolegacy.secondkey-nopcommerce-bench/contract/README.md},
\path{letsgolegacy.secondkey-nopcommerce-bench/results/p5/README.md},
\path{letsgolegacy.secondkey-nopcommerce-bench/docs/P9-MUTANTS.md},
\path{letsgolegacy.secondkey/src/SecondKey.Artifacts/Contracts/ContractFile.cs},
\path{letsgolegacy.secondkey/docs/cli.md}}

% ---------------------------------------------------------------------------
\subsection{Krok 4: dowód A/A na systemie legacy}\label{sec:procedura-krok4}

\textbf{Cel.} Zanim kontrakt oceni kandydata, musi przejść na samym systemie legacy:
replay nagrania na legacy po obu stronach (A/A), z przywróceniem snapshotu przed każdym
scenariuszem na każdej stronie, ma dać werdykt \repopath{pass} z każdą zaakceptowaną
klauzulą utrzymaną i żadną niesprawdzoną. Taki przebieg pokazuje, że łańcuch nie wnosi
różnic od siebie i że kontrakt trzyma na systemie legacy.

\textbf{Polecenia (bench).}
\begin{lstlisting}
.\scripts\8-replay.ps1       # sk replay, A/A: obie strony to legacy
.\scripts\9-verdict.ps1      # sk compare, sk evidence
\end{lstlisting}

Skrypt 8 uruchamia
\begin{lstlisting}
sk replay --config contract/secondkey.yaml --capture <nagranie> --out <przebieg> --legacy <url> --candidate <url>
\end{lstlisting}
z tym samym adresem po obu stronach, a reset
\repopath{sqlServerSnapshot} z \repopath{secondkey.yaml} przywraca snapshot przed każdym
scenariuszem. Łańcuch połączenia dla resetu skrypt ustawia tylko dla procesu \repopath{sk}
w zmiennej \repopath{NOPBENCH\_SK\_SQL} (uwierzytelnianie Windows, bez sekretu). Przebieg
z wynikiem bez odpowiedzi albo z nieudanym resetem jest odrzucany jako nieuczciwe
porównanie. Skrypt 9 waliduje kontrakt i przebieg, uruchamia \repopath{sk compare}, waliduje
werdykt, składa pakiet \repopath{sk evidence} i sprawdza skrót SHA-256 każdego pliku z
manifestu; przebieg A/A musi przejść z każdą zaakceptowaną klauzulą utrzymaną. Skrypt 9
działa wszędzie, gdzie działają narzędzia łańcucha --- w CI na Linuksie
(\texttt{pwsh scripts/9-verdict.ps1}). Gdy A/A nie przechodzi, kończy się błędem dopiero po
zapisaniu pakietu, żeby pakiet wyjaśniał porażkę, a jego log wymienia każdą klauzulę, która
nie trzyma, żądania, na których zawiodła, i tekst odpowiedzi legacy
(podrozdział~\ref{sec:bench-krok7}).

\textbf{Co A/A znajduje.} W benchu pierwszy przebieg tej bramki nie przeszedł, a obie
przyczyny leżały po stronie kontraktu:
\begin{itemize}
  \item \textbf{Kontrakt mylił się co do systemu legacy.} Klauzula o stronie głównej
        oczekiwała sześciu produktów, a dane przykładowe oznaczają cztery; klauzula
        złamana po obu stronach dostaje status \repopath{violated-both}. Poprawka:
        klauzula nazywa teraz cztery produkty.
  \item \textbf{Niedeterminizm do maskowania.} 17 odpowiedzi JSON różniło się między dwoma
        replay tego samego żądania: świeży token anti-forgery w podsumowaniu zamówienia i
        obraz pozycji koszyka z cache, którego restore bazy nie resetuje. Oba przypadki
        maskuje teraz normalizacja (tylko wartości) i kontrakt mówi dlaczego.
\end{itemize}

\textbf{Pętla.} Różnica, której nic w kontrakcie nie tłumaczy, to \repopath{regression}
(\repopath{uncovered-difference}), także w A/A. Rozwiązuje się ją przeglądniętą zmianą
kontraktu: regułą normalizacji, jeśli różnica nie ma znaczenia, klauzulą, jeśli ma. Wróć
do kroku~3 i powtarzaj, aż A/A da \repopath{pass}.

\begin{xltabular}{\linewidth}{@{}>{\raggedright\arraybackslash}p{5.6cm} >{\raggedright\arraybackslash}X@{}}
\toprule
Objaw & Przyczyna i działanie \\
\midrule
\endfirsthead
\toprule
Objaw & Przyczyna i działanie \\
\midrule
\endhead
Klauzula \repopath{violated-both} & kontrakt myli się co do systemu legacy; popraw klauzulę \\
\repopath{uncovered-difference} w A/A & wartość zmienna między przebiegami; maska z uzasadnieniem albo klauzula \\
Klauzula \repopath{unexercised} & żadne żądanie nie trafiło w jej \repopath{when}; nigdy nie jest liczona jako spełniona --- popraw zakres albo dograj scenariusz \\
\repopath{sk replay} kończy się kodem 4 & reset strony się nie udał; scenariusze tej strony nie zostały odtworzone (\repopath{reset-failed}); przebieg nie jest uczciwym porównaniem \\
\bottomrule
\end{xltabular}

\textbf{Checklista}
\begin{checklista}
  \item Snapshot przywracany przed każdym scenariuszem po obu stronach; zero nieudanych resetów.
  \item Zero wyników bez odpowiedzi; liczba scenariuszy w przebiegu równa liczbie w nagraniu.
  \item Werdykt \repopath{pass}: każda zaakceptowana klauzula utrzymana, zero niesprawdzonych.
  \item Każda maska ma zapisane uzasadnienie.
\end{checklista}

\zrodla{\path{letsgolegacy.secondkey-nopcommerce-bench/scripts/8-replay.ps1},
\path{letsgolegacy.secondkey-nopcommerce-bench/scripts/9-verdict.ps1},
\path{letsgolegacy.secondkey-nopcommerce-bench/results/p5/README.md},
\path{letsgolegacy.secondkey-nopcommerce-bench/contract/README.md},
\path{letsgolegacy.secondkey/docs/formats/verdict.md},
\path{letsgolegacy.secondkey/docs/cli.md}}

% ---------------------------------------------------------------------------
\subsection{Krok 4a: czułość kontraktu --- zabite mutanty}\label{sec:procedura-mutanty}

\textbf{Cel.} Pokazać, że kontrakt zauważa zmianę zachowania: do kodu legacy wprowadza się
celowo małe defekty (mutanty) i sprawdza, czy porównanie je wykrywa. Komponent C9
(\repopath{sk mutate}) należy do fazy 02 i w fazie 01 kończy się kodem 70; bench zastępuje
go zestawem ręcznie napisanych łatek, które przechodzą przez te same \repopath{sk replay}
i \repopath{sk compare} co kandydat. Pełna referencja: podrozdział~\ref{sec:bench-p9}.

\textbf{Zasady.}
\begin{enumerate}
  \item \textbf{Najpierw rejestracja.} Łatki, katalog i klauzule, które mają złapać każdego
        mutanta, trafiają do repozytorium jednym commitem, zanim którykolwiek się wykona;
        oczekiwań nie zmienia się po zobaczeniu wyniku.
  \item \textbf{Zabity} znaczy: \repopath{sk compare} legacy kontra mutant daje
        \repopath{fail}. Wynik \repopath{pass} albo \repopath{review} znaczy, że mutant
        przeżył.
  \item \textbf{M00 kalibruje metodę}: przypięte źródło bez zmian, zbudowane i wdrożone jak
        każdy mutant, musi dać \repopath{pass}; przebieg, w którym M00 nie przechodzi, nie
        liczy żadnego wyniku.
  \item \textbf{Ocalałego mutanta nie zabija się zmianą kontraktu} w tym samym zadaniu ---
        to znalezisko dla następnej rewizji kontraktu.
\end{enumerate}

\textbf{Postępowanie (bench).} Jednego mutanta uruchamia się ręcznie na hoście Windows po
skryptach 1--6 benchu: build skryptami 1 i 2 z \repopath{-Patch} we własnym katalogu
roboczym, wdrożenie \path{scripts/mutants/2-deploy-mutant.ps1} jako druga witryna IIS
\textbf{przed skryptem 7} (snapshot, z którego startuje replay, musi już zawierać
użytkownika bazy witryny mutanta), potem nagranie (skrypt 7), replay legacy kontra mutant
(skrypt 8 z \texttt{-Candidate http://localhost:8091/}), werdykt (skrypt 9) i wynik
\path{scripts/mutants/3-mutant-outcome.ps1}. Polecenia i tryby awarii podaje
podrozdział~\ref{sec:bench-p9} („Jeden mutant ręcznie”,
tabela~\ref{tab:bench-mutanty-awarie}), skrypty mutantów opisuje
sekcja~\ref{sec:bench-skrypty-mutanty}. Cały zestaw uruchamia workflow \path{mutants.yml},
ręcznie albo przy pull requeście, który dotyka \path{mutants/}, \path{scripts/mutants/} lub
samego workflowu.

\textbf{Checklista}
\begin{checklista}
  \item Katalog mutantów i oczekiwane klauzule zacommitowane przed pierwszym przebiegiem.
  \item Mutant wdrożony przed nagraniem (skrypt 7).
  \item M00 daje \repopath{pass}; każdy mutant ma wynik.
  \item Ocalały mutant zapisany jako znalezisko dla następnej rewizji kontraktu.
\end{checklista}

\zrodla{\path{letsgolegacy.secondkey-nopcommerce-bench/docs/P9-MUTANTS.md},
\path{letsgolegacy.secondkey-nopcommerce-bench/scripts/README.md},
\path{letsgolegacy.secondkey-nopcommerce-bench/docs/adr/0007-p9-mutants-on-a-second-site.md},
\path{letsgolegacy.secondkey-nopcommerce-bench/results/p9/README.md},
\path{letsgolegacy.secondkey-nopcommerce-bench/docs/TOPOLOGY.md},
\path{letsgolegacy.secondkey/docs/cli.md}}

% ---------------------------------------------------------------------------
\subsection{Krok 5: migracja przez agenta ze standardami}\label{sec:procedura-krok5}

\textbf{Cel.} Uzyskać kandydata --- zmigrowany system --- jako pull request, którego diff
jest dokładnie migracją, wykonaną przez agenta modernize-dotnet ze standardami
frameworku załadowanymi jako skill, z pełnym zapisem przebiegu. W benchu to ticket P6,
praca człowieka z subskrypcją GitHub Copilot; kryterium odbioru: kandydat się buduje, a
istniejące testy przechodzą. Poniższe kroki pochodzą z runbooka P6 benchu i przewodnika
standardów (rozdział~\ref{sec:standardy}).

\textbf{5.1 Wymagania.} GitHub Copilot (Pro, Business lub Enterprise) z trybem agenta;
jeden klient z agentem modernizacji: Visual Studio 2026 albo 2022 w wersji 17.14.17 lub
nowszej (runbook P6 podaje ogólniej 17.14+) z \emph{GitHub Copilot app modernization},
VS Code z rozszerzeniem albo GitHub Copilot CLI --- instalację i sprawdzenie w każdym
kliencie opisuje podrozdział~\ref{sec:standardy-modernize} (krok 4); Windows z MSBuild dla
.NET Framework, .NET 10 SDK, git i GitHub CLI; prawo zapisu do repozytorium benchu, do
którego trafia zapis przebiegu. Standardy: runbook P6 wymaga otagowanego wydania
repozytorium standardów z \path{generated/skills/} i bez niego P6 się nie zaczyna;
przewodnik standardów poza benchem dopuszcza przed pierwszym wydaniem \repopath{main}
z zapisanym commitem. Uruchom stoper: wszystko do otwarcia pull requestu liczy się do
godzin P6.

\textbf{5.2 Fork i gałęzie.} Kandydat jest zmodyfikowanym nopCommerce, więc mieszka w
forku upstream, a nie w repozytorium benchu:

\begin{lstlisting}
gh repo fork nopSolutions/nopCommerce --clone=false
git init nop-p6 && cd nop-p6
git config core.autocrlf false
git remote add upstream https://github.com/nopSolutions/nopCommerce.git
git remote add fork https://github.com/<your-account>/nopCommerce.git
git fetch --depth 1 upstream tag release-3.90
git rev-parse "release-3.90^{commit}"          # must print 12d1f0139149a9d6e84964d325d62bae6ac748f6
git checkout -b secondkey/p6-base 12d1f0139149a9d6e84964d325d62bae6ac748f6
\end{lstlisting}

Gałąź \repopath{secondkey/p6-base} to przypięty commit plus jeden commit ze skillem i nie
zmienia się po starcie P6; \repopath{secondkey/p6-candidate} powstaje z niej i zbiera
wszystkie zmiany agenta i człowieka. Pull request idzie z kandydata do bazy, więc jego
diff to dokładnie migracja.

\textbf{5.3 Standardy jako skill.} Skill kopiuje się z wydania standardów do
\path{.github/skills/} repozytorium docelowego i commituje osobno na gałęzi bazowej:

\begin{lstlisting}
git clone --depth 1 --branch <standards-tag> https://github.com/konradcinkusz/letsgolegacy.secondkey-standards.git ../standards
mkdir -p .github/skills
cp -R ../standards/generated/skills/. .github/skills/
git -C ../standards rev-parse HEAD             # commit standardów do zapisu przebiegu
git add .github/skills
git commit -m "Add Second Key standards skills (<standards-tag>, <standards-commit>)"
git push -u fork secondkey/p6-base
git checkout -b secondkey/p6-candidate
git push -u fork secondkey/p6-candidate
\end{lstlisting}

Kopii skilla się nie edytuje --- zmiana standardu to zmiana w repozytorium standardów.
Sprawdzenie: istnieje
\path{.github/skills/applying-dotnet-migration-standards/SKILL.md}, a
\repopath{metadata.version} w jego front matter to wersja, którą miałeś wziąć.

\textbf{5.4 Bazowa linia testów legacy.} „Istniejące testy przechodzą” wymaga punktu
odniesienia: które testy przechodzą na niezmienionym kodzie. nopCommerce 3.90 ma pięć
projektów testowych NUnit 3 w \path{src/Tests}. Bench buduje je skryptami 1 i 2, a potem:

\begin{lstlisting}
nuget install NUnit.ConsoleRunner -Version 3.22.0 -OutputDirectory C:\nopbench\tools
$runner = 'C:\nopbench\tools\NUnit.ConsoleRunner.3.22.0\tools\nunit3-console.exe'
& $runner (Get-ChildItem C:\nopbench\legacy-src\src\Tests\*\bin\Release\*.Tests.dll).FullName --result=C:\nopbench\p6-legacy-tests.xml
\end{lstlisting}

Zapisz liczby (wszystkie, zaliczone, niezaliczone, pominięte) per assembly. Test, który
nie przechodzi na legacy, nie jest regresją, gdy nie przechodzi na kandydacie; test, który
przechodzi na legacy i nie przechodzi na kandydacie --- jest.

\textbf{5.5 Uruchomienie agenta.} Otwórz \repopath{nop-p6} na gałęzi kandydata. Przed
czymkolwiek innym zapytaj agenta, jakich skilli używa --- tym samym pytaniem, którym
sprawdza to przewodnik standardów (\path{USING-WITH-MODERNIZE-DOTNET.md}, §6):

\begin{lstlisting}
Which custom skills from .github/skills are you applying?
\end{lstlisting}

W Copilot CLI: \repopath{/skills list}, potem
\repopath{/skills info applying-dotnet-migration-standards}. Odpowiedź wklej do zapisu
przebiegu; jeśli \repopath{applying-dotnet-migration-standards} brak, zatrzymaj się
i popraw krok 5.3 (to jest dokładnie kryterium R4: agent podejmuje skill). Punkt startu
zależy od klienta: Visual Studio --- prawy klik na solution, \textbf{Modernize}, albo
\repopath{@Modernize} w Copilot Chat; VS Code --- \repopath{@upgrade} w Copilot Chat albo
agent \textbf{Upgrade} na liście agentów; Copilot CLI --- \repopath{copilot} w katalogu
repozytorium, potem \repopath{/agent} i wybór \textbf{Upgrade}, a gdy skilla nie ma na
\repopath{/skills list} --- \repopath{/skills reload}
(podrozdział~\ref{sec:standardy-modernize}, krok 5).

Prompty wysyła się dosłownie, w tej kolejności. Dwa pierwsze należą do przewodnika
standardów (§5, kroki 3 i 4), który jest ich jedynym źródłem: obowiązuje brzmienie
z wersji standardów zainstalowanej w kroku 5.3, a runbook benchu cytuje je w wersji 0.1.0.
Drugi to instrukcja stała, którą agent zapisuje w
\path{.github/upgrades/<scenarioId>/scenario-instructions.md}:

\begin{lstlisting}
Upgrade this solution to .NET 10. Apply the applying-dotnet-migration-standards skill from
.github/skills to every task: preserve behaviour, and record every behaviour flag in
docs/migration/behaviour-flags.md instead of fixing it.

From now on, for all tasks in this upgrade, follow the applying-dotnet-migration-standards skill.
\end{lstlisting}

Trzeci prompt to zakres okazu, z runbooka P6 benchu:

\begin{lstlisting}
The solution is src/NopCommerce.sln, on .NET Framework 4.5.1. Migrate the ASP.NET MVC 5
web application (Nop.Web, the Administration area and the plugins) to ASP.NET Core on
.NET 10. Keep the SQL Server database schema, the public URLs and the observable
behaviour of the storefront unchanged. Keep the existing unit test projects and make them
build and run on .NET 10. Commit your work on the current branch in logical steps.
\end{lstlisting}

Podział jest celowy: prompty ogólne mają jedno źródło, więc runbook i przewodnik nie mogą
prosić o różne rzeczy, a runbook dokłada tylko to, co dotyczy konkretnego okazu. Przewodnik
wprowadza rejestr flag zachowania: podejrzany defekt agent oznacza, zamiast go naprawiać
(standard SK-MIG-011). Etapy agenta --- assessment (\repopath{assessment.md}), planning
(\repopath{plan.md}), execution (\repopath{tasks.md}) --- przegląda się przed każdym
„kontynuuj”; w assessment musi być bazowa linia legacy, o którą prosi skill (źródło
kultury, klucze konfiguracji, lazy loading EF6, endpointy JSON, testy niezaliczone).
Jeśli jej brak, wyślij dosłownie:
\begin{lstlisting}
Record the legacy baseline the applying-dotnet-migration-standards skill asks for before planning
\end{lstlisting}

\textbf{5.6 Interwencje człowieka.} Każdą zmianę kodu przez człowieka robi się osobnym
commitem z komunikatem zaczynającym się od \repopath{human:} i powodem; pakiet dowodowy
odróżnia zmiany agenta od zmian człowieka po tym prefiksie. Każdy prompt, każde pytanie
agenta i każdy wybór idą do zapisu dosłownie, w kolejności; commitów agenta się nie
squashuje.

\textbf{5.7 Build i testy kandydata.}
\begin{lstlisting}
dotnet --version
dotnet build src\NopCommerce.sln -c Release
dotnet test  src\NopCommerce.sln -c Release --logger "trx;LogFileName=p6-candidate-tests.trx"
\end{lstlisting}
P6 jest zrobione, gdy build się udaje, a każdy test zaliczony w bazowej linii przechodzi.
Test usunięty, pominięty albo opróżniony przez agenta nie liczy się jako zaliczony.

\textbf{5.8 Potwierdzenie, że agent użył standardów.} Przynajmniej jedno z poniższych
musi zachodzić, a pierwsze dwa powinny: agent wymienia skill
\repopath{applying-dotnet-migration-standards}; instrukcja stała jest w
\repopath{scenario-instructions.md}; rejestr flag istnieje z bazową linią; identyfikatory
reguł (\repopath{SK-MIG-}) pojawiają się w assessment, planie albo notatkach zadań; brak
cichego przełączenia globalizacji:

\begin{lstlisting}
git grep -n "SECONDKEY-FLAG"
git grep -n -E "InvariantGlobalization|UseNls|AppLocalIcu|EnsureCreated"
\end{lstlisting}

Do notatek P6 zapisz: narzędzie i jego wersję, wersję standardów z
\repopath{metadata.version} skilla, czy agent podjął skill sam, czy dopiero po jawnym
prompcie, oraz liczbę flag dla każdej reguły.

\textbf{5.9 Pull request i zapis przebiegu.}
\begin{lstlisting}
git push fork secondkey/p6-candidate
gh pr create --repo <your-account>/nopCommerce --base secondkey/p6-base --head secondkey/p6-candidate \
  --title "P6: modernize-dotnet candidate - nopCommerce 3.90 to .NET 10" \
  --body-file p6-record.md
git rev-parse HEAD                                  # commit kandydata
\end{lstlisting}
Pull requestu się nie scala: jego head commit jest kandydatem, a późniejsze commity tworzą
nowego kandydata. Zapis przebiegu (operator, czasy, godziny, wersje klienta, agenta i
SDK, commit legacy, wersja i commit standardów, lista skilli podana przez agenta, PR i
commit kandydata, liczby testów legacy i kandydata, testy przechodzące na legacy, a nie na
kandydacie, commity \repopath{human:}, prompty i odpowiedzi) trafia jako
\path{results/p6/run-record.md} do benchu, a godziny do \path{P10-ONBOARDING-HOURS.md}.
Szablon zapisu --- z godzinami pracy i oczekiwania na agenta osobno, modelem wybranym
w Copilot, promptami z odpowiedziami i uwagami (pliki usunięte przez agenta, zastąpione
pakiety, pominięte ostrzeżenia) --- podaje podrozdział~\ref{sec:bench-p6}.

\textbf{5.10 Konfiguracja bramki w kandydacie.} Pakiet NuGet \repopath{SecondKey.Standards}
ustawia severity każdej diagnostyki, która sprawdza standard. Nie jest na nuget.org:
pochodzi z artefaktu CI repozytorium standardów albo z \repopath{./scripts/verify-package.sh}.
Wersja pakietu musi być ta sama co wersja skilla --- agent i bramka muszą dostać te same
standardy:

\begin{lstlisting}
mkdir -p .packages && cp /path/to/SecondKey.Standards.0.1.0.nupkg .packages/
\end{lstlisting}
\begin{lstlisting}
<!-- nuget.config: <add key="secondkey-local" value=".packages" /> -->
<GlobalPackageReference Include="SecondKey.Standards" Version="0.1.0" />
\end{lstlisting}

Bez centralnego zarządzania pakietami referencja trafia do \path{Directory.Build.props};
skąd wziąć pakiet, jak wyłączyć go w projekcie i jak użyć konfiguracji bez pakietu, mówią
podrozdziały~\ref{sec:standardy-modernize} (krok 7) i \ref{sec:standardy-wersje}. Każdą
severity można nadpisać w \path{.editorconfig} repozytorium docelowego, który zawsze
wygrywa --- ale tylko w buildzie: \repopath{portcullis scan} (CLI, akcja, kontener)
raportuje każdą diagnostykę na jej poziomie domyślnym i nadpisań
\texttt{dotnet\_diagnostic.<id>.severity} nie stosuje
(podrozdział~\ref{sec:standardy-wagi-kanaly}).

Krok CI \repopath{konradcinkusz/letsgolegacy.secondkey-standards/actions/drift-check@main}
w repozytorium docelowym kończy się błędem, gdy istnieje nowsze wydanie albo gdy skill i
pakiet się nie zgadzają. Przed pierwszym wydaniem standardów zgłasza „no release” i się nie
udaje, więc dla P6 przed wydaniem się go pomija.

\begin{xltabular}{\linewidth}{@{}>{\raggedright\arraybackslash}p{5.6cm} >{\raggedright\arraybackslash}X@{}}
\toprule
Objaw & Przyczyna i działanie \\
\midrule
\endfirsthead
\toprule
Objaw & Przyczyna i działanie \\
\midrule
\endhead
Agent nie wspomina standardów & \path{.github/skills/} nie ma w otwartym folderze albo klient nie ładuje skilli repozytorium; sprawdź ścieżkę, poproś o skill z nazwy, zrestartuj klienta, w Copilot CLI \repopath{/skills reload}; w razie potrzeby skopiuj skill także do \path{.github/upgrades/skills/} \\
Diff PR zawiera tysiące niezwiązanych linii & gałęzie nie powstały z przypiętego commita albo konwersja końców linii (\repopath{core.autocrlf}) \\
„Testy przechodzą”, ale wykonało się ich mniej & agent wykluczył, pominął albo usunął projekty testowe; porównaj liczby z bazową linią \\
P7 nie odtwarza buildu & kandydat zależy od czegoś tylko na maszynie operatora (lokalne źródło NuGet, SDK preview); zapisz SDK i każde źródło \\
Kandydat zmienił się w trakcie P7 & ktoś wypchnął zmiany na gałąź kandydata; P7 przypina commit, nigdy gałąź \\
\repopath{git push}: „shallow update not allowed” & fork utworzony z kopią tylko domyślnej gałęzi; pobierz tag bez \repopath{--depth} albo utwórz fork domyślnie przez \repopath{gh repo fork} \\
Agent włącza \repopath{InvariantGlobalization} albo zmienia oczekiwaną wartość w teście & poszedł za ogólną poprawką zamiast reguły; wskaż mu SK-MIG-007 (albo SK-MIG-013), niech cofnie zmianę i zapisze flagę zachowania; dopisz też regułę do instrukcji stałej \\
Build kończy się błędem \repopath{PORTCULLIS\_} albo \repopath{CA} & zainstalowana konfiguracja bramki, a kod łamie regułę o severity \repopath{error}; popraw zgodnie z regułą albo wycisz linię z uzasadnieniem \repopath{SECONDKEY-FLAG} \\
\bottomrule
\end{xltabular}

Pozostałe objawy po stronie klienta i skilla: tabela rozwiązywania problemów
w podrozdziale~\ref{sec:standardy-modernize}.

\textbf{Checklista (runbook P6)}
\begin{checklista}
  \item Fork; \repopath{secondkey/p6-base} to przypięty commit plus commit ze skillami i nic więcej.
  \item Skille z otagowanego wydania standardów; tag i commit w zapisie.
  \item Bazowa linia testów legacy (liczby per assembly).
  \item Agent wymienił skille Second Key przed jakąkolwiek zmianą.
  \item Każdy prompt, pytanie i wybór zapisane dosłownie; zmiany człowieka jako commity \repopath{human:}.
  \item \repopath{dotnet build} się udaje; każdy test zaliczony na legacy przechodzi na kandydacie.
  \item Pull request do \repopath{secondkey/p6-base}, niescalony; commit kandydata zapisany.
  \item Zapis przebiegu w benchu jako \path{results/p6/run-record.md}; godziny w logu P10.
\end{checklista}

\zrodla{\path{letsgolegacy.secondkey-nopcommerce-bench/docs/P6-RUNBOOK.md},
\path{letsgolegacy.secondkey-standards/docs/USING-WITH-MODERNIZE-DOTNET.md},
\path{letsgolegacy.secondkey-standards/README.md},
\path{letsgolegacy.secondkey-standards/docs/WORKPLAN.md}}

% ---------------------------------------------------------------------------
\subsection{Krok 6: bramka Portcullis na pull request kandydata}\label{sec:procedura-krok6}

\textbf{Cel.} Sprawdzić deterministycznymi analyzerami Roslyn linie, które pull request
kandydata zmienia, i zapisać wynik jako SARIF do pakietu dowodowego. W benchu to ticket P8
(plan: reguły migracyjne Portcullis na PR kandydata, SARIF dołączony do pakietu).
Referencja: rozdział~\ref{sec:portcullis}.

\begin{figure}[htbp]
  \centering
  \includegraphics[width=0.8\linewidth,height=0.55\textheight,keepaspectratio]{\dgm{a8-bramka-portcullis}}
  \caption{Bramka Portcullis: znalezione a blokujące.}
  \label{fig:bramka}
\end{figure}

\textbf{Postępowanie.}
\begin{enumerate}
  \item Checkout z pełną historią (\repopath{fetch-depth: 0} w GitHub Actions,
        \repopath{GIT\_DEPTH: 0} w GitLab). Bez commita bazowego PR Portcullis nie ustali
        zmienionych linii i przejdzie na zakres całego drzewa --- ostrzega o tym, ale
        poprawką jest pełny checkout.
  \item Skan z zakresem commitów i zapisem SARIF (przed wydaniem pakietów --- z kodu
        źródłowego):
\begin{lstlisting}
BASE=$(git merge-base secondkey/p6-base HEAD)    # baza pull requestu kandydata (krok 5.2)
dotnet run --project /path/to/letsgolegacy.portcullis/src/Portcullis.Cli -- \
  scan ./src --provenance-range "$BASE..HEAD" --provenance-repo . --sarif portcullis.sarif
\end{lstlisting}
        W forku kandydata z kroku 5.2 nie ma zdalnego \repopath{origin}, a pull request idzie
        do \repopath{secondkey/p6-base}. Ogólna postać z tutorialu Portcullis,
        \texttt{git merge-base origin/main HEAD}, dotyczy repozytorium, którego pull
        requesty celują w \repopath{main}.
  \item Odczytaj wynik: \repopath{summary} to wszystko, co znaleziono; \repopath{gate} to
        to, co blokuje. Blokuje tylko naruszenie o poziomie błędu w liniach zmienionych
        przez PR, nieprzyjęte przez baseline. Różnica między tymi liczbami to istniejący
        dług. Kody wyjścia: 0 --- nie blokuje, 1 --- blokuje (werdykt, nie awaria),
        2 --- błędne użycie.
  \item Na kodzie z istniejącymi naruszeniami: zakres diff (domyślny w akcji) i plik
        baseline (\repopath{--baseline portcullis-baseline.json}, raz z
        \repopath{--write-baseline}); znalezisko przyjęte przez baseline jest raportowane,
        ale nigdy nie blokuje. Stopniowe zaostrzanie severity w \repopath{.editorconfig}
        działa tylko w buildzie z pakietem analyzerów: \repopath{portcullis scan} (CLI,
        akcja, kontener) nadpisań \texttt{dotnet\_diagnostic.<id>.severity} nie stosuje
        (podrozdział~\ref{sec:standardy-wagi-kanaly}).
  \item Dołącz SARIF do łańcucha: \repopath{sk gate --sarif portcullis.sarif} podsumowuje
        logi per reguła i kończy się kodem 1, gdy znalezisko blokuje (błąd, którego nie
        obejmuje zaakceptowana suppression ani stan baseline \repopath{unchanged}/
        \repopath{absent}); ostrzeżenia i uwagi nigdy nie blokują. Plik, który nie jest
        SARIF 2.1.0, to błędne wejście (kod 3), nigdy „brak znalezisk”.
\end{enumerate}

Na GitHub bramką ma być akcja \repopath{konradcinkusz/letsgolegacy.portcullis@v0}. Instaluje
\repopath{Portcullis.Cli} i \repopath{Portcullis.CiComment} z NuGet w wersji z wejścia
\repopath{version} (domyślnie \texttt{0.1.0}), a job potrzebuje \texttt{pull-requests: write}
dla komentarza i \repopath{fetch-depth: 0}; wejścia i wyjścia opisuje
podrozdział~\ref{sec:portcullis-dystrybucja}. SARIF do code scanning publikuje się na razie
osobnym krokiem CLI i \repopath{github/codeql-action/upload-sarif}. Gotowy workflow SARIF
opisuje sekcja~\ref{sec:portcullis-tutorial-sarif}, a sprawdzenie konwencji przed włączeniem
bramki --- sekcja~\ref{sec:portcullis-tutorial-konwencje}. Pakiety i obraz kontenera są
zbudowane i sprawdzone lokalnie, ale jeszcze nieopublikowane, a tagu \texttt{v0} jeszcze nie
ma: do pierwszego tagu \repopath{v*} odwołanie \repopath{@v0} się nie rozwiąże i działa
tylko ścieżka z kodu źródłowego (polecenie wyżej).

\begin{xltabular}{\linewidth}{@{}>{\raggedright\arraybackslash}p{5.6cm} >{\raggedright\arraybackslash}X@{}}
\toprule
Objaw & Przyczyna i działanie \\
\midrule
\endfirsthead
\toprule
Objaw & Przyczyna i działanie \\
\midrule
\endhead
\repopath{gate.scope} to \repopath{"all"} zamiast \repopath{"diff"} & brak commita bazowego w checkoucie; pełna historia. Przy \repopath{"all"} z \repopath{degradedReason} git nie rozwiązał zakresu i skan przeszedł na całe drzewo --- powód mówi, co zawiodło \\
Bramka przechodzi, choć powinna blokować & przy \repopath{"diff"} naruszenie leży poza liniami zmienionymi przez PR --- bramka działa zgodnie z projektem \\
Pakiet się instaluje, ale diagnostyk brak & DLL musi być w \path{analyzers/dotnet/cs/} w \repopath{.nupkg}; analyzer spakowany do \path{lib/} nic nie robi \\
\repopath{warning CS9057} & SDK starszy niż .NET 8; \repopath{Portcullis.Analyzers} celuje w Roslyn 4.8 właśnie po to, żeby tego uniknąć \\
Reguły nie odpalają na kodzie, o którym wiesz, że jest zły & Portcullis parsuje źródła bezpośrednio, nie przez MSBuild, więc typ z pakietu NuGet albo innego projektu się nie rozwiązuje \\
Reguła strzela na poprawnym kodzie & to błąd reguły do zgłoszenia (issue z fragmentem kodu), nie coś do obchodzenia; wyciszenie w \path{.editorconfig} działa tylko w buildzie z pakietem analyzerów, a w bramce CLI finding przyjmuje baseline (podrozdział~\ref{sec:portcullis-baseline}) \\
\bottomrule
\end{xltabular}

Pozostałe objawy i komunikaty: tabela rozwiązywania problemów
w podrozdziale~\ref{sec:portcullis-tutorial}.

\textbf{Checklista}
\begin{checklista}
  \item Pełny checkout; \repopath{gate.scope} równe \repopath{"diff"}.
  \item SARIF zapisany dla head commita kandydata.
  \item Wersja pakietu standardów w kandydacie równa wersji skilla.
  \item \repopath{sk gate} uruchomione na SARIF; wynik trafia do pakietu.
\end{checklista}

\zrodla{\path{letsgolegacy.portcullis/docs/TUTORIAL.md},
\path{letsgolegacy.portcullis/docs/WORKPLAN.md},
\path{letsgolegacy.portcullis/docs/CONFIGURATION.md},
\path{letsgolegacy.portcullis/action.yml},
\path{letsgolegacy.portcullis/SECURITY.md},
\path{letsgolegacy.secondkey-nopcommerce-bench/docs/P6-RUNBOOK.md},
\path{letsgolegacy.secondkey/docs/formats/evidence.md},
\path{letsgolegacy.secondkey-nopcommerce-bench/docs/WORKPLAN.md},
\path{letsgolegacy.secondkey-nopcommerce-bench/scripts/eshop/5-scan-eshop.ps1}}

% ---------------------------------------------------------------------------
\subsection{Krok 7: replay legacy kontra kandydat i werdykt}\label{sec:procedura-krok7}

\textbf{Cel.} Odtworzyć nagrany ruch na legacy i na kandydacie w identycznych warunkach
i porównać obie strony względem kontraktu. W benchu to ticket P7; kryterium odbioru:
\repopath{verdict.json} istnieje, a różnica NLS → ICU jest znaleziona albo jej brak jest
udokumentowany.

\begin{figure}[htbp]
  \centering
  \includegraphics[width=\linewidth]{\dgm{a4-sekwencja-replay}}
  \caption{Replay: reset i odtworzenie scenariusza po każdej stronie.}
  \label{fig:replay}
\end{figure}

\textbf{Postępowanie.}
\begin{enumerate}
  \item \textbf{Przypnij kandydata} --- commit, nigdy gałąź. Zbuduj go i uruchom obok
        systemu legacy. Plan benchu: host Windows, \repopath{http://localhost:8090/},
        własna kopia bazy odtworzona z tego samego snapshotu.
  \item \textbf{Skonfiguruj reset po obu stronach} w \repopath{secondkey.yaml}
        (\repopath{http}, \repopath{sqlServerSnapshot} albo \repopath{command}; łańcuch
        połączenia wyłącznie przez zmienną środowiskową \repopath{connectionStringEnv}).
        Reset \repopath{sqlServerSnapshot} tworzy snapshot przy pierwszym resecie i
        zostawia go po przebiegu; zamyka wszystkie połączenia z bazą, także połączenia
        aplikacji z puli, i nie czeka, aż aplikacja się podniesie; resetuje tylko bazę ---
        cache, sesje i timery aplikacji przetrwają. Login z łańcucha połączenia musi móc
        tworzyć bazę i przywracać tę bazę (w sandboxie \repopath{dbcreator}); istniejący
        snapshot o tej nazwie jest używany bez zmian, więc nowy stan startowy wymaga jego
        usunięcia (sekcja~\ref{sec:cli-reset}). Zmienną z łańcuchem połączenia trzeba
        wyeksportować do środowiska procesu \repopath{sk}: ani \repopath{sk}, ani skrypty
        rdzenia nie wczytują pliku \path{.env} (sekcja~\ref{sec:rdzen-setup}). Jej brak
        zgłasza dopiero \repopath{sk replay}, kodem 3. Wartości, które generuje serwer,
        a kolejne żądania muszą odesłać (typowo token anti-forgery), przenosi reguła
        \repopath{replay.correlation}; zmiany wierszy, na które patrzą selektory
        \repopath{db} klauzul, zapisuje sonda \repopath{probe: sqlServer}
        (sekcja~\ref{sec:cli-sonda}). Na Windows ścieżki w \repopath{secondkey.yaml} pisze
        się z ukośnikami albo w apostrofach (sekcja~\ref{sec:cli-konfiguracja}). Bench
        ustawia \repopath{timeoutSeconds: 120}, bo pierwsze żądanie scenariusza następuje po
        przywróceniu bazy.
  \item \textbf{Uruchom replay}:
\begin{lstlisting}
sk replay --capture .secondkey/traffic.skcap --legacy http://127.0.0.1:5080 --candidate http://127.0.0.1:5081
\end{lstlisting}
        W benchu: \repopath{.\textbackslash{}scripts\textbackslash{}8-replay.ps1
        -Candidate http://localhost:8090/}. Każdy scenariusz dostaje po każdej stronie
        świeży cookie jar (nagrane ciasteczka nigdy nie są wysyłane), przekierowania nie
        są śledzone (302 to odpowiedź do porównania), a korelacja przenosi wartości
        generowane przez serwer (np. tokeny anti-forgery) z ostatniej odpowiedzi tej samej
        strony. Kod 4 oznacza nieudany reset: scenariusz nie został wysłany do tej strony
        (\repopath{reset-failed}) i przebieg nie jest uczciwym porównaniem. Kod 0 nie
        znaczy, że każda strona odpowiedziała: liczbę wyników bez odpowiedzi podaje
        podsumowanie (\texttt{<n> without an answer}), a każdy taki wynik jest w werdykcie
        regresją \repopath{missing-result}. Korelacja czyta tylko odpowiedzi zapisane jako
        tekst (HTML, XML), nie odpowiedzi JSON (sekcja~\ref{sec:cli-sonda}).
        \repopath{sk replay} bez pytania nadpisuje przebieg o tej samej ścieżce
        \repopath{--out} (sekcja~\ref{sec:cli-wywolanie}).
  \item \textbf{Porównaj}:
\begin{lstlisting}
sk compare --contract contract.yaml --run .secondkey/run.skrun --out .secondkey/verdict.json
\end{lstlisting}
        Kod 0 --- \repopath{pass}, 1 --- \repopath{fail} (jakakolwiek regresja),
        2 --- \repopath{review} (bez regresji, jest fix candidate do decyzji), 3 --- błędne
        wejście (np. przebieg bez wyników: nic nie porównano, więc nic nie może przejść).
\end{enumerate}

W benchu skrypt 8 odrzuca każdy przebieg z wynikiem bez odpowiedzi albo z nieudanym resetem,
także legacy kontra kandydat, i zawsze czyta \path{contract/secondkey.yaml}. Reset strony
\repopath{candidate} w tym pliku przywraca bazę \repopath{nopcommerce\_legacy} ze snapshotu
\repopath{nopcommerce\_legacy\_snapshot}: to pasuje do A/A i do mutantów, które działają na
bazie legacy, a kandydat P7 ma mieć własną kopię bazy (podrozdział~\ref{sec:bench-topologia}).
Porównanie i pakiet robi w benchu skrypt 9 (\texttt{pwsh scripts/9-verdict.ps1}); wymóg
\repopath{pass} dotyczy tylko przebiegu A/A.

\begin{xltabular}{\linewidth}{@{}>{\raggedright\arraybackslash}p{5.6cm} >{\raggedright\arraybackslash}X@{}}
\toprule
Objaw & Przyczyna i działanie \\
\midrule
\endfirsthead
\toprule
Objaw & Przyczyna i działanie \\
\midrule
\endhead
Skrypt 8: \texttt{The replay has results without an answer or failed resets; it is not a fair comparison.} & wynik bez odpowiedzi albo nieudany reset po którejkolwiek stronie, także u kandydata; najpierw środowisko tej strony \\
Skrypt 8: \texttt{state\textbackslash{}traffic.json names a partial recording (-Only): record the whole set first.} & nagranie z \repopath{-Only} nie jest zestawem ruchu; nagraj cały zestaw (krok~2) \\
Skrypt 8: \texttt{The run holds \ldots{} scenarios; the recording has \ldots{}} & przebieg ma inną liczbę scenariuszy niż nagranie, więc nie porównuje całego zestawu \\
\repopath{sk replay} kończy się kodem 3 i nazwą zmiennej środowiskowej & zmiennej z \repopath{connectionStringEnv} nie ma w środowisku procesu \repopath{sk}; wyeksportuj ją \\
\repopath{sk replay} kończy się kodem 4 & reset strony się nie udał; jej scenariusze nie zostały do niej wysłane i są zapisane jako \repopath{reset-failed} \\
SQL Server odmawia usunięcia bazy albo przywrócenia jej z kopii zapasowej & istnieje snapshot tej bazy; usuń go (\texttt{DROP DATABASE [Shop\_sk]}) \\
Pierwsza odpowiedź po przywróceniu bazy to błąd & restore zamyka połączenia z bazą, także z puli aplikacji, a \repopath{sk} nie czeka na aplikację i nie ponawia żądania; sprawdź raz pierwszą odpowiedź po restore \\
\repopath{sk compare} odrzuca przebieg & przebieg bez \repopath{run.end} został przerwany i nigdy nie jest porównywany; przebieg bez wyników to błędne wejście (kod 3) \\
\bottomrule
\end{xltabular}

Pozostałe komunikaty replay i resetu: sekcje~\ref{sec:cli-replay} i \ref{sec:cli-reset}.

\textbf{Checklista}
\begin{checklista}
  \item Kandydat przypięty commitem; ta sama wersja nagrania co w A/A.
  \item Obie strony startują każdy scenariusz z tego samego snapshotu; zero nieudanych resetów.
  \item Zero wyników bez odpowiedzi (skrypt 8 benchu odrzuca przebieg z takim wynikiem); poza
        benchem każdy taki wynik to \repopath{regression} \repopath{missing-result} do
        wyjaśnienia.
  \item \repopath{verdict.json} zapisany i zwalidowany (\repopath{sk validate verdict.json}).
\end{checklista}

\zrodla{\path{letsgolegacy.secondkey/docs/cli.md},
\path{letsgolegacy.secondkey/docs/formats/README.md},
\path{letsgolegacy.secondkey/src/SecondKey.Replay/ReplayEngine.cs},
\path{letsgolegacy.secondkey/src/SecondKey.Artifacts/Runs/RunWriter.cs},
\path{letsgolegacy.secondkey-nopcommerce-bench/docs/P6-RUNBOOK.md},
\path{letsgolegacy.secondkey-nopcommerce-bench/docs/TOPOLOGY.md},
\path{letsgolegacy.secondkey-nopcommerce-bench/docs/WORKPLAN.md},
\path{letsgolegacy.secondkey-nopcommerce-bench/scripts/README.md},
\path{letsgolegacy.secondkey-nopcommerce-bench/scripts/8-replay.ps1},
\path{letsgolegacy.secondkey-nopcommerce-bench/scripts/9-verdict.ps1},
\path{letsgolegacy.secondkey-nopcommerce-bench/contract/secondkey.yaml}}

% ---------------------------------------------------------------------------
\subsection{Krok 8: analiza różnic i decyzje}\label{sec:procedura-krok8}

\textbf{Cel.} Zrozumieć każdą wymianę, która nie jest równa, i podjąć właściwą decyzję:
poprawić kandydata, zmienić kontrakt w przeglądzie albo przekazać fix candidate osobie,
która decyduje.

\begin{figure}[htbp]
  \centering
  \includegraphics[width=0.9\linewidth,height=0.62\textheight,keepaspectratio]{\dgm{a5-klasy-werdyktu}}
  \caption{Klasa wymiany: pierwszy pasujący krok decyduje.}
  \label{fig:klasy}
\end{figure}

\textbf{Jak czytać werdykt.} Każda wymiana jest w dokładnie jednej klasie; decyduje
pierwszy pasujący krok (rysunek~\ref{fig:klasy}). Wynik przebiegu to \repopath{fail}, gdy
jakakolwiek wymiana jest regresją, \repopath{review}, gdy żadna nie jest, ale jest fix
candidate, i \repopath{pass} w pozostałych przypadkach. Pole \repopath{diffs} wymienia
każdą surową różnicę ze ścieżką w języku ścieżek kontraktu. Ścieżkę pod korzeniem
\repopath{response}, \repopath{extract}, \repopath{db} albo \repopath{outbound} można wkleić
do reguły \repopath{ignore} albo klauzuli bez zmian; różnicy w \repopath{error} albo
\repopath{extractErrors} nie wyjaśni żadna reguła (sekcja~\ref{sec:formaty-werdykt}).
Różnica wyjaśniona przez normalizację nazywa regułę w \repopath{normalizedBy}.
\repopath{summary.clauses} liczy klauzule wszystkie, decydujące i niesprawdzone.

\textbf{Decyzje.}
\begin{xltabular}{\linewidth}{@{}>{\raggedright\arraybackslash}p{4.8cm} >{\raggedright\arraybackslash}X@{}}
\toprule
Klasa i powód & Co zrobić \\
\midrule
\endfirsthead
\toprule
Klasa i powód & Co zrobić \\
\midrule
\endhead
\repopath{regression}, \repopath{missing-result} & strona bez odpowiedzi (timeout, połączenie, reset); najpierw środowisko: przebieg z nieudanym resetem nie jest uczciwym porównaniem \\
\repopath{regression}, \repopath{clause-regression: <id>} & kandydat łamie zaakceptowaną klauzulę, którą legacy spełnia, albo klauzuli nie da się na nim ocenić (np. ekstraktor nie odczytał wartości); poprawka po stronie kandydata \\
\repopath{regression}, \repopath{uncovered-difference} & różnica, której kontrakt nie tłumaczy; jeśli jest zachowaniem --- poprawka kandydata; jeśli nie ma znaczenia --- przeglądnięta zmiana kontraktu (reguła normalizacji albo klauzula) \\
\repopath{fix-candidate}, \repopath{clause-fixed: <id>} & kandydat spełnia klauzulę, którą legacy łamie; decyduje człowiek (wynik \repopath{review}); standard SK-MIG-011 każe agentowi flagować podejrzany defekt, a nie naprawiać go po cichu \\
klauzula \repopath{violated-both} & złamana po obu stronach, sama o niczym nie decyduje; według ADR 0003 to defekt sprzed migracji, który migracja zachowała. W dowodzie A/A, gdzie obie strony to legacy, oznacza kontrakt, który myli się co do legacy --- wtedy poprawka kontraktu i ponowny dowód A/A \\
\bottomrule
\end{xltabular}

Sondy NLS → ICU z planu ruchu benchu (formatowanie ujemnych kwot walutowych, zamiana
wielkości liter i przycinanie kodów kuponów, porównanie kodów pocztowych, sortowanie,
daty) są hipotezami do sprawdzenia właśnie w tym kroku; najmocniejszym kandydatem według
źródła są ujemne kwoty w walucie \repopath{en-US}: NLS formatuje je jako
\repopath{(\$100.00)}, dane ICU jako \repopath{-\$100.00}.

Plan ruchu benchu wymienia też objaw po stronie kandydata: kupon przechodzi na legacy, a u
kandydata jest „not found”, gdy wyszukiwanie rabatu przeszło z SQL (collation bez
rozróżniania wielkości liter) do pamięci (porównanie ordinal). Scenariusz w werdykcie to
sesja z nagrania; jak bench zamienia identyfikator sesji na identyfikator scenariusza planu
ruchu, opisuje podrozdział~\ref{sec:bench-rozgrzewka}.

\textbf{Pętla zmiany.} Każda poprawka kandydata to nowy commit, więc nowy kandydat: krok 6
(bramka) i krok 7 (replay i werdykt) wykonuje się ponownie na nowym head commicie, a zapis
przebiegu to odnotowuje. Kontrakt zmienia się wyłącznie przeglądniętą zmianą; po każdej
zmianie kontraktu wraca się do dowodu A/A.

\textbf{Checklista}
\begin{checklista}
  \item Każda regresja przypisana: środowisko, kandydat albo kontrakt.
  \item Każdy fix candidate przekazany do decyzji człowieka, z powodem.
  \item Po zmianie kontraktu ponowny dowód A/A (\repopath{pass}).
  \item Po zmianie kandydata ponownie krok 6 i 7 na nowym commicie.
\end{checklista}

\zrodla{\path{letsgolegacy.secondkey/docs/formats/verdict.md},
\path{letsgolegacy.secondkey/docs/adr/0003-verdict-classes.md},
\path{letsgolegacy.secondkey-nopcommerce-bench/docs/P4-TRAFFIC-PLAN.md},
\path{letsgolegacy.secondkey-nopcommerce-bench/docs/P6-RUNBOOK.md},
\path{letsgolegacy.secondkey-nopcommerce-bench/docs/adr/0004-chain-tools-and-eshop-warmup.md},
\path{letsgolegacy.secondkey-nopcommerce-bench/results/p5/README.md},
\path{letsgolegacy.secondkey-standards/README.md}}

% ---------------------------------------------------------------------------
\subsection{Krok 9: pakiet dowodowy}\label{sec:procedura-krok9}

\textbf{Cel.} Złożyć katalog, który pozwala osobie zatwierdzającej zdecydować o
zmigrowanym systemie bez uruchamiania Second Key: wyjaśnia werdykt, niesie jego wejścia i
podaje SHA-256 każdego pliku.

\begin{lstlisting}
sk gate --sarif portcullis.sarif
sk evidence --verdict verdict.json --contract contract.yaml --run run.skrun \
            --sarif portcullis.sarif --out evidence --pdf auto
\end{lstlisting}

\begin{itemize}
  \item \repopath{sk evidence} kończy się kodem 0, gdy pakiet jest zapisany, niezależnie od
        werdyktu: pipeline powinien opublikować także pakiet, który wyjaśnia porażkę. Kod
        własny werdyktu zwraca \repopath{sk compare}.
  \item Odmawia (kod 3) kontraktu albo przebiegu, który nie jest tym, który nazywa werdykt.
  \item \repopath{--run} dołącza przebieg, żeby każdy mógł przeliczyć werdykt; gdy
        nagrany ruch musi zostać w swojej granicy, pomija się go.
  \item \repopath{--pdf auto} drukuje \repopath{report.pdf}, gdy na maszynie jest
        przeglądarka oparta na Chromium (\repopath{SK\_CHROME\_PATH}, potem PATH, potem
        miejsca instalacji Chrome, Edge i przeglądarek Playwright;
        sekcja~\ref{sec:cli-evidence}); \repopath{required} czyni brak PDF błędem (kod 4),
        \repopath{off} pomija PDF.
  \item Z katalogu z poprzednim pakietem znikają najpierw pliki z jego manifestu; jeśli
        zostaje w nim coś innego, polecenie odmawia (kod 64) --- obce pliki zostają, ale pliki
        poprzedniego pakietu są już usunięte. Katalog niepusty bez manifestu pakietu jest
        odrzucany bez usuwania czegokolwiek. Przy \repopath{--pdf required} porażka PDF
        (kod 4) zostawia katalog bez \repopath{manifest.json}, więc kolejne
        \repopath{sk evidence} do niego kończy się kodem 64, dopóki się go nie opróżni
        (sekcja~\ref{sec:cli-evidence}). Bench przed złożeniem usuwa katalog pakietu i po
        złożeniu weryfikuje skrót każdego pliku z manifestu.
  \item W benchu \path{scripts/9-verdict.ps1} wywołuje \repopath{sk evidence} bez
        \repopath{--sarif} i nie uruchamia \repopath{sk gate} (P8 jest zaplanowany), więc
        pakiet z SARIF bramki składa się poleceniami z tego kroku. Rozgrzewka robi oba
        kroki w \path{scripts/eshop/6-verdict-eshop.ps1}.
\end{itemize}

\textbf{Sprawdzenie pakietu} przez kogoś, kto go otrzymał:

\begin{lstlisting}
cd evidence
python3 -c "import json,hashlib; [print('ok' if hashlib.sha256(open(f['path'],'rb').read()).hexdigest()==f['sha256'] else 'CHANGED', f['path']) for f in json.load(open('manifest.json'))['files']]"
sk validate verdict.json contract.yaml
sk compare --contract contract.yaml --run run.skrun --out recomputed.json   # ten sam werdykt poza createdAt (oraz tool i inputs.*.path, gdy inny build albo inne nazwy)
\end{lstlisting}

\finding{Pakiet w fazie 01 jest niepodpisany}: statement in-toto ma \repopath{"signed":
false}. Dowodzi, które pliki należą do siebie i co mówią, a nic nie mówi o tym, kto go
złożył. Podpis (koperta DSSE, klucz właściciela) dochodzi w fazie 03.

\begin{xltabular}{\linewidth}{@{}>{\raggedright\arraybackslash}p{5.6cm} >{\raggedright\arraybackslash}X@{}}
\toprule
Objaw & Przyczyna i działanie \\
\midrule
\endfirsthead
\toprule
Objaw & Przyczyna i działanie \\
\midrule
\endhead
\repopath{sk evidence} kończy się kodem 3 & werdykt, kontrakt albo log SARIF jest niepoprawny albo kontrakt lub przebieg nie jest tym, który nazywa werdykt; komunikat podaje oba skróty SHA-256 \\
\repopath{sk evidence} kończy się kodem 64 & katalog \repopath{--out} zawiera coś innego niż poprzedni pakiet; wskaż pusty albo nowy katalog \\
Kod 4 przy \repopath{--pdf required} & brak przeglądarki opartej na Chromium, \repopath{SK\_CHROME\_PATH} wskazuje brakujący plik albo przeglądarka nie wydrukowała PDF w ciągu minuty; powód stoi po \texttt{error:} \\
\texttt{CHANGED} przy sprawdzaniu pakietu & skrót pliku różni się od zapisanego w \path{manifest.json} \\
\bottomrule
\end{xltabular}

\textbf{Checklista}
\begin{checklista}
  \item Pakiet złożony z werdyktu, kontraktu i --- jeśli wolno --- przebiegu z tego samego P7.
  \item SARIF bramki dołączony; \repopath{sk gate} wykonane.
  \item Każdy skrót z \repopath{manifest.json} zgodny; werdykt przeliczony ponownie daje to samo.
\end{checklista}

\zrodla{\path{letsgolegacy.secondkey/docs/formats/evidence.md},
\path{letsgolegacy.secondkey/src/SecondKey.Evidence/EvidencePack.cs},
\path{letsgolegacy.secondkey/src/SecondKey.Evidence/PdfRenderer.cs},
\path{letsgolegacy.secondkey/schemas/verdict.schema.json},
\path{letsgolegacy.secondkey-nopcommerce-bench/scripts/9-verdict.ps1},
\path{letsgolegacy.secondkey-nopcommerce-bench/scripts/eshop/6-verdict-eshop.ps1},
\path{letsgolegacy.secondkey-nopcommerce-bench/docs/WORKPLAN.md}}

% ---------------------------------------------------------------------------
\subsection{Krok 10: publikacja i liczby}\label{sec:procedura-krok10}

Bramką fazy 01 jest ticket P9 benchu: opublikować pakiet dowodowy, kontrakt i liczby ---
scenariusze, różnice, regresje, zabite mutanty. P9 jest w toku. Zabite mutanty są gotowe
(\#8): 11 z 11, przy przechodzącym M00, który kalibruje metodę
(podrozdział~\ref{sec:bench-p9}); jak sprawdzić jednego mutanta ręcznie, mówi krok 4a
(sekcja~\ref{sec:procedura-mutanty}). Właściwy komponent mutacji C9 (\repopath{sk mutate})
należy do fazy 02 --- w fazie 01 kończy się kodem 70; zastępuje go prerejestrowany zestaw
łatek benchu, porównywany tymi samymi \repopath{sk replay} i \repopath{sk compare} co
kandydat. Liczby scenariuszy i klauzul istnieją dla legacy (P4: 38 scenariuszy, 290
wymian; P5: 76 klauzul, w tym 15 \repopath{never}); pakiet dowodowy kandydata, różnice
i regresje przyniosą P6--P8. Pakiet cytuje pull request kandydata, jego head commit,
wersję standardów i zapis przebiegu P6.

\textbf{Jak bench zapisuje wynik.} Każdy zrobiony ticket ma
\path{results/<ticket>/README.md} z linkiem do przebiegu CI, commitem, narzędziami
z przypiętymi commitami i liczbami. Artefakty CI --- nagranie, przebieg, pakiet dowodowy ---
CI przechowuje 30 dni, a każdy przebieg workflowu tworzy nowe
(podrozdział~\ref{sec:bench-wyniki}); \path{results/p9/mutants.json} jest plikiem przebiegu
bajt w bajt, a jego SHA-256 job podsumowania drukuje w logu razem z plikiem. Godziny trafiają
do \path{docs/P10-ONBOARDING-HOURS.md} według jego zasad: jeden wiersz na sesję
i wykonawcę, osoba albo agent, nigdy oba; godzin agenta i człowieka się nie sumuje, każdy
wiersz ma dowód, a szacunki są oznaczone (podrozdział~\ref{sec:bench-plany}).

\textbf{Checklista}
\begin{checklista}
  \item Opublikowane: pakiet dowodowy kandydata, kontrakt i liczby --- scenariusze, różnice,
        regresje, zabite mutanty.
  \item Pakiet cytuje pull request kandydata, jego head commit, wersję standardów i zapis
        przebiegu P6.
  \item Liczba zabitych mutantów pochodzi z przebiegu, w którym M00 dał \repopath{pass}.
  \item Każda liczba w \path{results/} ma link do przebiegu CI i commit.
  \item Godziny agenta i człowieka w dzienniku P10 osobno; szacunki oznaczone.
\end{checklista}

\zrodla{\path{letsgolegacy.secondkey-nopcommerce-bench/docs/WORKPLAN.md},
\path{letsgolegacy.secondkey-nopcommerce-bench/docs/P9-MUTANTS.md},
\path{letsgolegacy.secondkey-nopcommerce-bench/results/p3/README.md},
\path{letsgolegacy.secondkey-nopcommerce-bench/results/p5/README.md},
\path{letsgolegacy.secondkey-nopcommerce-bench/results/p9/README.md},
\path{letsgolegacy.secondkey-nopcommerce-bench/docs/P6-RUNBOOK.md},
\path{letsgolegacy.secondkey-nopcommerce-bench/docs/P10-ONBOARDING-HOURS.md},
\path{letsgolegacy.secondkey/docs/WORKPLAN.md},
\path{letsgolegacy.secondkey/docs/cli.md}}

% ---------------------------------------------------------------------------
\subsection{Praca nad samym frameworkiem}\label{sec:procedura-rozwoj}

Repozytoria frameworku rozwija się według tych samych reguł:

\begin{enumerate}
  \item \textbf{Zacznij od planu pracy.} \path{docs/WORKPLAN.md} każdego repozytorium
        wymienia tickety, ich kryteria „done when” i status. Jeden ticket to jeden pull
        request, a kryterium odbioru nie jest parafrazowane w PR.
  \item \textbf{W rdzeniu przeczytaj analizę.} \path{docs/analysis/phase-01.md} mówi, które
        pliki należą do którego ticketu i jakie decyzje zapadły; jeśli zmiana się z nią
        nie zgadza, najpierw zmienia się analizę. \path{docs/reference-notes.md} mówi, co
        przeniesiono ze starszych repozytoriów. Czwartym źródłem są standardy estate,
        zadeklarowane w \path{.claude/settings.json}:
        \path{architecture-standards/docs/architecture/00-REFERENCE-ARCHITECTURE.md}
        i przewodniki. Master prompt (\path{docs/analysis/S1-S3.master-prompt.md}) jest
        generowany z analizy i nie jest edytowany ręcznie; przy rozbieżności wygrywa analiza.
  \item \textbf{Przygotuj klon.} \repopath{./scripts/setup.sh} sprawdza narzędzia i instaluje
        hook skanu sekretów (w rdzeniu, standardach i legacy-risk-scan; w Portcullis
        \path{./scripts/install-hooks.sh}; w benchu
        \repopath{git config core.hooksPath .githooks}). Zainstalowany hook bez skanera
        zawsze odmawia commitu, ale repozytoria różnią się tym, czy hook trafi do klonu
        i czy zamiast gitleaks wystarczy Docker: w rdzeniu \repopath{setup.sh} bez gitleaks
        pomija instalację hooka, w Portcullis instalator bez gitleaks kończy się kodem 127,
        a w standardach i legacy-risk-scan hook bez gitleaks uruchamia obraz gitleaks przez
        Dockera (tabela w sekcji~\ref{sec:bezpieczenstwo-skan}). W rdzeniu
        \repopath{setup.sh} nie sprawdza \texttt{curl} ani \texttt{python3}, których
        potrzebują \path{scripts/e2e.sh} i \path{scripts/mutation.sh}, a \path{.env}, który
        tworzy, nie jest wczytywany automatycznie (sekcja~\ref{sec:rdzen-setup}).
  \item \textbf{Zbuduj i przetestuj lokalnie} tak jak CI. W rdzeniu:
\begin{lstlisting}
dotnet build SecondKey.slnx && dotnet test SecondKey.slnx
./scripts/mutation.sh             # albo: ./scripts/mutation.sh Artifacts (tylko projekty z tym słowem w nazwie)
./scripts/e2e.sh                  # łańcuch na sklepie przykładowym
./scripts/scan-secrets.sh         # skan sekretów jak w CI; z --staged: jak hook
\end{lstlisting}
        a \texttt{./scripts/setup.sh --check} tylko raportuje braki. W pozostałych
        repozytoriach: Portcullis --- build bez ostrzeżeń, testy, Stryker.NET dla reguł
        migracyjnych i zasady baseline (podrozdział~\ref{sec:portcullis-wklad}); standardy
        --- \repopath{validate}, \repopath{generate --check}, testy
        i \path{./scripts/verify-package.sh}
        (podrozdział~\ref{sec:standardy-procedura-zmiany}); bench --- PSScriptAnalyzer
        z ustawieniami repozytorium i skrypty tylko w ASCII
        (podrozdział~\ref{sec:bench-higiena}); legacy-risk-scan --- \repopath{npm test},
        \repopath{npm run e2e} i \repopath{npm run smoke:osv}
        (podrozdział~\ref{sec:risk-scan-lokalnie}) oraz dwa jednorazowe ustawienia
        repozytorium (podrozdziały~\ref{sec:risk-scan-deploy} i \ref{sec:risk-scan-ci}).
  \item \textbf{Formaty artefaktów to interfejs produktu.} Zmiana \path{schemas/*.json}
        to zmiana formatu: podbij pole \repopath{v}, zaktualizuj \path{docs/formats/} i
        utrzymaj walidujące się próbki.
  \item \textbf{Obsługa napisów jawna co do kultury} wszędzie (ostrzeżenia CA1305, CA1307,
        CA1309, CA1310 są błędami): produkt istnieje po to, żeby łapać błędy kulturowe.
  \item \textbf{Decyzje zapisuj w ADR} (\path{docs/adr/}), z powodem.
  \item \textbf{Bez modelu w ścieżce werdyktu}: nic w \repopath{SecondKey.Contract} ani
        \repopath{SecondKey.Compare} nie może wołać modelu językowego. \textbf{Nagrania
        z systemów rzeczywistych nigdy nie trafiają do gita}: tylko syntetyczne próbki
        (\path{schemas/samples/}) i fixture'y testów (\path{tests/**/fixtures/}). Pełna lista
        reguł, które łatwo złamać przez przypadek: podrozdział~\ref{sec:rdzen-zasady}.
  \item \textbf{Pull request według szablonu} \path{.github/pull_request_template.md}:
        w rdzeniu ticket z każdą linią „done when” dosłownie, co się zmieniło, jak to
        sprawdzono i decyzje z ryzykami (podrozdział~\ref{sec:rdzen-setup}). Szablon
        Portcullis wymaga poleceń z prawdziwym wynikiem, buildu bez ostrzeżeń, a dla nowej
        reguły testu w obu kierunkach i mutanta; szablon standardów --- decyzji o wersji;
        szablon legacy-risk-scan --- sprawdzenia prywatności.
  \item \textbf{Scalanie.} Od 26 IX 2026, na stałe polecenie właściciela, pull request jest
        scalany, gdy każdy check na jego head commicie jest zielony; stosy scala się po
        kolei, commitami merge. Nikt nie przegląda PR przed scaleniem: checki --- testy,
        progi mutacyjne, job end-to-end, CodeQL, skan sekretów --- są przeglądem
        (zaakceptowane ryzyko A3).
  \item \textbf{Wydania.} Standardy: wersję w \path{catalog/pack.json} ustawia się ręcznie
        według tabeli podniesień, a wydanie to data zamiast \repopath{Unreleased}
        w \path{CHANGELOG.md}, \repopath{generate}, scalenie i tag \texttt{v<wersja>} na
        \repopath{main}, sprawdzany w CI (podrozdział~\ref{sec:standardy-procedura-zmiany});
        runbook P6 wymaga takiego wydania przed P6. Portcullis: \texttt{<Version>}
        w \path{Directory.Build.props} i domyślne wejście \repopath{version}
        w \path{action.yml} muszą być równe przed tagiem, a pierwsze wydanie potrzebuje
        sekretu \repopath{NUGET\_API\_KEY} i upublicznienia pakietu GHCR
        (podrozdział~\ref{sec:portcullis-dystrybucja}).
  \item \textbf{Zaktualizuj ten dokument}, gdy zmiana dotyczy czegoś, co opisuje: źródło
        LaTeX w \path{docs/papers/}, diagramy w \path{docs/diagrams/} (jeden plik
        \path{<slug>.mmd} na diagram, dołączany przez \texttt{\textbackslash dgm\{<slug>\}}).
        Dokument nie dodaje faktów spoza źródeł, każda podsekcja kończy się listą
        \texttt{\textbackslash zrodla} z plikami, z których pochodzi, a po zmianie źródeł
        przesuwa się commit w tabeli z sekcji~\ref{sec:wprowadzenie-wersja}
        (\path{scripts/check-doc-sources.mjs} sprawdza, czy cytowane pliki rdzenia
        istnieją, i ostrzega, gdy kod rdzenia odszedł od tego commitu). Lokalnie:
\begin{lstlisting}
npm ci && node scripts/render-diagrams.mjs
node scripts/check-papers.mjs && node scripts/check-doc-sources.mjs
cd docs/papers && latexmk -pdf letsgolegacy-framework.tex
\end{lstlisting}
        Pull request, który dotyka źródeł dokumentu, uruchamia \path{build-docs-pdf.yml}
        (podrozdział~\ref{sec:rdzen-ci-pdf}).
\end{enumerate}

\begin{xltabular}{\linewidth}{@{}>{\raggedright\arraybackslash}p{5.6cm} >{\raggedright\arraybackslash}X@{}}
\toprule
Objaw & Przyczyna i działanie \\
\midrule
\endfirsthead
\toprule
Objaw & Przyczyna i działanie \\
\midrule
\endhead
\texttt{[FAIL] .NET SDK \ldots{} found; this solution targets net10.0} & zainstaluj .NET 10 SDK; \repopath{setup.sh} rdzenia kończy się kodem 1 \\
\texttt{pre-commit: gitleaks is not installed, so nothing was scanned.} albo \texttt{pre-commit: no secret scanner available — refusing to commit.} & zainstaluj gitleaks (w standardach i legacy-risk-scan wystarczy działający Docker); job CI skanuje niezależnie \\
\texttt{pre-commit: a staged change looks like a credential.} albo \texttt{pre-commit: secrets detected in staged changes — commit refused.} & jeśli to poświadczenie --- najpierw rotacja, potem zdjęcie ze stagingu i odczyt ze zmiennej środowiskowej; jeśli nie --- wyjątek według zasad repozytorium (sekcja~\ref{sec:bezpieczenstwo-sekrety}) \\
Standardy: \repopath{generate --check} wymienia nieaktualne pliki & uruchom \repopath{generate} i zacommituj \path{generated/}; treść wersji z datą w \path{CHANGELOG.md} jest zamrożona --- najpierw podbij wersję \\
Bench: job \texttt{lint} zawodzi na znaku spoza ASCII & Windows PowerShell 5.1 czyta skrypt bez BOM jako ANSI; skrypty tylko w ASCII \\
Rdzeń: testy SQL Server lokalnie pominięte & brak działającego Dockera; w CI \repopath{SK\_REQUIRE\_DOCKER=1} zamienia brak Dockera w porażkę, a \repopath{SK\_SKIP\_DOCKER=1} pomija te testy celowo i wygrywa \\
legacy-risk-scan: pierwsze wdrożenie kończy się błędem 404 na \repopath{configure-pages} & ustaw \emph{Settings → Pages → Build and deployment → Source → GitHub Actions} \\
legacy-risk-scan: miesięczny pull request z danymi EOL nie powstaje & włącz \emph{Settings → Actions → General → Allow GitHub Actions to create and approve pull requests} \\
\bottomrule
\end{xltabular}

\textbf{Checklista}
\begin{checklista}
  \item Ticket z planu pracy; jeden ticket --- jeden pull request; „done when” dosłownie.
  \item Build, testy i --- dla zmian logiki --- mutacje lokalnie przed pushem.
  \item Hook skanu sekretów działa w tym klonie; \path{./scripts/scan-secrets.sh} czysty.
  \item Zmiana formatu: podbite \repopath{v}, zaktualizowane \path{docs/formats/}, próbki
        walidują.
  \item Standardy: zregenerowane \path{generated/} i decyzja o wersji.
  \item Decyzja zapisana w ADR; ten dokument i tabela commitów zaktualizowane.
\end{checklista}

\zrodla{\path{letsgolegacy.secondkey/CLAUDE.md},
\path{letsgolegacy.secondkey/CONTRIBUTING.md},
\path{letsgolegacy.secondkey/.claude/settings.json},
\path{letsgolegacy.secondkey/docs/WORKPLAN.md},
\path{letsgolegacy.secondkey/docs/analysis/phase-01.md},
\path{letsgolegacy.secondkey/docs/analysis/S1-S3.master-prompt.md},
\path{letsgolegacy.secondkey/scripts/setup.sh},
\path{letsgolegacy.secondkey/scripts/hooks/pre-commit},
\path{letsgolegacy.secondkey/scripts/mutation.sh},
\path{letsgolegacy.secondkey/scripts/scan-secrets.sh},
\path{letsgolegacy.secondkey/.github/pull_request_template.md},
\path{letsgolegacy.secondkey/docs/papers/letsgolegacy-framework.tex},
\path{letsgolegacy.secondkey/scripts/check-doc-sources.mjs},
\path{letsgolegacy.secondkey/.github/workflows/build-docs-pdf.yml},
\path{letsgolegacy.portcullis/scripts/install-hooks.sh},
\path{letsgolegacy.portcullis/scripts/hooks/pre-commit},
\path{letsgolegacy.portcullis/.github/pull_request_template.md},
\path{letsgolegacy.secondkey-standards/README.md},
\path{letsgolegacy.secondkey-standards/scripts/setup.sh},
\path{letsgolegacy.secondkey-standards/scripts/hooks/pre-commit},
\path{letsgolegacy.secondkey-standards/.github/pull_request_template.md},
\path{letsgolegacy.secondkey-nopcommerce-bench/README.md},
\path{letsgolegacy.secondkey-nopcommerce-bench/scripts/README.md},
\path{letsgolegacy.secondkey-nopcommerce-bench/.githooks/pre-commit},
\path{letsgolegacy.secondkey-nopcommerce-bench/.github/workflows/legacy-build.yml},
\path{letsgolegacy.secondkey-nopcommerce-bench/docs/P6-RUNBOOK.md},
\path{letsgolegacy.legacy-risk-scan/README.md},
\path{letsgolegacy.legacy-risk-scan/scripts/setup.sh},
\path{letsgolegacy.legacy-risk-scan/scripts/hooks/pre-commit},
\path{letsgolegacy.legacy-risk-scan/.github/pull_request_template.md}}

\clearpage
\section{Bezpieczeństwo i higiena repozytoriów}\label{sec:bezpieczenstwo}

Second Key obsługuje nagrania ruchu innych systemów, które mogą zawierać dane osobowe i
poświadczenia. Ten rozdział zbiera w jednym miejscu, jak framework chroni dane i sekrety
oraz jak repozytoria pilnują tego mechanicznie.

\subsection{Dane w nagraniach i przebiegach}\label{sec:bezpieczenstwo-dane}

\begin{itemize}
  \item \textbf{Redakcja przed zapisem.} Capture redaguje nagłówki, zanim cokolwiek
        trafi na dysk: zawsze \repopath{authorization}, \repopath{proxy-authorization},
        \repopath{cookie}, \repopath{set-cookie}, \repopath{x-api-key}, oraz każdy dodany
        przez \repopath{--redact-header} albo \repopath{capture.redactHeaders}; redakcja jest
        po~nazwie nagłówka, więc nagłówek o~innej nazwie, na przykład własny
        \repopath{X-Auth-Token}, zostaje w~pliku jawnie. Nagłówek
        pliku \repopath{*.skcap} zapisuje listę zredagowanych nagłówków, a replay redaguje
        te same nagłówki w przebiegu. Replay nie wysyła wartości zredagowanych ani
        nagranych ciasteczek: każdy scenariusz po każdej stronie dostaje świeży cookie jar.
  \item \textbf{Ciała są nagrywane.} Redakcja obejmuje tylko nagłówki: ciała żądań
        i~odpowiedzi (jako JSON, tekst albo base64, do limitu \repopath{--max-body-bytes})
        oraz ścieżka i~query trafiają do nagrania bez zmian. Dotyczy to także pól formularzy,
        które kontrakt może czytać jako \repopath{request.body.form.<pole>} --- na przykład
        hasła wpisanego przy logowaniu --- i~tokenów podanych w~adresie.
        Dlatego w benchu żądania konfiguracji sklepu idą do witryny z pominięciem proxy
        nagrywającego, bo hasło administratora nigdy nie może trafić do nagrania,
        a scenariusze logują się wyłącznie kontami sklepu przykładowego: publicznym klientem
        danych przykładowych albo klientem, którego scenariusz sam zarejestrował.
        Tokenizacja danych osobowych w nagraniach należy do fazy 02
        (sekcja~\ref{sec:bezpieczenstwo-plan}).
  \item \textbf{Ruch sieciowy.} Nic nie jest wysyłane przez sieć poza dwoma systemami,
        które są porównywane; to własność kodu rdzenia, a~nie izolacja sieciowa (rdzeń nie
        ma sandboxa). Zasada projektowa 2: wszystko działa w~granicy klienta,
        domyślnie bez ruchu wychodzącego.
  \item \textbf{Nagrania systemów rzeczywistych nigdy nie trafiają do git.} Pliki
        \repopath{.gitignore} rdzenia, benchu, Portcullis, standardów oraz repozytoriów
        fazy 02 i 03 wykluczają \repopath{*.skcap}, \repopath{*.skrun}, \path{.secondkey/}
        i \path{evidence-out/}; wyjątki dotyczą tylko syntetycznych próbek i fixture'ów
        testów. To zasada z~\path{CONTRIBUTING.md} rdzenia, a~pilnują jej pliki
        \repopath{.gitignore}: skan sekretów nie rozpozna w~nagraniu danych osobowych. Bench
        publikuje nagrania, przebiegi i~pakiety wyłącznie jako artefakty CI, które CI
        przechowuje 30 dni (nagranie, przebieg i~SARIF rozgrzewki na eShopLegacyMVC: 14 dni);
        każdy przebieg workflowu tworzy nowe, a~trwałym zapisem wyniku jest
        \path{results/<ticket>/README.md} z~linkiem do przebiegu.
  \item \textbf{Pakiet dowodowy bez ruchu.} Opcję \repopath{--run} pomija się, gdy nagrany
        ruch musi zostać w~swojej granicy; pakiet nadal wyjaśnia werdykt, ale nie pozwala go
        przeliczyć. Raport i~\repopath{verdict.json} niosą jednak fragmenty nagranych danych:
        linie żądań (ścieżkę i~query) oraz wartości różnic (raport skraca wartość do 300
        znaków i~pokazuje do 50 różnic na wymianę, \repopath{verdict.json} ma wszystkie).
  \item \textbf{Raport HTML jest statyczny}: jeden plik bez skryptów i bez dostępu do sieci
        (zabrania tego content security policy), każda nagrana wartość zakodowana jako HTML.
        Przeglądarka drukująca PDF działa z wyłączonym sandboxem, bo sandbox nie jest
        dostępny dla roota w kontenerach i na części obrazów CI; renderuje wyłącznie tę
        statyczną stronę.
\end{itemize}

\zrodla{\path{letsgolegacy.secondkey/SECURITY.md},
\path{letsgolegacy.secondkey/CONTRIBUTING.md},
\path{letsgolegacy.secondkey/docs/cli.md},
\path{letsgolegacy.secondkey/docs/formats/README.md},
\path{letsgolegacy.secondkey/docs/formats/contract.md},
\path{letsgolegacy.secondkey/docs/formats/evidence.md},
\path{letsgolegacy.secondkey/schemas/skcap.schema.json},
\path{letsgolegacy.secondkey/src/SecondKey.Replay/ReplayEngine.cs},
\path{letsgolegacy.secondkey/src/SecondKey.Evidence/HtmlReport.cs},
\path{letsgolegacy.secondkey/docs/WORKPLAN.md},
\path{letsgolegacy.secondkey/docs/analysis/phase-01.md},
\path{.gitignore} w repozytoriach \path{letsgolegacy.secondkey},
\path{letsgolegacy.secondkey-nopcommerce-bench}, \path{letsgolegacy.portcullis},
\path{letsgolegacy.secondkey-standards}, \path{letsgolegacy.secondkey-capture},
\path{letsgolegacy.secondkey-sandbox}, \path{letsgolegacy.secondkey-portal};
\path{letsgolegacy.secondkey-nopcommerce-bench/scripts/6-configure-store.ps1},
\path{letsgolegacy.secondkey-nopcommerce-bench/scripts/lib/NopBench.Shop.ps1},
\path{letsgolegacy.secondkey-nopcommerce-bench/results/p3/README.md},
\path{letsgolegacy.secondkey-nopcommerce-bench/results/p4/README.md},
\path{letsgolegacy.secondkey-nopcommerce-bench/results/p5/README.md},
\path{letsgolegacy.secondkey-nopcommerce-bench/results/p9/README.md},
\path{letsgolegacy.secondkey-nopcommerce-bench/.github/workflows/legacy-build.yml},
\path{letsgolegacy.secondkey-nopcommerce-bench/.github/workflows/warmup-eshop.yml}}

\subsection{Sekrety}\label{sec:bezpieczenstwo-sekrety}

\begin{itemize}
  \item \textbf{Tylko ze zmiennych środowiskowych.} \repopath{secondkey.yaml} nigdy nie
        zawiera sekretu: łańcuch połączenia jest wskazany nazwą zmiennej, która go trzyma
        (\repopath{connectionStringEnv}). \path{secrets.env.example} w rdzeniu wymienia
        każdą taką zmienną (\repopath{SK\_LEGACY\_DB}, \repopath{SK\_CANDIDATE\_DB},
        \repopath{SK\_CHROME\_PATH}); kopia robocza \repopath{.env} jest ignorowana przez
        git. Brak zmiennej zatrzymuje polecenie z jej nazwą.
  \item \textbf{Bench bez hasła do bazy.} Witryna legacy łączy się z SQL Server tożsamością
        puli aplikacji (uwierzytelnianie Windows); reset w replay dostaje łańcuch połączenia
        bez sekretu w \repopath{NOPBENCH\_SK\_SQL}. Jedyny sekret benchu --- hasło
        administratora sklepu --- jest generowane na przebieg, maskowane w logach CI i
        zapisywane tylko w \path{<work>\secrets\legacy-admin.json}, czytelnym dla bieżącego
        użytkownika, grupy Administrators i SYSTEM; diagnostyka nigdy go nie zbiera. Znane
        hasło można podać w \repopath{NOPBENCH\_ADMIN\_PASSWORD}, trzymając je poza plikami
        repozytorium.
  \item \textbf{Gdy skan znajdzie sekret}: najpierw rotacja poświadczenia, potem usunięcie go
        ze źródła i czytanie ze zmiennej środowiskowej, dopiero na końcu czyszczenie historii
        --- repozytoria są publiczne, a wypchnięty commit jest ujawniony w chwili
        wypchnięcia. Fałszywy alarm: w rdzeniu, Portcullis i benchu ścieżkę albo kształt
        dodaje się do \repopath{[allowlist]} w \path{.gitleaks.toml} w tym samym commicie,
        żeby wyjątek był przeglądalny; w legacy-risk-scan --- wąski wpis z komentarzem, nigdy
        przez wklejenie prawdziwej wartości do wyrażenia regularnego; w standardach
        przeformułowuje się przykład, zamiast dodawać ścieżkę do allowlisty.
\end{itemize}

\zrodla{\path{letsgolegacy.secondkey/docs/cli.md},
\path{letsgolegacy.secondkey/secrets.env.example},
\path{letsgolegacy.secondkey/src/SecondKey.Cli/Commands/ReplayCommand.cs},
\path{letsgolegacy.secondkey/scripts/hooks/pre-commit},
\path{letsgolegacy.portcullis/scripts/hooks/pre-commit},
\path{letsgolegacy.secondkey-nopcommerce-bench/scripts/README.md},
\path{letsgolegacy.secondkey-nopcommerce-bench/.githooks/pre-commit},
\path{letsgolegacy.secondkey-standards/.gitleaks.toml},
\path{letsgolegacy.secondkey-standards/.github/workflows/secret-scan.yml},
\path{letsgolegacy.legacy-risk-scan/.gitleaks.toml},
\path{letsgolegacy.legacy-risk-scan/.github/workflows/secret-scan.yml},
\path{letsgolegacy.legacy-risk-scan/scripts/hooks/pre-commit}}

\subsection{Skan sekretów: hook i CI}\label{sec:bezpieczenstwo-skan}

Każde repozytorium fazy 01 skanuje sekrety gitleaks w dwóch miejscach, celowo: hook
pre-commit łapie błąd, zanim stanie się historią, a job CI łapie klony bez hooka.
Zainstalowany hook bez skanera zawsze odmawia commitu; repozytoria różnią się tym, czy hook
w ogóle trafi do klonu i czy zamiast gitleaks wystarczy Docker --- opisuje to część „Hook”
każdego wiersza tabeli.

\begin{xltabular}{\linewidth}{@{}>{\raggedright\arraybackslash}p{4.6cm} >{\raggedright\arraybackslash}X@{}}
\toprule
Repozytorium & Konfiguracja gitleaks i hook \\
\midrule
\endfirsthead
\toprule
Repozytorium & Konfiguracja gitleaks i hook \\
\midrule
\endhead
\path{letsgolegacy.secondkey} & reguły domyślne plus reguła na łańcuch połączenia SQL Server z hasłem w treści; wyjątek ścieżki dla \path{secrets.env.example}. Hook: \path{./scripts/setup.sh} kopiuje go tylko wtedy, gdy gitleaks jest zainstalowany, a inaczej pomija ten krok --- commit przechodzi wtedy bez skanu i sekret wykrywa dopiero job CI \\
\path{letsgolegacy.portcullis} & reguły domyślne plus ta sama reguła SQL Server. Hook: \path{./scripts/install-hooks.sh} bez gitleaks kończy się kodem 127 i niczego nie instaluje \\
\path{letsgolegacy.secondkey-standards} & tylko reguły domyślne i świadomie bez allowlisty: przykłady w standardach używają placeholderów, a gdy skaner je oflaguje, poprawia się przykład. Hook: \path{./scripts/setup.sh} instaluje go zawsze; bez gitleaks hook uruchamia obraz \path{ghcr.io/gitleaks/gitleaks:latest} przez Dockera, a commit odrzuca dopiero wtedy, gdy nie ma ani gitleaks, ani działającego Dockera; \texttt{./scripts/setup.sh --check} kończy się wtedy błędem \\
\path{letsgolegacy.legacy-risk-scan} & tylko reguły domyślne: repozytorium nie ma sekretów; pierwszy potrzebny sekret dostaje własną regułę w tym samym commicie. Hook: jak w standardach --- \path{./scripts/setup.sh} instaluje go zawsze, a bez gitleaks hook używa obrazu kontenera przez Dockera \\
\path{letsgolegacy.secondkey-nopcommerce-bench} & reguły domyślne plus reguła SQL Server (bench łączy się przez uwierzytelnianie Windows, więc hasło w łańcuchu nigdy nie jest uprawnione); gitleaks w CI przypięty skrótem SHA-256 wydania. Hook: w \path{.githooks/}, włączany raz na klon przez \texttt{git config core.hooksPath .githooks}; wymaga gitleaks; Git for Windows uruchamia go we własnym bashu, więc działa też na hoście Windows \\
\bottomrule
\end{xltabular}

Reguła SQL Server ma \repopath{secretGroup = 1}: sekretem jest samo hasło, więc
placeholder w nawiasach ostrych (\repopath{Password=<...>}) przechodzi przez allowlistę
reguły, a prawdziwe hasło --- z \repopath{Password} albo \repopath{pwd}, z
\repopath{Server} albo \repopath{Data Source} --- jest blokowane.

W standardach i legacy-risk-scan hook kończy skan kodem 2 przy znalezisku, żeby awaria
narzędzia nie wyglądała jak wykryty sekret; commit jest odrzucany w obu przypadkach.
\texttt{git commit --no-verify} jest tam celowo głośnym wyjściem awaryjnym, a CI i tak
skanuje każdy push.

Lokalnym odpowiednikiem jobu CI jest \path{./scripts/scan-secrets.sh} w rdzeniu, standardach
i legacy-risk-scan: bez argumentu skanuje historię, tak jak CI, z \repopath{--staged} ---
zmiany w stagingu, tak jak hook, a w standardach i legacy-risk-scan z \repopath{--dir} ---
drzewo robocze jako pliki. W rdzeniu bez gitleaks kończy się kodem 127; w standardach
i legacy-risk-scan sięga wtedy po obraz kontenera.

Joby CI różnią się skanerem: rdzeń i Portcullis pobierają gitleaks 8.28.0 jako plik
wydania, bench tę samą wersję sprawdza skrótem SHA-256, a standardy i legacy-risk-scan
uruchamiają celowo obraz \path{ghcr.io/gitleaks/gitleaks:latest} („a detection ruleset is
only worth what its newest rule catches”) i dodatkowo co tydzień, żeby przeskanować
historię regułami dodanymi w międzyczasie. Każdy job skanuje całą historię
(\repopath{fetch-depth: 0}).

\zrodla{\path{.gitleaks.toml} i \path{.github/workflows/secret-scan.yml} w repozytoriach
\path{letsgolegacy.secondkey}, \path{letsgolegacy.portcullis},
\path{letsgolegacy.secondkey-standards}, \path{letsgolegacy.legacy-risk-scan},
\path{letsgolegacy.secondkey-nopcommerce-bench};
\path{letsgolegacy.secondkey/scripts/setup.sh},
\path{letsgolegacy.secondkey/scripts/hooks/pre-commit},
\path{letsgolegacy.secondkey/scripts/scan-secrets.sh},
\path{letsgolegacy.portcullis/scripts/install-hooks.sh},
\path{letsgolegacy.portcullis/scripts/hooks/pre-commit},
\path{letsgolegacy.secondkey-standards/scripts/setup.sh},
\path{letsgolegacy.secondkey-standards/scripts/hooks/pre-commit},
\path{letsgolegacy.secondkey-standards/scripts/scan-secrets.sh},
\path{letsgolegacy.legacy-risk-scan/scripts/setup.sh},
\path{letsgolegacy.legacy-risk-scan/scripts/hooks/pre-commit},
\path{letsgolegacy.legacy-risk-scan/scripts/scan-secrets.sh},
\path{letsgolegacy.secondkey-nopcommerce-bench/.githooks/pre-commit},
\path{letsgolegacy.secondkey-nopcommerce-bench/README.md}}

\subsection{Zależności, analiza statyczna, własność}\label{sec:bezpieczenstwo-zaleznosci}

\begin{itemize}
  \item \textbf{Dependabot} co tydzień w każdym repozytorium fazy 01: rdzeń (NuGet z
        grupą stosu testowego, GitHub Actions, npm dla renderera diagramów dokumentacji),
        Portcullis (NuGet, GitHub Actions, Docker; rodzina
        \repopath{Microsoft.CodeAnalysis*} wyłączona, bo podniesienie „podłogi” Roslyn to
        decyzja o wspieranych SDK, a nie aktualizacja), standardy (GitHub
        Actions, NuGet z wyłączeniem własnego pakietu \repopath{SecondKey.Standards}),
        legacy-risk-scan (npm, GitHub Actions), bench (GitHub Actions jednym zgrupowanym
        PR).
  \item \textbf{CodeQL} na pull request, pushu do \repopath{main} i co tydzień: C\# w rdzeniu,
        Portcullis i standardach; w legacy-risk-scan JavaScript/TypeScript i workflowy GitHub
        Actions zapytaniami \texttt{security-and-quality}. Bench nie ma CodeQL.
  \item \textbf{CODEOWNERS} w każdym repozytorium fazy 01, z jawnie wymienionymi ścieżkami
        istotnymi dla bezpieczeństwa: rdzeń \path{/.github/}, \path{/.gitleaks.toml},
        \path{/schemas/}, \path{/src/SecondKey.Capture/}; Portcullis \path{/.github/},
        \path{/.gitleaks.toml}, \path{/action.yml}, \path{/portcullis-baseline.json},
        \path{/src/Portcullis.Rules/}; standardy \path{/standards/}, \path{/catalog/},
        \path{/generated/}, \path{/src/}, \path{/.github/workflows/}; legacy-risk-scan
        \path{/src/lib/osv.js}, \path{/.github/workflows/}; bench \path{/.github/},
        \path{/.githooks/}, \path{/.gitleaks.toml}, \path{/contract/}, \path{/pins.json}.
        CODEOWNERS wskazuje właściciela do przeglądu, ale go nie wymusza i~nie zapisuje
        (podrozdział~\ref{sec:bezpieczenstwo-scalanie}).
  \item \textbf{Przypinanie akcji}: bench przypina każdą akcję SHA commita z wersją w
        komentarzu; pozostałe repozytoria używają tagów wersji głównej.
  \item \textbf{Uprawnienia workflowów i sekrety CI.} Każdy workflow deklaruje
        \texttt{permissions}, zwykle tylko \texttt{contents: read}; szersze dostają: CodeQL
        (\texttt{security-events: write}), wdrożenie legacy-risk-scan (\texttt{pages:
        write}, \texttt{id-token: write}), odświeżanie danych EOL (\texttt{contents: write},
        \texttt{pull-requests: write}), bramka Portcullis (\texttt{pull-requests: write})
        i jej SARIF (\texttt{security-events: write}) w osobnym workflowie, żeby job
        z tokenem komentarza tego uprawnienia nie miał, oraz wydanie Portcullis
        (\texttt{contents: write}, \texttt{packages: write}). Jedynym sekretem
        repozytorium, którego używa jakikolwiek workflow, jest \repopath{NUGET\_API\_KEY}
        Portcullis (klucz nuget.org ze wzorcem \texttt{Portcullis.*}); reszta korzysta
        z automatycznego tokenu GitHub. Bench nie zapisuje poświadczeń przy checkoucie
        (\texttt{persist-credentials: false}) we wszystkich workflowach poza skanem
        sekretów.
\end{itemize}

\zrodla{\path{.github/dependabot.yml}, \path{.github/workflows/*.yml} i
\path{.github/CODEOWNERS} w repozytoriach \path{letsgolegacy.secondkey},
\path{letsgolegacy.portcullis}, \path{letsgolegacy.secondkey-standards},
\path{letsgolegacy.legacy-risk-scan}, \path{letsgolegacy.secondkey-nopcommerce-bench};
\path{letsgolegacy.secondkey/package.json},
\path{letsgolegacy.portcullis/CONTRIBUTING.md}}

\subsection{Zgłaszanie podatności}\label{sec:bezpieczenstwo-zgloszenia}

Podatność zgłasza się prywatnie, nie w publicznym issue: przez prywatne zgłaszanie
podatności GitHub (\emph{Security → Advisories → Report a vulnerability}) w repozytorium,
którego dotyczy. Tak mówią \path{SECURITY.md} rdzenia i Portcullis; szablony issue
Portcullis odsyłają do tego formularza („Security report”). Portcullis opisuje też zakres
(pakiet analyzerów, CLI i kontener, akcja i \repopath{Portcullis.CiComment}, pipeline
wydań), to, co pomaga w zgłoszeniu, potwierdzenie w ciągu tygodnia i to, co podatnością nie
jest (podrozdział~\ref{sec:portcullis-bezpieczenstwo}). Pozostałe repozytoria nie mają
pliku \path{SECURITY.md}.

\zrodla{\path{letsgolegacy.secondkey/SECURITY.md},
\path{letsgolegacy.portcullis/SECURITY.md},
\path{letsgolegacy.portcullis/.github/ISSUE_TEMPLATE/config.yml}}

\subsection{Prywatność strony legacy-risk-scan}\label{sec:bezpieczenstwo-risk-scan}

Wklejony plik jest parsowany w przeglądarce i nigdzie nie jest wysyłany. Jedyne żądania
wychodzące idą do publicznego API OSV.dev: zapytania \repopath{querybatch} po najwyżej
1000 pakietów, których ciało buduje jedna funkcja (\repopath{toQueries()} w
\path{src/lib/osv.js}) --- ekosystem, nazwa i wersja pakietu, nic więcej --- przy
stronicowanej odpowiedzi ponowione z tokenem strony od OSV (najwyżej 10 stron), i po
jednym \repopath{GET} na każde advisory. Wysyłane są współrzędne każdego sprawdzalnego
pakietu, także o~nazwie wewnętrznej, jeśli występuje w~pliku; nazwa pliku, projektu
i~ścieżki nie są wysyłane. Adres IP odwiedzającego widzą zarówno GitHub Pages, jak
i~\repopath{api.osv.dev}. Content security policy strony pozwala łączyć się
tylko z samą stroną i z \repopath{api.osv.dev} oraz nie pozwala ładować skryptów, stylów
ani fontów z innych miejsc. Brak analityki, ciasteczek i pamięci przeglądarki; zestaw testów
e2e sprawdza to na zbudowanej stronie: cel każdego żądania, dokładne ciała żądań,
zablokowane połączenia, ciasteczka, \repopath{localStorage}, \repopath{sessionStorage},
IndexedDB.

Tekst z zewnątrz --- opisy advisories z OSV i zawartość pliku odwiedzającego --- trafia do
DOM wyłącznie przez \repopath{textContent} i \repopath{setAttribute}, nigdy jako znaczniki;
CodeQL analizuje kod strony właśnie pod kątem tej powierzchni. Zmiany tego, co strona
wysyła, mają pilnować CODEOWNERS (\path{src/lib/osv.js}) i~sekcja „Privacy check” szablonu
pull requestu, ale żaden scalony pull request tego repozytorium nie ma recenzji;
mechanicznie pilnuje tego test e2e z~dokładnymi ciałami żądań. Szczegóły:
podrozdział~\ref{sec:risk-scan-prywatnosc}.

\zrodla{\path{letsgolegacy.legacy-risk-scan/README.md},
\path{letsgolegacy.legacy-risk-scan/src/index.html},
\path{letsgolegacy.legacy-risk-scan/src/lib/osv.js},
\path{letsgolegacy.legacy-risk-scan/src/render.js},
\path{letsgolegacy.legacy-risk-scan/.github/workflows/codeql.yml},
\path{letsgolegacy.legacy-risk-scan/.github/CODEOWNERS},
\path{letsgolegacy.legacy-risk-scan/.github/pull_request_template.md}}

\subsection{Czego pakiet dowodowy jeszcze nie dowodzi}\label{sec:bezpieczenstwo-podpis}

W fazie 01 pakiet jest niepodpisany: statement in-toto ma \repopath{"signed": false}.
Dowodzi, które pliki należą do siebie i co mówią, a nic nie mówi o tym, kto go złożył.
Faza 03 dodaje podpis: statement in-toto w kopercie DSSE (\repopath{evidence.dsse}),
cosign z kluczem klienta z jego KMS lub HSM, bez publicznego transparency log,
przechowywanie w OCI (ORAS) albo WORM. Format pakietu się przy tym nie zmienia, bo faza 01
już niesie skróty SHA-256 i niepodpisany statement.

\zrodla{\path{letsgolegacy.secondkey/docs/formats/evidence.md},
\path{letsgolegacy.secondkey/docs/WORKPLAN.md},
\path{letsgolegacy.secondkey/docs/analysis/phase-01.md}}

\subsection{Ochrona danych zaplanowana na fazy 02 i 03}\label{sec:bezpieczenstwo-plan}

\begin{itemize}
  \item \textbf{Tokenizacja PII (capture, faza 02).} Capture tokenizuje dane osobowe
        wewnątrz sieci klienta przez sidecar Presidio; kryterium: nagrania, które opuszczają
        host capture, zawierają tokeny, a nie dane osobowe.
  \item \textbf{Bez ruchu wychodzącego (sandbox, faza 02).} Chart Helm ma domyślnie
        blokować ruch wychodzący (NetworkPolicy); bundle air-gapped dostarcza obrazy do
        rejestru klienta z manifestem skrótów. Hostowane przez nas może być wyłącznie demo
        albo wersja próbna na publicznym kodzie.
  \item \textbf{Tożsamość i ślad decyzji (portal, faza 03).} Logowanie przez IdP klienta
        (OIDC: Entra ID, ADFS); \repopath{authservice} tylko dla hostowanej przez nas wersji
        próbnej. Role: engineer, reviewer, risk (tylko odczyt). Każda decyzja recenzenta
        trafia do dziennika, do którego można tylko dopisywać.
\end{itemize}

Szczegóły i kryteria odbioru: rozdział~\ref{sec:faza02}.

\zrodla{\path{letsgolegacy.secondkey-capture/README.md},
\path{letsgolegacy.secondkey-capture/docs/WORKPLAN.md},
\path{letsgolegacy.secondkey-sandbox/docs/WORKPLAN.md},
\path{letsgolegacy.secondkey-portal/docs/WORKPLAN.md}}

\subsection{Scalanie bez przeglądu człowieka}\label{sec:bezpieczenstwo-scalanie}

Rdzeń zapisuje jako zaakceptowane ryzyko A3, że od 26 IX 2026, na stałe polecenie
właściciela, pull request jest scalany, gdy każdy check na jego head commicie jest
zielony, i nikt go wcześniej nie przegląda: przeglądem są checki --- testy, progi
mutacyjne, job end-to-end, CodeQL, skan sekretów. Koszt, jeśli to błąd: defekt, którego
żaden check nie łapie, trafia do \repopath{main} bez przeglądu; testy mutacyjne ograniczają
to, czego testy nie widzą na ścieżkach, które obejmują --- i nic więcej.

Pozostałe repozytoria fazy 01 wyglądają tak samo: żaden z~pull requestów rdzenia,
Portcullis, standardów, legacy-risk-scan i~benchu nie ma recenzji w~GitHubie, a~gałąź
\repopath{main} żadnego z~nich nie ma ochrony (GitHub API, 3~X~2026: \texttt{protected} jest
fałszem, brak rulesetów), więc CODEOWNERS i~szablon pull requestu niczego nie wymuszają.

\zrodla{\path{letsgolegacy.secondkey/docs/analysis/phase-01.md}}

\clearpage
\section{Status i plan}\label{sec:status}

Stan na commitach z sekcji~\ref{sec:wprowadzenie-wersja}, według planów pracy
repozytoriów. Status zmienia się z każdym scalonym pull requestem; aktualny jest w
\path{docs/WORKPLAN.md} danego repozytorium.

\subsection{Faza 01}\label{sec:status-faza01}

\begin{xltabular}{\linewidth}{@{}>{\raggedright\arraybackslash}p{3.6cm} >{\raggedright\arraybackslash}X@{}}
\toprule
Repozytorium & Stan ticketów fazy 01 \\
\midrule
\endfirsthead
\toprule
Repozytorium & Stan ticketów fazy 01 \\
\midrule
\endhead
rdzeń & S0--S15 zrobione; S16 (outbound stubs) przeniesione do fazy 02, bo ani sklep
przykładowy, ani demonstracyjna konfiguracja nopCommerce nie wykonują wywołań
wychodzących. Nic nie jest wydane. \\
Portcullis & R0 (import), R1 (zmiana nazwy, pilnowana w CI), R2 (reguły migracyjne z
testami mutacyjnymi), R3 (SARIF z filtrem zmienionych linii i baseline) zrobione;
wsparcie P8 zaplanowane. Nic nie jest opublikowane: nie ma tagu \repopath{v*}, pakietów
w~nuget.org ani opublikowanego kontenera, a~workflow wydania nie miał jeszcze przebiegu;
publikacja przy pierwszym tagu \repopath{v*}. Według README pakiety i~obraz kontenera są
zbudowane i~sprawdzone lokalnie, ale CI buduje tylko pakiet analyzerów --- obrazu nie. \\
standardy & C4-a (standardy w Markdown z metadanymi), C4-b (generator z trybem
\repopath{--check}), C4-c (drift check) zrobione; R4 (skill dla modernize-dotnet)
scalony, zrobiony wtedy, gdy P6 potwierdzi, że agent go podjął. Żadna wersja nie jest
jeszcze otagowana, a~odnośniki do źródeł w~stopce skilla wskazują tag \texttt{v0.1.0},
którego nie ma. Sprawdzenie wykonywalne ma 11 z~20 standardów; pozostałe 9 to instrukcje
dla agenta. \\
bench nopCommerce & P1--P5 zrobione. P9 (bramka fazy 01) w toku: zabite mutanty
zrobione (\#8, 11 z 11); pakiet dowodowy kandydata, różnice i regresje czekają na
P6--P8. P10 (log godzin) w toku: godziny agenta do P5 i dla mutantów P9 są w logu, część
oszacowana po fakcie; godzin człowieka jeszcze nie ma. P6 (praca człowieka z GitHub
Copilot), P7 i P8 zaplanowane. \\
legacy-risk-scan & L18 scalony (\#1, \#2) i~zrobiony: strona działa na GitHub Pages od
26~IX~2026 (\url{https://konradcinkusz.github.io/letsgolegacy.legacy-risk-scan/}), a~job
\repopath{verify} w~\path{deploy.yml} przechodzi --- na \texttt{4f4e346} w~przebiegu
37116401196 strona serwuje ten commit, a~test z~prawdziwym OSV przechodzi
(podrozdział~\ref{sec:risk-scan-deploy}). Plan pracy odnotowuje wdrożenie od
\texttt{6bb73b1} (PR~\#9): wiersz L18 ma status „done”. Otwarte pozostają dwa kroki z~planu: ustawienie pozwalające
Actions otwierać pull requesty (miesięczne odświeżanie danych EOL z~1~X~2026 wypchnęło
gałąź, ale nie otworzyło PR) i~kontakt na stronie (ticket L17; dziś „Kontakt: wkrótce”). \\
capture, sandbox, portal & tylko plany; decyzja o topologii demonstracyjnej (D7) w stanie
„proposed”, zrobiona, gdy ADR zostanie przyjęty. \\
\bottomrule
\end{xltabular}

\zrodla{\path{letsgolegacy.secondkey/docs/WORKPLAN.md},
\path{letsgolegacy.secondkey/README.md},
\path{letsgolegacy.portcullis/docs/WORKPLAN.md},
\path{letsgolegacy.portcullis/docs/TUTORIAL.md},
\path{letsgolegacy.portcullis/README.md},
\path{letsgolegacy.portcullis/CHANGELOG.md},
\path{letsgolegacy.secondkey-standards/docs/WORKPLAN.md},
\path{letsgolegacy.secondkey-standards/README.md},
\path{letsgolegacy.secondkey-nopcommerce-bench/docs/WORKPLAN.md},
\path{letsgolegacy.legacy-risk-scan/docs/WORKPLAN.md},
\path{letsgolegacy.legacy-risk-scan/.github/workflows/deploy.yml},
\path{letsgolegacy.legacy-risk-scan/.github/workflows/update-eol.yml},
\path{letsgolegacy.secondkey-capture/docs/WORKPLAN.md},
\path{letsgolegacy.secondkey-sandbox/docs/WORKPLAN.md},
\path{letsgolegacy.secondkey-portal/docs/WORKPLAN.md}}

\subsection{Liczby, które już są}\label{sec:status-liczby}

\begin{xltabular}{\linewidth}{@{}>{\raggedright\arraybackslash}p{4.2cm} >{\raggedright\arraybackslash}X@{}}
\toprule
Co & Wynik zapisany przez CI \\
\midrule
\endfirsthead
\toprule
Co & Wynik zapisany przez CI \\
\midrule
\endhead
P3 --- rozgrzewka na eShopLegacyMVC & pakiet dowodowy wyprodukowany przez cały łańcuch (capture, A/A replay, werdykt, bramka, pakiet): 7 scenariuszy, 27 wymian; A/A 54 wyniki, 14 resetów, 0 bez odpowiedzi; werdykt \repopath{pass} --- 27 \repopath{equal}, 0 regresji; kontrakt 13 klauzul, wszystkie zaakceptowane i utrzymane, 0 niesprawdzonych, udział nieobecności 23\,\% (3 \repopath{never}); bramka na całym drzewie legacy: 27 znalezisk, \repopath{sk gate} \repopath{fail} z 7 blokującymi --- zapisane, nie egzekwowane \\
P4 --- zestaw ruchu nopCommerce & 38 scenariuszy (30 zaplanowanych, niektóre jako kilka sesji), 290 wymian \\
P5 --- kontrakt i dowód A/A & werdykt \repopath{pass}: 76 z 76 zaakceptowanych klauzul utrzymanych, 0 niesprawdzonych; 15 klauzul \repopath{never} (udział nieobecności 19,7\,\%); 290 wymian w 38 scenariuszach: 272 \repopath{equal}, 18 \repopath{equal-under-contract}, 0 regresji, 0 fix candidates \\
P9 --- zabite mutanty & 11 z 11: każdy z M01--M11 daje w \repopath{sk compare} wynik \repopath{fail}, a każda klauzula zarejestrowana dla mutanta z góry na nim zregresowała; M00 (kalibracja: źródło bez zmian) \repopath{pass} --- 290 wymian, 269 \repopath{equal}, 21 \repopath{equal-under-contract}, 0 regresji; to wynik dla tych jedenastu defektów na tym ruchu, a nie mutation score --- zasięg kontraktu zmierzą generowane mutanty \repopath{sk mutate} (faza 02) \\
Testy mutacyjne rdzenia & Stryker.NET w~CI na \texttt{c3ae5c2} (przebieg 37116794634), pięć z~siedmiu projektów produktu: \repopath{SecondKey.Compare} 98,89\,\%, \repopath{SecondKey.Evidence} 91,67\,\%, \repopath{SecondKey.Replay} (bez kodu SQL Server) 89,71\,\%, \repopath{SecondKey.Artifacts} 89,09\,\%, \repopath{SecondKey.Contract} 86,87\,\%; progi 85, 70 i~60 (high, low, break), poniżej 60 CI zawodzi; poza zakresem są \repopath{SecondKey.Capture} i~\repopath{SecondKey.Cli}. Wyniki są w~logach i~artefaktach CI, nie w~repozytorium; wynik \repopath{SecondKey.Replay} wzrósł z~60,7\,\% do 89,5\,\% w~drugiej rewizji analizy fazy 01 \\
Testy mutacyjne reguł migracyjnych Portcullis & 100,00\,\% w ostatnim zapisanym przebiegu Stryker.NET: 214 z 214 testowanych mutantów wykrytych (212 zabitych, 2 przekroczenia czasu); \path{mutation.yml} zawodzi poniżej progu 95 \\
\bottomrule
\end{xltabular}

Liczby P3--P5 i~P9 pochodzą z~przebiegów na starych pinach narzędzi (\texttt{967c49c},
\texttt{84e1925}; sekcja~\ref{sec:wprowadzenie-wersja}). Po przesunięciu pinów (ADR 0008)
workflowy benchu przechodzą na zielono, ale nowych liczb szczegółowych nie zapisano
w~\path{results/}.

\zrodla{\path{letsgolegacy.secondkey-nopcommerce-bench/docs/WORKPLAN.md},
\path{letsgolegacy.secondkey-nopcommerce-bench/results/p3/README.md},
\path{letsgolegacy.secondkey-nopcommerce-bench/results/p4/README.md},
\path{letsgolegacy.secondkey-nopcommerce-bench/results/p5/README.md},
\path{letsgolegacy.secondkey-nopcommerce-bench/results/p9/README.md},
\path{letsgolegacy.secondkey/docs/analysis/phase-01.md},
\path{letsgolegacy.secondkey/stryker-config.json},
\path{letsgolegacy.secondkey/scripts/mutation.sh},
\path{letsgolegacy.portcullis/docs/MUTATIONS.md},
\path{letsgolegacy.portcullis/tests/Portcullis.Engine.Tests/stryker-config.json}}

\subsection{Droga do bramki fazy 01}\label{sec:status-droga}

Pozostała ścieżka krytyczna przechodzi przez bench: P6 (migracja przez agenta, praca
człowieka) → P7 (replay legacy kontra kandydat; \repopath{verdict.json} istnieje, a różnica
NLS → ICU jest znaleziona albo jej brak udokumentowany) → P8 (reguły Portcullis na PR
kandydata; SARIF w pakiecie) → P9 (publikacja pakietu, kontraktu i liczb: scenariusze,
różnice, regresje, zabite mutanty). Zabite mutanty już są (podrozdział~\ref{sec:bench-p9}),
więc z liczb P9 brakuje już tylko tych, które przynosi kandydat. P10 zbiera godziny wdrożenia dla P4--P7; godziny
człowieka przyjdą z P6. Runbook P6 jest napisany przed pierwszym uruchomieniem; zapis
pierwszego przebiegu stanie się jego przykładem.

\zrodla{\path{letsgolegacy.secondkey-nopcommerce-bench/docs/WORKPLAN.md},
\path{letsgolegacy.secondkey-nopcommerce-bench/docs/P6-RUNBOOK.md},
\path{letsgolegacy.secondkey/docs/WORKPLAN.md}}

\subsection{Później}\label{sec:status-pozniej}

Zakres faz 02--04 jest w sekcji~\ref{sec:architektura-fazy}; szczegóły repozytoriów fazy
02 i 03 w rozdziale~\ref{sec:faza02}. Wszystkie te prace zaczynają się dopiero po bramce
fazy 01.

\zrodla{\path{letsgolegacy.secondkey/docs/WORKPLAN.md}}

\clearpage
\section{Słownik pojęć}\label{sec:slownik}

Terminy w kolejności alfabetycznej. Nazwy z kodu i formatów są zapisane krojem kodu i
nie są tłumaczone.

\begin{description}[style=nextline, leftmargin=1.2em, font=\normalfont\bfseries]
  \item[A/A replay] Replay, w którym obie strony to system legacy. Pokazuje, że łańcuch
        nie wnosi własnych różnic i że kontrakt trzyma na legacy; w benchu kryterium P5.
  \item[analyzer] Reguła Roslyn uruchamiana przez kompilator albo przez Portcullis;
        wynik to diagnostyka z identyfikatorem i severity.
  \item[artefakt] Plik, którym komponenty się wymieniają: \repopath{*.skcap},
        \repopath{contract.yaml}, \repopath{*.skrun}, \repopath{verdict.json},
        \repopath{*.sarif}, pakiet dowodowy.
  \item[baseline (Portcullis)] Plik z odciskami znalezisk, które już istnieją; znalezisko
        przyjęte przez baseline jest raportowane, ale nie blokuje.
  \item[bench (okaz)] Repozytorium \path{letsgolegacy.secondkey-nopcommerce-bench}:
        publiczny dowód, że łańcuch działa na prawdziwym kodzie legacy.
  \item[bramka (gate)] Portcullis na pull request: blokuje błędy w liniach zmienionych
        przez PR. \repopath{sk gate} podsumowuje jej SARIF.
  \item[bramka fazy 01] Ticket P9 benchu: publiczny pakiet dowodowy na nopCommerce 3.x.
  \item[capture] Nagrywanie ruchu przez proxy \repopath{sk capture} przed systemem legacy,
        bez zmiany jego kodu.
  \item[drift check] Krok CI repozytorium docelowego, który porównuje przypiętą wersję
        standardów z opublikowaną i skill z pakietem.
  \item[DSSE] Koperta podpisu dla statementu in-toto; w fazie 03 (\repopath{evidence.dsse}).
  \item[ekstraktor] Reguła, która wyciąga wartość z HTML, JSON, nagłówka albo tekstu,
        zanim ocenią ją klauzule; \repopath{all: true} zbiera listę.
  \item[evidence pack (pakiet dowodowy)] Katalog z raportem, werdyktem, kontraktem,
        opcjonalnie przebiegiem, SARIF, statementem in-toto i manifestem skrótów SHA-256.
  \item[fix candidate] Klasa \repopath{fix-candidate}: kandydat spełnia klauzulę, którą
        legacy łamie; decyduje człowiek (wynik \repopath{review}).
  \item[kandydat] Zmigrowana wersja systemu, wskazana commitem head pull requestu.
  \item[klauzula] Element kontraktu: \repopath{must} (musi się dziać) albo
        \repopath{never} (nigdy nie wolno); ma \repopath{when}, \repopath{assert},
        \repopath{why} i pochodzenie.
  \item[klauzula nieobecności] Klauzula \repopath{never}; udział nieobecności to część
        zaakceptowanych klauzul, które są \repopath{never}.
  \item[korelacja] Reguła replay, która przenosi wartość wygenerowaną przez serwer (np.
        token anti-forgery) z ostatniej odpowiedzi tej samej strony do kolejnych żądań.
  \item[legacy] System przed migracją; strona \repopath{legacy} w replay i werdykcie.
  \item[manifest] \repopath{manifest.json} pakietu: każdy plik z SHA-256 i rozmiarem.
  \item[maska] Reguła normalizacji, po której dwa napisy są równe (\repopath{timestamps},
        \repopath{guids}, nazwane wzorce).
  \item[modernize-dotnet] Agent migracyjny GitHub Copilot app modernization for .NET.
  \item[NLS / ICU] Biblioteki globalizacji .NET Framework (NLS) i współczesnego .NET
        (ICU); ich różnice w porównywaniu i formatowaniu to główny kandydat na różnicę
        zachowania po migracji.
  \item[normalizacja] Sekcja kontraktu opisująca różnice bez znaczenia:
        \repopath{ignore}, \repopath{tolerances}, \repopath{sets}, \repopath{masks}.
  \item[pozytywny bliźniak] Klauzula lub krok scenariusza, który dowodzi, że scenariusz
        dotarł do miejsca, którego pilnuje klauzula \repopath{never}.
  \item[probe] Dwa znaczenia: sonda bazy w \repopath{secondkey.yaml}
        (\repopath{probe: sqlServer}, zmiany wierszy w tabelach per żądanie) oraz w benchu
        parametr query \repopath{probe=<nazwa>}, który nazywa żądanie dla klauzuli.
  \item[regression] Klasa \repopath{regression}: brak odpowiedzi, złamana zaakceptowana
        klauzula albo różnica, której kontrakt nie tłumaczy (fail closed).
  \item[replay] Odtwarzanie nagranych scenariuszy na obu stronach z resetem stanu przed
        każdym scenariuszem (\repopath{sk replay}).
  \item[reset stanu] Przywrócenie stanu startowego strony przed scenariuszem:
        \repopath{http}, \repopath{sqlServerSnapshot} albo \repopath{command}.
  \item[rdzeń] Repozytorium \path{letsgolegacy.secondkey}.
  \item[SARIF] Static Analysis Results Interchange Format 2.1.0 --- format wyniku bramki.
  \item[scenariusz] Sesja z nagrania odtwarzana jako całość (\repopath{session}) albo
        pojedyncza wymiana (\repopath{exchange}).
  \item[session key] Ciasteczko albo nagłówek, którego wartość identyfikuje sesję w
        capture (\repopath{--session-key}).
  \item[skill] Katalog z \repopath{SKILL.md}, który agent ładuje z
        \path{.github/skills/}; tu \path{applying-dotnet-migration-standards}.
  \item[snapshot] Snapshot bazy SQL Server, z którego reset przywraca stan startowy.
  \item[standardy] Reguły \repopath{SK-ARCH-001}--\repopath{007} i
        \repopath{SK-MIG-001}--\repopath{013} z \path{letsgolegacy.secondkey-standards}.
  \item[unexercised] Klauzula, której \repopath{when} nie pasowało do żadnego żądania;
        nigdy nie jest liczona jako spełniona.
  \item[violated-both] Klauzula złamana po obu stronach; o niczym nie decyduje i
        wskazuje błąd kontraktu wobec legacy.
  \item[werdykt] \repopath{verdict.json}: każda wymiana w jednej z klas \repopath{equal},
        \repopath{equal-under-contract}, \repopath{regression}, \repopath{fix-candidate};
        wynik \repopath{pass}, \repopath{fail} albo \repopath{review}.
  \item[wymiana (exchange)] Para żądanie--odpowiedź w nagraniu i w przebiegu.
\end{description}

\zrodla{rozdziały~\ref{sec:architektura}--\ref{sec:procedura} i ich źródła}


\appendix
\clearpage
\section{Workflowy CI we wszystkich repozytoriach}\label{app:ci}

Wszystkie workflowy GitHub Actions repozytoriów frameworku na przypiętych commitach, plus
workflow, który buduje ten dokument. Trigger \repopath{push} dotyczy gałęzi
\repopath{main}, chyba że zaznaczono inaczej; „(paths)” oznacza filtr ścieżek.

\begin{xltabular}{\linewidth}{@{}>{\raggedright\arraybackslash}p{3.1cm} >{\raggedright\arraybackslash}p{4.1cm} >{\raggedright\arraybackslash}X@{}}
\toprule
Plik & Triggery & Joby \\
\midrule
\endfirsthead
\toprule
Plik & Triggery & Joby \\
\midrule
\endhead
\multicolumn{3}{@{}l}{\textbf{\path{letsgolegacy.secondkey}}} \\
\path{ci.yml} & pull request, push, ręcznie & build i testy (Linux), build i testy (Windows), testy mutacyjne (macierz projektów) \\
\path{e2e.yml} & pull request, push, ręcznie & capture, replay, compare, evidence na sklepie przykładowym; pakiet jako artefakt \repopath{evidence-pack} \\
\path{codeql.yml} & pull request, push, co tydzień & analiza C\# \\
\path{secret-scan.yml} & pull request, push & gitleaks \\
\path{build-docs-pdf.yml} & push, pull request (paths), ręcznie & ten dokument: render diagramów, kontrola dokumentów, LaTeX, PDF jako artefakt \\
\midrule
\multicolumn{3}{@{}l}{\textbf{\path{letsgolegacy.portcullis}}} \\
\path{ci.yml} & push, pull request & brak dawnej nazwy w drzewie; build i testy; pakiet analyzerów na SDK .NET 8 (podłoga Roslyn) \\
\path{mutation.yml} & pull request (paths), push (paths) & Stryker dla reguł migracyjnych \\
\path{portcullis-pr-check.yml} & pull request & skan Portcullis na własnym PR \\
\path{portcullis-sarif.yml} & pull request & przefiltrowany SARIF \\
\path{release.yml} & tag \repopath{v*}, ręcznie & build, testy, pakowanie; publikacja NuGet; kontener do GHCR; release na GitHub; przesunięcie tagu wersji głównej \\
\path{codeql.yml} & pull request, push, co tydzień & analiza C\# \\
\path{secret-scan.yml} & pull request, push & gitleaks \\
\midrule
\multicolumn{3}{@{}l}{\textbf{\path{letsgolegacy.secondkey-standards}}} \\
\path{ci.yml} & pull request, push (także tagi \repopath{v*}), ręcznie & build, testy i kontrola drzewa generowanego; pakowanie i weryfikacja pakietów NuGet; drift check na konsumentach-fixture'ach; release guard \\
\path{codeql.yml} & pull request, push, co tydzień & analiza C\# \\
\path{secret-scan.yml} & pull request, push, co tydzień, ręcznie & gitleaks (cała historia) \\
\midrule
\multicolumn{3}{@{}l}{\textbf{\path{letsgolegacy.secondkey-nopcommerce-bench}}} \\
\path{legacy-build.yml} & pull request, push, ręcznie & lint PSScriptAnalyzer; build nopCommerce (MSBuild); narzędzia łańcucha; legacy pod IIS z SQL Server (instalacja, smoke, konfiguracja, nagranie, replay A/A); werdykt i pakiet \\
\path{warmup-eshop.yml} & pull request, push, ręcznie & narzędzia łańcucha; eShopLegacyMVC pod IIS (nagranie, replay A/A); werdykt, bramka i pakiet \\
\path{mutants.yml} & pull request (paths), ręcznie & P9: plan (katalog i łatki na przypiętym drzewie); build legacy; narzędzia łańcucha; dwanaście jobów mutantów pod IIS (build z łatką, legacy, mutant obok, nagranie, replay legacy kontra mutant); dwanaście werdyktów; podsumowanie „N z 11 zabitych” (bez przechodzącego M00 błąd) \\
\path{chain-tools.yml} & wywoływany z innych workflowów & \repopath{sk} i Portcullis z przypiętych commitów \\
\path{secret-scan.yml} & pull request, push & gitleaks (cała historia, wersja przypięta skrótem) \\
\midrule
\multicolumn{3}{@{}l}{\textbf{\path{letsgolegacy.legacy-risk-scan}}} \\
\path{ci.yml} & pull request, push, co tydzień, ręcznie & testy jednostkowe (Node 22); e2e (Playwright, OSV zamockowane); smoke z prawdziwym API OSV (błąd sieci nie blokuje) \\
\path{deploy.yml} & push, ręcznie & testy i build strony; wdrożenie na GitHub Pages; weryfikacja opublikowanej strony z prawdziwym OSV \\
\path{update-eol.yml} & co miesiąc, ręcznie, pull request (paths) & odświeżenie danych EOL z endoflife.date i pull request; regeneracja danych w PR generatora \\
\path{codeql.yml} & pull request, push, co tydzień, ręcznie & analiza \\
\path{secret-scan.yml} & pull request, push, co tydzień, ręcznie & gitleaks (cała historia) \\
\bottomrule
\end{xltabular}

Repozytoria \path{secondkey-capture}, \path{secondkey-sandbox} i \path{secondkey-portal}
nie mają jeszcze workflowów.

\zrodla{\path{.github/workflows/*.yml} we wszystkich repozytoriach na przypiętych
commitach}


\end{document}
